\documentclass[11pt,a4paper]{article}
\usepackage[utf8]{inputenc}
\usepackage[T1]{fontenc}
\usepackage[english]{babel}
\usepackage[margin=2.2cm]{geometry}
\usepackage{amsmath,amssymb,amsthm}
\usepackage{parskip}
\usepackage{enumitem}
\setlist{nosep,leftmargin=1.4em}
\usepackage{booktabs}
\usepackage{array}
\usepackage{xcolor}
\definecolor{linkblue}{RGB}{20,60,150}
\definecolor{defbg}{RGB}{240,244,252}
\definecolor{defframe}{RGB}{100,130,190}
\usepackage[most]{tcolorbox}
\newtcolorbox{defbox}[1][]{breakable,colback=defbg,colframe=defframe,boxrule=0.6pt,arc=1.5mm,left=2mm,right=2mm,top=1mm,bottom=1mm,title={#1},fonttitle=\bfseries}
\newtheorem{theorem}{Theorem}[section]
\theoremstyle{definition}
\newtheorem{definition}{Definition}[section]
\usepackage[colorlinks=true,linkcolor=linkblue,urlcolor=linkblue]{hyperref}
\usepackage{algorithm}
\usepackage{algpseudocode}

% Quotation macro carried over from the Czech original (babel-czech \uv).
\providecommand{\uv}[1]{``#1''}

\title{\textbf{Multiagent Systems}\\[0.3em]\large Study Text for the Final State Examination}
\author{Final State Examination in Artificial Intelligence, MFF UK}
\date{July 2026}

\begin{document}
\maketitle

\vspace{-1.5em}

\begin{center}
\begin{minipage}{0.93\textwidth}
\small
This text covers the examination topic \uv{Multiagent Systems} of the master's final state
examination in Artificial Intelligence at the Faculty of Mathematics and Physics, Charles
University. It is based on the lecture course NAIL106 \emph{Multiagent Systems} (MFF UK) and on
the textbooks M.~Wooldridge: \emph{An Introduction to MultiAgent Systems} (2nd ed., 2009),
Y.~Shoham, K.~Leyton-Brown: \emph{Multiagent Systems} (2009)
and S.~Russell, P.~Norvig: \emph{Artificial Intelligence: A~Modern Approach}.
This is an English translation of the Czech study text \emph{Multiagentní systémy};
terminology follows the standard usage of the sources listed above.
\end{minipage}
\end{center}

\section*{Mapping of the examination topic to chapters}
\addcontentsline{toc}{section}{Mapping of the examination topic to chapters}

\begin{center}
\small
\begin{tabular}{p{0.53\textwidth} p{0.40\textwidth}}
\toprule
\textbf{Item of the official examination topic} & \textbf{Chapter in this text} \\
\midrule
Architecture of an autonomous agent; action selection mechanism
  & \ref{sec:uvod}~Introduction, \ref{sec:architektura}~Architecture of an autonomous agent \\
Methods for controlling agents; symbolic reactive planning
  & \ref{sec:deliberativni}~Deliberative architectures, \ref{sec:reaktivni}.1--\ref{sec:reaktivni}.3 \\
Connectionist reactive planning
  & \ref{sec:reaktivni}.4 Agent network architecture (Maes) \\
Hybrid approaches
  & \ref{sec:reaktivni}.6 Hybrid architectures, \ref{sec:reaktivni}.7 IDA \\
The pathfinding problem
  & \ref{sec:cesta}~Pathfinding (A*, D*~Lite, MAPF) \\
Communication and knowledge in MAS, ontologies
  & \ref{sec:ontologie}~Knowledge and ontologies \\
Speech acts, FIPA-ACL, protocols
  & \ref{sec:komunikace}~Agent communication \\
Distributed problem solving, cooperation
  & \ref{sec:cdps}~Distributed problem solving (CNET, BBS, DCOP) \\
Nash equilibria, Pareto efficiency
  & \ref{sec:hry}~Agent interaction and game theory (incl.\ coalitional games) \\
(social choice as a complement to preference aggregation)
  & \ref{sec:volby}~Social choice and voting \\
Resource allocation, auctions
  & \ref{sec:aukce}~Resource allocation, auctions, negotiation (+ mechanism design) \\
Methods for agent learning; reinforcement learning
  & \ref{sec:uceni}~Agent learning \\
Design methodologies, MAS languages and environments
  & \ref{sec:metodologie}~Methodologies, languages, and platforms \\
\bottomrule
\end{tabular}
\end{center}

\tableofcontents
\newpage

\section{Introduction: agent and multiagent system}
\label{sec:uvod}

\subsection{Motivation}

Multiagent systems (MAS) emerged as an independent area of computer science in the 1990s in
response to several concurrent trends: the \emph{ubiquity} of computing technology,
\emph{connectivity and distribution} (computation as interaction), the growing \emph{complexity}
of tasks, the \emph{delegation} of tasks to software, and the pressure towards ease of use. The
classical paradigm \uv{a program receives an input, computes an output, and terminates} ceases to
be adequate the moment software is permanently embedded in an environment that changes
independently of it and in which other independent components run alongside it.

MAS can be viewed from four different perspectives, which explains the breadth of the topic:
\begin{itemize}
\item \textbf{Software engineering.} An agent is the next level of abstraction in the sequence
  machine code $\to$ structured programs $\to$ objects $\to$ distributed objects (CORBA)
  $\to$ agents. The classic quip: \uv{objects do it for free, agents do it because they want to}.
\item \textbf{Distributed and parallel systems.} Decades of research have solved mutual exclusion,
  deadlock, and consensus. The difference is that agents are \emph{autonomous}, coordination
  mechanisms are not given in advance, and everything has to be resolved at run time.
\item \textbf{Artificial intelligence.} Traditionally: MAS is a subfield of AI. According to
  Russell and Norvig the goal of AI is conversely the design of intelligent agents (hence
  AI $\subseteq$ MAS); Etzioni's pointed version is \uv{MAS is 1~\% AI and 99~\% computer
  science}. MAS uses AI techniques (knowledge representation, planning) but emphasises the
  \emph{social} aspects that classical AI overlooked.
\item \textbf{Game theory and economics.} A formal apparatus for rational decision making by
  multiple parties. Mathematical game theory deals mainly with questions of \emph{existence},
  whereas MAS deals with the \emph{computational} aspects (complexity, practical usability).
\item \textbf{Sociology.} A key notion of MAS is \emph{agent societies}, and sociology studies
  human societies --- so the inspiration suggests itself. The same caveat applies here as to the
  relation between AI and human intelligence: inspiration yes, necessity no (aeroplanes do not
  flap their wings). In the opposite direction the link is stronger: sociology, economics, and
  ecology all use agents as a building block of their simulation models --- \emph{agent-based
  modelling} (ABM), see the platforms in chapter~\ref{sec:metodologie}.
\end{itemize}

\subsection{What an agent is}

There is no single agreed definition of an agent --- Franklin and Graesser (1996), in the paper
\emph{Is it an agent, or just a program?}, collected dozens of mutually incompatible definitions.
The crucial point is to distinguish an agent from an ordinary procedure or object. The most
frequently cited formulations:

\begin{defbox}[Agent --- working definitions]
\textbf{Russell, Norvig (AIMA):} An agent is anything that can be viewed as an entity
\emph{perceiving} its environment through sensors and \emph{acting} upon that environment through
effectors.

\textbf{Pattie Maes:} Autonomous agents are computational systems that inhabit some complex
dynamic environment, sense and act autonomously in this environment, and by doing so realise a
set of goals or tasks for which they were designed.

\textbf{Wooldridge, Jennings:} An agent is a computer system situated in some environment that is
capable of \emph{autonomous action} in that environment in order to meet its delegated goals.
\end{defbox}

\begin{defbox}[Multiagent system]
A multiagent system is a system composed of several agents that \emph{interact} with one another
--- most often by exchanging messages over a computer network. The agents \emph{need not share a
common goal}; they may have their own, even conflicting, interests. This is why it is necessary
to study mechanisms of negotiation, coordination, and dynamic resource allocation.
\end{defbox}

\textbf{Autonomy.} A human being is fully autonomous, a Java method not at all. An agent lies in
between: it should autonomously choose the \emph{way} in which a given problem is to be solved
(that is, the choice of subgoals), but not the main goal --- that is delegated to it.

\subsection{Properties of an intelligent agent}

\begin{itemize}
\item \textbf{Reactivity} --- it perceives the environment and is able to respond to changes in
  it within a reasonable (real) time.
\item \textbf{Proactiveness} --- it has its own goals and actively pursues them, on its own
  initiative, not merely as a response to a stimulus.
\item \textbf{Social ability} --- it is able to communicate with other agents and with humans.
\end{itemize}

Achieving \emph{purely} reactive or \emph{purely} goal-oriented behaviour is fairly easy. What is
hard is to find the \textbf{balance between reactivity and proactiveness}: an agent must not
blindly follow a plan that has become pointless, but neither may it constantly reconsider its
goals and never get anywhere. This trade-off recurs in various guises throughout the whole text
(see in particular commitment strategies in chapter~\ref{sec:deliberativni} and hybrid
architectures in chapter~\ref{sec:reaktivni}).

\subsection{Intentional systems and the intentional stance}

When we talk about agents, we ascribe \emph{mental states} to them: \uv{the agent believes
that\dots}, \uv{the agent wants\dots}. Daniel Dennett defines an \textbf{intentional system} as an
entity whose behaviour can be predicted by ascribing to it properties such as beliefs, desires,
and rationality. Intentional systems form a hierarchy: a \emph{first-order} system has beliefs
and desires, a \emph{second-order} system has beliefs and desires about beliefs and desires (its
own as well as those of others), and so on.

Dennett distinguishes three \emph{stances} towards predicting the behaviour of a system: the
\emph{physical} one (the laws of physics suffice --- if I throw a stone I need not speak of its
desires, mass, velocity, and gravity are enough), the \emph{design stance} (I do not know all the
physical laws, but I know the purpose for which the system was designed --- I can set an alarm
clock without understanding its electronics), and the \emph{intentional} one (I describe the
system in terms of beliefs, desires, and rationality).

Shoham gives a provocative example: \uv{The light switch in the room is a very cooperative agent
that is willing of its own free will to conduct electric current, but does so only when it
believes that we really want it to.} The description is consistent and elegant, yet most people
find it absurd --- \emph{because for a light switch we have a simpler and equally accurate
description}. The intentional stance pays off where a physical or design description does not
exist or is impractical (even with the schematic of a computer at hand it is hard to derive why
clicking an icon opens a window).
\textbf{From the point of view of computer science, the intentional stance is above all an
abstraction for managing complexity} --- and for many programmers that is precisely the main
contribution of MAS.

There remains the philosophical question Dennett poses: does the intentionality of entities exist
\emph{in itself}, or does it arise only \emph{in the process of interpretation} by an observer?
Dennett leans towards the second option --- no intentionality without interpretation. For us this
has a practical consequence: ascribing beliefs and desires to an agent is \emph{our description},
not a property of the implementation, and it is legitimate exactly to the extent that it helps us
predict behaviour.

\subsection{Exam summary}

\begin{itemize}
\item Agent = an autonomous system situated in an environment, which perceives and acts in it in
  order to meet delegated goals. MAS = a system of interacting agents \emph{without a necessarily
  shared goal}.
\item The key problem in agent design is \textbf{action selection}; an agent architecture is a
  software architecture that realises this selection.
\item Three properties of an intelligent agent: reactivity, proactiveness, social ability;
  the hard part is balancing them.
\item Dennett: the intentional system and the three stances (physical, design, intentional);
  the intentional stance is an abstraction for managing complexity.
\item Distinguishing MAS from: distributed systems (autonomy, coordination only at run time),
  objects (an agent performs an action only if it wants to), classical AI (emphasis on social
  aspects), game theory (emphasis on computational aspects).
\end{itemize}

\section{Architecture of an autonomous agent and the action selection mechanism}
\label{sec:architektura}

\subsection{The environment and its properties}

An agent's behaviour is fundamentally determined by the type of environment it finds itself in.
The standard classification (Russell, Norvig):

\begin{center}
\begin{tabular}{p{0.28\textwidth} p{0.64\textwidth}}
\toprule
\textbf{Dichotomy} & \textbf{Meaning} \\
\midrule
fully / partially observable
  & The agent's sensors give it the complete state of the environment --- or only part of it
    (hidden information, noise, limited range). \\
static / dynamic
  & The environment changes only as a result of the agent's actions --- or also otherwise (other
    agents, natural processes). In a static environment the agent can \uv{think its next move over
    at leisure}. \\
deterministic / non-deterministic
  & Every action has exactly one possible outcome --- or several (with a known or unknown
    probability distribution; in that case we speak of a stochastic environment). \\
discrete / continuous
  & A finite number of states and actions --- or continuous quantities (position, velocity). \\
\bottomrule
\end{tabular}
\end{center}

Further commonly listed dichotomies: \emph{episodic / sequential} (does today's action affect
future rewards?), \emph{single-agent / multiagent}, \emph{known / unknown} (does the agent know the
transition function in advance?).

\textbf{An unpleasant observation:} most interesting environments (the real world, the internet)
are continuous, non-deterministic, dynamic, and only partially observable. Precisely for this
reason it is not enough to program an agent as a fixed table of rules.

\medskip
\noindent\emph{Example --- a thermostat.} A not particularly intelligent agent, situated in an
environment (a room). Two percepts (\uv{it is cold}, \uv{the temperature is fine}), two actions
(switch on, switch off), and a decision unit: cold $\Rightarrow$ switch on, OK $\Rightarrow$ switch
off. The environment is nevertheless non-deterministic (someone opens the door).

\noindent\emph{Example --- a software daemon.} For instance the Unix \texttt{xbiff}: the
environment is the operating system, it perceives by calling system services, and its actions are
software operations (start a mail client, change an icon). Its decision algorithm is at the level
of the thermostat.

\subsection{An abstract agent architecture}

We now formalise the view of an agent and its interaction with the environment. We assume the
environment is discrete --- which is not restrictive, as every continuous environment can be
discretised to arbitrary precision.

\begin{defbox}[Environment, agent, run]
The \textbf{environment} can be in one of finitely many discrete states,
$E = \{e, e', \dots\}$. The \textbf{agent} has a finite set of actions at its disposal,
$A = \{a, a', \dots\}$.

A \textbf{run} of an agent in an environment is a sequence of alternating environment states and
actions:
\[
r = e_0 \xrightarrow{a_0} e_1 \xrightarrow{a_1} e_2 \xrightarrow{a_2} \cdots
\xrightarrow{a_{n-1}} e_n .
\]
Let $R$ denote the set of all finite runs (over $E$ and $A$), $R^A \subseteq R$ those that end
with an action, and $R^E \subseteq R$ those that end with an environment state.
\end{defbox}

\begin{defbox}[State transformer function of the environment]
\[
\tau : R^A \to 2^E
\]
The function $\tau$ maps a run ending in an action to the set of states the environment may reach
after that action. Two consequences follow from this definition:
\begin{itemize}
\item \textbf{The environment has a history} --- the next state depends not only on the last
  action but on the entire history of the run.
\item \textbf{The environment is non-deterministic} --- the result is a set of states, and we do
  not know in advance which one will occur.
\end{itemize}
If $\tau(r) = \emptyset$, the run has ended. In what follows we assume that all runs terminate.

An \textbf{environment} is a triple $\mathit{Env} = \langle E, e_0, \tau\rangle$, where $e_0$ is
the initial state.
\end{defbox}

\begin{defbox}[Agent and system]
An \textbf{agent} is a function mapping runs that end in an environment state to actions:
\[
Ag : R^E \to A .
\]
An agent therefore decides on the basis of the \emph{entire history} of the system, and it decides
\emph{deterministically}. Let $\mathcal{AG}$ denote the set of all agents.

A \textbf{system} is a pair $\langle Ag, \mathit{Env}\rangle$. The sequence
$(e_0, a_0, e_1, a_1, \dots)$ is a \emph{run of agent $Ag$ in environment $\mathit{Env}$} if and
only if
\begin{enumerate}
\item $e_0$ is the initial state of $\mathit{Env}$,
\item $a_0 = Ag(e_0)$,
\item for every $u > 0$ we have $e_u \in \tau\big((e_0,a_0,\dots,a_{u-1})\big)$
  and $a_u = Ag\big((e_0,a_0,\dots,e_u)\big)$.
\end{enumerate}
The set of all runs of agent $Ag$ in $\mathit{Env}$ is denoted $R(Ag,\mathit{Env})$.
\end{defbox}

\begin{defbox}[Functional equivalence of agents]
Agents $Ag_1$ and $Ag_2$ are \textbf{functionally equivalent with respect to $\mathit{Env}$} if
and only if $R(Ag_1,\mathit{Env}) = R(Ag_2,\mathit{Env})$.

They are \textbf{simply functionally equivalent} if they are functionally equivalent with respect
to \emph{all} environments.
\end{defbox}

Functional equivalence matters as a tool for comparing the \emph{expressive power} of
architectures: it says that the internal implementation is irrelevant as long as it leads to the
same observable runs.

\subsection{Purely reactive agents and agents with internal state}

\begin{defbox}[Purely reactive (tropistic) agent]
\[
Ag : E \to A
\]
It reacts only to the current state of the environment and ignores the history.

\textbf{Relation to the general agent:} for every purely reactive agent there exists a
functionally equivalent standard agent; the converse does \emph{not} hold.

\emph{Thermostat:} $Ag(e) = \mathrm{off}$ for $e = $ \uv{temperature OK}, otherwise
$Ag(e) = \mathrm{on}$.
\end{defbox}

\begin{defbox}[Agent with internal state]
Let $I$ be the set of internal states of the agent and $\mathit{Per}$ the set of its percepts. The
agent is determined by a triple of functions:
\begin{align*}
\mathit{see} &: E \to \mathit{Per} && \text{(perception)}\\
\mathit{next} &: I \times \mathit{Per} \to I && \text{(internal state update)}\\
\mathit{action} &: I \to A && \text{(action selection)}
\end{align*}
\textbf{The cycle:} the agent starts in state $i_0$; it perceives the environment
($\mathit{see}(e)$); it updates its state ($i \leftarrow \mathit{next}(i, \mathit{see}(e))$); it
selects an action ($a \leftarrow \mathit{action}(i)$); it executes the action and repeats.

\textbf{Theorem (without proof):} For every agent with internal state there exists a functionally
equivalent standard agent. Internal state therefore does \emph{not increase expressive power}, but
it is essential for the \emph{efficiency} and clarity of the implementation.
\end{defbox}

The function $\mathit{see}$ at the same time formalises \textbf{observability}: if $\mathit{see}$
maps two different states $e_1 \ne e_2$ to the same percept, the agent cannot distinguish them ---
the environment is only partially observable for it. The extreme cases are
$|\mathit{Per}| = |E|$ with $\mathit{see}$ injective (fully observable) versus
$|\mathit{Per}| = 1$ (the agent sees nothing).

\subsection{How to tell an agent what to do: utility}

We do not want agents with hard-wired behaviour. We want to tell the agent \emph{what} to do, not
\emph{how}. The standard AI approach is to define the goal \emph{indirectly}, by a measure of
success.

\begin{defbox}[Utility]
\textbf{Utility of a state:} $u : E \to \mathbb{R}$. The agent's goal is to reach states with a
high utility value.

Evaluating individual states is short-sighted. We therefore introduce the \textbf{utility of a
run}:
\[
u : R \to \mathbb{R} .
\]
This corresponds better to agents that run over the long term. The utility of a run can be derived
from states, e.g.\ as the utility of the \emph{worst} state visited or as the \emph{average}
utility of the states visited --- which is better depends on the task.
\end{defbox}

\begin{defbox}[Optimal agent]
Let $P(r \mid Ag, \mathit{Env})$ be the probability of run $r$ of agent $Ag$ in environment
$\mathit{Env}$, where $\sum_{r \in R(Ag,\mathit{Env})} P(r \mid Ag,\mathit{Env}) = 1$. An optimal
agent maximises \emph{expected utility}:
\[
Ag_{\mathrm{opt}} = \arg\max_{Ag \in \mathcal{AG}} \sum_{r \in R(Ag,\mathit{Env})}
u(r)\, P(r \mid Ag, \mathit{Env}).
\]
\end{defbox}

The definition is elegant but \textbf{non-constructive}: it does not say how to design such an
agent. Moreover, it is often difficult to specify the utility function at all. (A note: money is
not a good utility for a human being --- the utility of money is non-linear and differs for the
rich and the poor.)

\medskip
\noindent\textbf{Example: Tileworld.} A classical testbed for agent architectures. Agents move
around a chessboard in four directions; the board contains holes, obstacles, and tiles, and the
goal is to push the tiles into the holes. The environment is \emph{dynamic} --- holes, obstacles,
and tiles appear and disappear at random. The utility of a run is
\[
u(r) = \frac{N_f}{N_a},
\]
where $N_f$ is the number of holes filled by the agent during run $r$ and $N_a$ the number of all
holes that appeared during $r$. The agent has to be able both to \emph{react} to changes (the tile
I am pushing has disappeared) and to \emph{exploit opportunities} (a new tile has appeared next to
me). Historically, Tileworld was used precisely to investigate the trade-off between reactivity and
proactiveness (see agent commitment, section~\ref{ssec:commitment}).

\subsection{Predicate task specification}

A special case of utility is a mapping into $\{0,1\}$: a run is either successful ($u(r)=1$) or
not.

\begin{defbox}[Task environment]
Let $\Psi : R \to \{0,1\}$ be a \textbf{predicate specification}: $\Psi(r)$ holds if and only if
$u(r) = 1$. A \textbf{task environment} is a pair $\langle \mathit{Env}, \Psi \rangle$; it
specifies both the properties of the system (through $\mathit{Env}$) and the criterion for
completing the task (through $\Psi$).

Let $R_\Psi(Ag,\mathit{Env}) = \{ r \in R(Ag,\mathit{Env}) \mid \Psi(r) \}$.
\end{defbox}

When does an agent \emph{solve} a task? There are three readings:
\begin{itemize}
\item \textbf{Pessimistic:} $R_\Psi(Ag,\mathit{Env}) = R(Ag,\mathit{Env})$ --- \emph{all} runs
  satisfy $\Psi$.
\item \textbf{Optimistic:} $\exists r \in R(Ag,\mathit{Env}): \Psi(r)$ --- \emph{at least one} run.
\item \textbf{Realistic:} we extend $\tau$ with a probability distribution and measure success by
  the probability
  $P(\Psi \mid Ag,\mathit{Env}) = \sum_{r \in R_\Psi(Ag,\mathit{Env})} P(r \mid Ag,\mathit{Env})$.
\end{itemize}

\begin{defbox}[Types of tasks]
\textbf{Achievement tasks} --- \uv{find something}. A set of goal states $G \subseteq E$ is given;
$\Psi(r)$ holds if at least one state from $G$ occurs in run $r$. Typically every classical AI task
(search).

\textbf{Maintenance tasks} --- \uv{avoid something}. A set of \uv{bad} states $B \subseteq E$ is
given; $\Psi(r)$ fails to hold if a state from $B$ occurs in $r$. Typically games, where $B$ are
the losing positions and the environment is the opponent.

\textbf{Combination:} reach a state from $G$ while avoiding the states in $B$.
\end{defbox}

\subsection{The action selection mechanism --- an overview}
\label{ssec:vyber-akci}

\begin{defbox}[Action selection mechanism]
The procedure by which an agent decides at each moment which of the available actions to perform,
on the basis of percepts from the environment and its own internal state. An \textbf{agent
architecture} is a software architecture that realises such a mechanism --- according to Maes
(1991) \uv{a particular methodology for building agents that specifies how the agent can be
decomposed into a set of components and how these components should interact so that together they
answer the question of how the sensor data and the current internal state determine the actions and
the future internal state of the agent.}
\end{defbox}

The following two chapters discuss the main families of action selection mechanisms:

\begin{center}
\small
\begin{tabular}{p{0.20\textwidth} p{0.34\textwidth} p{0.38\textwidth}}
\toprule
\textbf{Family} & \textbf{Action selection mechanism} & \textbf{Typical representative} \\
\midrule
symbolic / deliberative
  & logical deduction over a database of formulae; planning
  & deductive agent, STRIPS, BDI, PRS \\
reactive (symbolic)
  & a set of \uv{situation $\to$ action} rules with fixed priorities
  & subsumption architecture (Brooks) \\
reactive (connectionist)
  & activation spreading in a network of modules, the most active one wins
  & agent network architecture (Maes) \\
hybrid
  & layers of different abstraction + arbitration between them
  & TouringMachines, InteRRaP \\
coalitional
  & competition of dynamically arising coalitions of processes for \uv{consciousness}
  & IDA / global workspace \\
learned
  & maximisation of a learned $Q$-function / stochastic policy
  & Q-learning, actor-critic (chapter~\ref{sec:uceni}) \\
\bottomrule
\end{tabular}
\end{center}

\subsection{Exam summary}

\begin{itemize}
\item Environment: 4 dichotomies (observability, static/dynamic, determinism, discreteness);
  real environments are in the hardest variant.
\item Formalism: $E$, $A$, run $r = e_0 a_0 e_1 \dots e_n$, $R^A$/$R^E$,
  state transformer function $\tau : R^A \to 2^E$ (hence history + non-determinism),
  $\mathit{Env}=\langle E,e_0,\tau\rangle$, agent $Ag : R^E \to A$, system
  $\langle Ag,\mathit{Env}\rangle$.
\item Functional equivalence (with respect to an environment / simple).
\item Purely reactive agent $Ag: E \to A$ (weaker); agent with internal state
  ($\mathit{see}$, $\mathit{next}$, $\mathit{action}$) --- both reducible to a standard agent.
\item Utility: $u : E \to \mathbb{R}$ vs.\ $u : R \to \mathbb{R}$; an optimal agent maximises
  \emph{expected} utility --- the definition is non-constructive.
\item Task environment $\langle \mathit{Env}, \Psi\rangle$; pessimist/optimist/realist;
  achievement tasks ($G$) and maintenance tasks ($B$).
\item Tileworld and the utility $N_f/N_a$ as the standard benchmark of reactivity vs.\
  proactiveness.
\end{itemize}

\section{Deliberative architectures: symbolic reasoning, BDI, PRS}
\label{sec:deliberativni}

\subsection{The classical AI approach and the deductive agent}

The classical way in which AI builds an \uv{intelligent system} rests on two pillars:
\begin{enumerate}
\item \textbf{Symbolic representation} of the environment and of behaviour --- here specifically by
  logical formulae.
\item \textbf{Syntactic manipulation} of this representation --- that is, logical deduction,
  theorem proving.
\end{enumerate}
The theory of the agent's behaviour is then also its program --- an \emph{executable specification}
that generates concrete actions. This approach runs into two classical problems:
\begin{itemize}
\item \textbf{The transduction problem}: how to translate the real world into a symbolic
  representation (computer vision, natural language processing, learning).
  In Goethe's words: \uv{All theory, dear friend, is grey, but the golden tree of life springs ever
  green.}
\item \textbf{The representation and reasoning problem}: how to represent knowledge and how to
  reason over it efficiently (theorem provers, planners).
\end{itemize}

\medskip
\noindent\textbf{A digression: deduction vs.\ induction.}
\emph{Deduction} is the derivation of conclusions that are true whenever the premises are true ---
a move from the general to the particular (the syllogism: \uv{All men are mortal, Socrates is a man,
therefore Socrates is mortal}). \emph{Induction} means that, given true premises, the conclusion is
more likely true than false --- from the particular to the general. Sherlock Holmes in fact uses
induction, even though he calls it deduction; mathematical induction, on the contrary, is deduction.

\begin{defbox}[Deductive agent (the agent as a theorem prover)]
Let $L$ be a set of first-order predicate logic formulae and $D = 2^L$ the set of databases of
such formulae. The agent's internal state is a database $\mathit{DB} \in D$. Reasoning is realised
by inference rules $\rho$:
\[
\mathit{DB} \vdash_\rho \varphi \quad\text{---\ formula } \varphi \text{ is derivable from }
\mathit{DB} \text{ using the rules } \rho .
\]
The agent is then a triple of functions
$\mathit{see} : E \to \mathit{Per}$, $\mathit{next} : D \times \mathit{Per} \to D$,
$\mathit{action} : D \to A$, where $\mathit{action}$ is defined by the procedure below.
\end{defbox}

\noindent\textbf{Action selection as theorem proving.}
\begin{algorithmic}[1]
\Function{action}{$\mathit{DB} : D$} : $A$
  \ForAll{$a \in A$}
    \If{$\mathit{DB} \vdash_\rho \mathit{Do}(a)$} \Return $a$ \EndIf
  \EndFor
  \ForAll{$a \in A$}
    \If{$\mathit{DB} \not\vdash_\rho \neg \mathit{Do}(a)$} \Return $a$ \EndIf
  \EndFor
  \State \Return $\mathit{null}$
\EndFunction
\end{algorithmic}

The first loop looks for an action that can be \emph{proved} to be desirable ($\mathit{Do}(a)$ is
derived). If there is no such action, the second loop looks for an action that is at least
\emph{consistent} with the database (it cannot be proved that it should not be done). If that fails
too, the agent does nothing.

\medskip
\noindent\textbf{Example: the 3$\times$3 vacuum world.}
Predicates: $\mathit{In}(x,y)$ (the agent is at position $(x,y)$), $\mathit{Dirt}(x,y)$ (there is
dirt at $(x,y)$), $\mathit{Facing}(d)$ (the agent faces direction $d$). Rules for deriving an
action:
\begin{align*}
\mathit{In}(x,y) \wedge \mathit{Dirt}(x,y) &\Rightarrow \mathit{Do}(\mathit{suck}) \\
\mathit{In}(0,0) \wedge \mathit{Facing}(\mathit{north}) \wedge \neg\mathit{Dirt}(0,0)
  &\Rightarrow \mathit{Do}(\mathit{forward}) \\
\mathit{In}(0,2) \wedge \mathit{Facing}(\mathit{north}) \wedge \neg\mathit{Dirt}(0,2)
  &\Rightarrow \mathit{Do}(\mathit{turn}) \qquad \dots
\end{align*}
The agent thus systematically traverses the world in a \uv{snake} pattern
00--01--02--12--11--10--\dots
Note that cleaning has higher priority than movement (the rule for $\mathit{suck}$ is listed first
and the other rules have the negation of $\mathit{Dirt}$ in their premises).

\medskip
\noindent\textbf{Pros and cons.}
\begin{itemize}
\item[$+$] Elegant, with clear (and beautiful) semantics; specification = program.
\item[$-$] Practically unusable for non-trivial tasks: proving is slow and need not terminate.
\item[$-$] \textbf{The problem of calculative rationality}: the agent derives the optimal action for
  the state of the environment as it was \emph{at the beginning} of the derivation --- meanwhile the
  environment has changed and the action is no longer optimal. A disaster in a rapidly changing
  environment.
\item[$-$] The function $\mathit{see}$ is hard to design (how to turn an image into formulae, how to
  represent time).
\end{itemize}
Nevertheless, individual elements of this approach reappear in later architectures, particularly in
BDI.

\subsection{Practical reasoning: the BDI architecture}

\begin{defbox}[Practical reasoning]
Inspired by human decision making. \emph{Theoretical} reasoning (Socrates) leads only to what we
think about the world; \emph{practical} reasoning leads to \textbf{actions}. It has two phases:
\begin{itemize}
\item \textbf{Deliberation} --- \emph{what we want to achieve}
  (\uv{I want to finish my master's degree}).
\item \textbf{Means-ends reasoning}, that is planning --- \emph{how we are going to achieve it}
  (construct a plan for finishing the degree).
\end{itemize}
And neither may take too long.
\end{defbox}

\begin{defbox}[Beliefs --- Desires --- Intentions]
\textbf{Beliefs} --- the informational state of the agent, its knowledge about the world.
In MAS we deliberately do not call them \uv{knowledge}, in order to stress that they are
\emph{subjective}, \emph{need not be true}, and \emph{may change in the future}. They may also
contain inference rules that permit forward chaining.

\textbf{Desires} --- the motivational state of the agent, the goals or situations it would like to
achieve. Desires \emph{may be mutually inconsistent} (I want both to study and to go for a beer).

\textbf{Intentions} --- states of the world that the agent has \emph{committed} itself to achieving.
Intentions lead to actions: it is for their sake that the agent plans. Intentions \textbf{persist}
until
\begin{enumerate}
\item the agent achieves them, or
\item it comes to believe that they are unachievable, or
\item the reasons that led to them cease to hold.
\end{enumerate}
Moreover, intentions \emph{constrain further deliberation} (I will not reconsider options
incompatible with an intention I have adopted) and \emph{influence future beliefs} (I count on my
intention succeeding).
\end{defbox}

\noindent\textbf{The difference between a desire and an intention} is illustrated by Bratman (1990):
\uv{My desire to play basketball this afternoon is merely a potential influencer of my conduct this
afternoon. It must vie with my other relevant desires\dots{} In contrast, once I
\emph{intend} to play basketball this afternoon, the matter is settled: I normally need not continue
to weigh the pros and cons; when the afternoon arrives, I will normally just proceed to execute my
intention.}

\begin{defbox}[The deliberation functions]
Let $B \subseteq \mathit{Bel}$ denote the agent's current beliefs, $D \subseteq \mathit{Des}$ its
current desires, and $I \subseteq \mathit{Int}$ its current intentions. Deliberation consists of
three functions:
\begin{itemize}
\item $\mathit{brf} : 2^{\mathit{Bel}} \times \mathit{Per} \to 2^{\mathit{Bel}}$ ---
  the \emph{belief revision function}, updating beliefs according to a percept;
\item $\mathit{options} : 2^{\mathit{Bel}} \times 2^{\mathit{Int}} \to 2^{\mathit{Des}}$ ---
  generating options (desires) from beliefs and existing intentions;
\item $\mathit{filter} : 2^{\mathit{Bel}} \times 2^{\mathit{Des}} \times 2^{\mathit{Int}}
  \to 2^{\mathit{Int}}$ --- selecting intentions from the options.
\end{itemize}
Splitting deliberation into \emph{option generation} and \emph{filtering} is essential: generation
is cheap and broad, whereas filtering is the actual decision making (it also takes existing
commitments into account).
\end{defbox}

\subsection{Planning: STRIPS}

\begin{defbox}[Means-ends reasoning / planning]
The process of deciding how to achieve a goal (an intention) using the available means (actions).
\begin{itemize}
\item \textbf{Input:} the goal (= intention $I$); the current state of the environment
  (= beliefs $B$); the actions available to the agent ($\mathit{Ac}$).
\item \textbf{Output:} a plan = a sequence of actions after whose execution the goal is satisfied.
\end{itemize}
\end{defbox}

\textbf{STRIPS} (Fikes, Nilsson, 1971) models the world by a set of first-order logic formulae.
Each action is described by a \emph{schema} with preconditions and effects (facts added and
deleted). The planning algorithm finds the difference between the goal and the current state and
reduces it by a suitable action (\emph{means-ends analysis}). This is elegant but not very
practical --- the algorithm often gets bogged down in low-level details.

\begin{defbox}[Action descriptor, planning problem, plan]
An \textbf{action descriptor} $a$ is a triple $\langle P_a, D_a, A_a\rangle$:
$P_a$ --- preconditions, $D_a$ --- the delete list, $A_a$ --- the add list; for simplicity these
are ground atomic formulae (without connectives and variables).

A \textbf{planning problem} is a triple $\langle B_0, O, G\rangle$, where $B_0$ are the initial
beliefs, $O = \{\langle P_a,D_a,A_a\rangle : a \in \mathit{Ac}\}$ the descriptors of all actions,
and $G$ a set of formulae describing the goal.

A \textbf{plan} $\pi = (a_1,\dots,a_n)$, $a_i \in \mathit{Ac}$, determines a sequence of databases
\[
B_i = (B_{i-1} \setminus D_{a_i}) \cup A_{a_i}, \qquad i = 1,\dots,n .
\]
A plan is \textbf{admissible} if $B_{i-1} \models P_{a_i}$ for every $i$.
A plan is \textbf{correct (valid)} if it is admissible and moreover $B_n \models G$.

\textbf{Planning} means, for a given $\langle B_0,O,G\rangle$, finding a correct plan or declaring
that none exists.
\end{defbox}

\medskip
\noindent\textbf{Example: the blocks world.}
Predicates: $\mathit{On}(x,y)$ ($x$ is on $y$), $\mathit{OnTable}(x)$, $\mathit{Clear}(x)$
(nothing is on $x$), $\mathit{Holding}(x)$ (the arm is holding $x$), $\mathit{ArmEmpty}$.
Actions:
\begin{center}
\small
\begin{tabular}{p{0.20\textwidth} p{0.36\textwidth} p{0.36\textwidth}}
\toprule
\textbf{Action} & \textbf{Pre / Del} & \textbf{Add} \\
\midrule
$\mathit{Stack}(x,y)$
  & Pre: $\{\mathit{Clear}(y), \mathit{Holding}(x)\}$ \newline
    Del: $\{\mathit{Clear}(y), \mathit{Holding}(x)\}$
  & $\{\mathit{ArmEmpty}, \mathit{On}(x,y)\}$ \\
$\mathit{UnStack}(x,y)$
  & Pre: $\{\mathit{On}(x,y),$ $\mathit{Clear}(x),\mathit{ArmEmpty}\}$ \newline
    Del: $\{\mathit{On}(x,y),\mathit{ArmEmpty}\}$
  & $\{\mathit{Holding}(x), \mathit{Clear}(y)\}$ \\
$\mathit{Pickup}(x)$
  & Pre: $\{\mathit{OnTable}(x),$ $\mathit{Clear}(x),\mathit{ArmEmpty}\}$ \newline
    Del: $\{\mathit{OnTable}(x),\mathit{ArmEmpty}\}$
  & $\{\mathit{Holding}(x)\}$ \\
$\mathit{PutDown}(x)$
  & Pre: $\{\mathit{Holding}(x)\}$ \newline
    Del: $\{\mathit{Holding}(x)\}$
  & $\{\mathit{ArmEmpty}, \mathit{OnTable}(x)\}$ \\
\bottomrule
\end{tabular}
\end{center}
Initial state $B_0 = \{\mathit{Clear}(A), \mathit{On}(A,B), \mathit{OnTable}(B),
\mathit{OnTable}(C), \mathit{Clear}(C), \mathit{ArmEmpty}\}$,
goal $G = \{\mathit{OnTable}(A), \mathit{OnTable}(B), \mathit{OnTable}(C)\}$.

\noindent A correct plan: $\pi = (\mathit{UnStack}(A,B),\ \mathit{PutDown}(A))$. Verification:
\begin{itemize}
\item $B_0 \models \{\mathit{On}(A,B),\mathit{Clear}(A),\mathit{ArmEmpty}\}$, so the first action
  is admissible; after it
  \[
  B_1 = \big(B_0 \setminus \{\mathit{On}(A,B),\mathit{ArmEmpty}\}\big) \cup
  \{\mathit{Holding}(A),\mathit{Clear}(B)\}.
  \]
\item $B_1 \models \{\mathit{Holding}(A)\}$, so the second action is admissible too; after it
  \[
  B_2 = \big(B_1 \setminus \{\mathit{Holding}(A)\}\big) \cup
  \{\mathit{ArmEmpty},\mathit{OnTable}(A)\}.
  \]
\item $B_2 \models G$, so the plan is correct.
\end{itemize}

\begin{defbox}[Auxiliary functions over plans]
$\mathit{Plan}$ --- the set of all plans over $\mathit{Ac}$;
$\mathit{pre}(\pi)$ --- the precondition of a plan; $\mathit{body}(\pi)$ --- the body of a plan
(the sequence of actions); $\mathit{empty}(\pi)$ --- is the plan empty?;
$\mathit{execute}(\pi)$ --- executes all actions of the plan;
$\mathit{hd}(\pi)$ --- the first action; $\mathit{tail}(\pi)$ --- the rest;
$\mathit{sound}(\pi,I,B)$ --- plan $\pi$ is correct with respect to intentions $I$ and beliefs $B$.

The agent's planning function:
$\mathit{plan} : 2^{\mathit{Bel}} \times 2^{\mathit{Int}} \times 2^{\mathit{Ac}} \to \mathit{Plan}$.
\end{defbox}

\textbf{The plan library.} The agent need not construct plans online --- that is expensive. Very
often $\mathit{plan}()$ is implemented as a \emph{library of ready-made plans}: it suffices to
traverse it once and check whether the preconditions of a plan match the current beliefs and
whether the effects of the plan match the goal. This idea is the core of the PRS architecture
(section~\ref{ssec:prs}).

\subsection{Implementing a BDI agent}

\noindent\textbf{The main control loop of a BDI agent:}
\begin{algorithmic}[1]
\State $B \gets B_0;\quad I \gets I_0$
\While{$\mathbf{true}$}
  \State $v \gets \mathit{see}()$
  \State $B \gets \mathit{brf}(B,v)$
  \State $D \gets \mathit{options}(B,I)$
  \State $I \gets \mathit{filter}(B,D,I)$
  \State $\pi \gets \mathit{plan}(B,I,\mathit{Ac})$
  \While{$\neg\big(\mathit{empty}(\pi) \vee \mathit{succeeded}(I,B) \vee \mathit{impossible}(I,B)\big)$}
    \State $a \gets \mathit{hd}(\pi);\quad \mathit{execute}(a);\quad \pi \gets \mathit{tail}(\pi)$
    \State $v \gets \mathit{see}();\quad B \gets \mathit{brf}(B,v)$
    \If{$\mathit{reconsider}(I,B)$}
      \State $D \gets \mathit{options}(B,I);\quad I \gets \mathit{filter}(B,D,I)$
    \EndIf
    \If{$\neg\,\mathit{sound}(\pi,I,B)$}
      \State $\pi \gets \mathit{plan}(B,I,\mathit{Ac})$ \Comment{replanning}
    \EndIf
  \EndWhile
\EndWhile
\end{algorithmic}

The meaning of the predicates: $\mathit{succeeded}(I,B)$ --- the intentions $I$ hold given the
beliefs $B$; $\mathit{impossible}(I,B)$ --- the intentions $I$ cannot hold given $B$.
The inner loop therefore terminates when the plan is exhausted, the goal is achieved, or the goal
has become unachievable.

\subsection{Agent commitment}
\label{ssec:commitment}

\begin{defbox}[Commitment strategies towards intentions]
When should an agent drop a goal?
\begin{itemize}
\item \textbf{Blind commitment} (also \emph{fanatical}) --- the agent maintains an intention until
  it \emph{comes to believe that it has achieved it}. (It never gives up.)
\item \textbf{Single-minded commitment} --- it maintains an intention until it comes to believe
  that it has achieved it \emph{or} that it is no longer achievable.
\item \textbf{Open-minded commitment} --- an intention persists as long as the agent believes it is
  still achievable (and still desired) --- so it may also be dropped when the motivation changes.
\end{itemize}
\end{defbox}

\textbf{Commitment to a plan} versus \textbf{commitment to a goal}. The agent in the algorithm above
is committed to a single plan; it abandons it when it believes the goal is satisfied, that it is
unachievable, or when the plan is empty. But the harder question remains: when should the agent
\emph{reconsider the intention itself}?

The classical dilemma: deliberation \textbf{takes time} and the environment changes while it is
going on.
\begin{itemize}
\item An agent that \emph{does not reconsider} its intentions may cling to goals that can no longer
  be achieved.
\item An agent that reconsiders \emph{too often} is so busy reasoning that it gets nothing done.
\end{itemize}

\textbf{Meta-level control} (Wooldridge, Parsons, 1999): a function $\mathit{reconsider}()$ that is
\emph{computationally cheaper} than the reconsideration itself. The quality criterion:
$\mathit{reconsider}()$ works well if, whenever it proposes reconsideration, the agent
\emph{actually} changes its intention after deliberating (and vice versa).

\begin{defbox}[Bold vs.\ cautious agent]
A \textbf{bold agent} --- reconsiders its intentions only after completing the entire plan (it never
interrupts a plan for the sake of deliberation).

A \textbf{cautious agent} --- reconsiders after every action of the plan.

The \textbf{level of boldness} --- how many actions the agent performs between two deliberations.

The \textbf{rate of world change} --- how many times the environment changes during one cycle of the
agent.

\textbf{Agent effectiveness} $=$ the number of intentions achieved $/$ the number of all intentions.
\end{defbox}

\textbf{Experimental result (Kinny, Georgeff, in Tileworld):}
\begin{itemize}
\item If the world changes \emph{slowly}, \textbf{bold} agents are more effective --- deliberation
  would be needless overhead.
\item If the world changes \emph{rapidly}, \textbf{cautious} agents outperform them --- plans become
  obsolete quickly.
\end{itemize}
This is a concrete quantitative form of the general reactivity vs.\ proactiveness trade-off.

\subsection{The Procedural Reasoning System (PRS)}
\label{ssec:prs}

\begin{defbox}[PRS --- \emph{Procedural Reasoning System}]
Georgeff et al., 1980s, Stanford. \textbf{The first implementation of the BDI architecture} and one
of the most successful agent architectures in practice --- it has been reimplemented many times:
\textbf{AgentSpeak/Jason}, \textbf{JAM}, \textbf{JACK}, \textbf{Jadex}.

Practical deployments: OASIS (the air traffic control system in Sydney), SPOC (business process
organisation), SWARMM (an air combat simulator of the Australian air force).
\end{defbox}

\textbf{Components of a PRS agent:} a database of \emph{beliefs} (FOL atoms of the Prolog kind),
\emph{desires} (goals), a \emph{plan library}, a stack of \emph{intentions}, and an
\emph{interpreter} that ties everything together. Sensor data enter the beliefs, and the interpreter
issues actions.

\begin{defbox}[A plan in PRS]
The agent has a library of ready-made plans representing its \emph{procedural knowledge}.
\textbf{Nothing is planned from scratch}; plans are only \emph{selected} from the library. A plan
has three parts:
\begin{itemize}
\item \textbf{Goal} (also the \emph{invocation condition} / \emph{cue}) --- the condition that is to
  hold after the plan has been executed; it determines which goal the plan is suited to.
\item \textbf{Context} --- the condition necessary for launching the plan (it must be consistent
  with the beliefs).
\item \textbf{Body} --- the actions to be performed.
\end{itemize}
The body of a plan is \emph{not} merely a linear sequence of actions --- it may contain further
\emph{goals}: \uv{achieve $f$}, \uv{achieve $f$ or $g$}, \uv{keep achieving $f$ until $g$ occurs}.
This gives rise to the hierarchical decomposition of goals into subgoals.
\end{defbox}

\textbf{The run of the interpreter:}
\begin{enumerate}
\item Initialisation: the initial beliefs $B_0$ and a top-level goal.
\item The intention stack contains the current goals in various stages of elaboration.
\item The interpreter takes the intention from the top of the stack and searches the plan library
  for plans with a matching goal.
\item Among the plans found it selects those whose \emph{context} is consistent with the current
  beliefs --- these constitute the current \emph{options} / \emph{desires}.
\item \textbf{Deliberation} = selecting one intention from the options:
  \begin{itemize}
  \item The original PRS used \emph{meta-plans} --- plans about plans, which modified the agent's
    intentions. This turned out to be too complicated.
  \item A more practical variant: each plan carries a numerical \emph{rating} (its expected
    utility), and the plan with the highest value is selected.
  \end{itemize}
\item The selected plan is executed, which may lead to further intentions being pushed onto the
  stack.
\item If a plan fails, the agent selects another intention from the options and continues.
\end{enumerate}

\medskip
\noindent\textbf{An illustration --- Jason/AgentSpeak} (an agent that brings its owner beer from the
fridge):
\begin{quote}\small\ttfamily
\noindent available(beer,fridge).\ \ limit(beer,10).\ \ \ \ /* initial beliefs */\\
too\_much(B) :- .count(consumed(\_,\_,\_,B),Q) \& limit(B,L) \& Q > L.\ \ /* rule */\\[0.3em]
+!has(owner,beer)\ \ \ \ \ \ \ \ \ \ \ \ \ \ \ \ \ \ \ \ \ \ \ /* triggering event */\\
\hspace*{1em}: available(beer,fridge) \& not too\_much(beer)\ \ \ \ /* context */\\
\hspace*{1em}<- !at(robot,fridge); open(fridge); get(beer); close(fridge);\ \ /* body */\\
\hspace*{2em} !at(robot,owner); hand\_in(beer).\\[0.3em]
+!at(robot,P) : not at(robot,P) <- move\_towards(P); !at(robot,P).
\end{quote}
The notation \texttt{+!g : c <- b} reads: \uv{if the event of adding goal \texttt{g} occurs and the
context \texttt{c} holds, execute the body \texttt{b}}. The symbol \texttt{!} denotes an
\emph{achievement goal} (achieve), \texttt{+}/\texttt{-} the addition/removal of a belief or a goal,
and \texttt{.send} and \texttt{.count} are built-in actions.

\subsection{Exam summary}

\begin{itemize}
\item \textbf{Deductive agent}: $\mathit{DB} \in 2^L$, $\mathit{DB} \vdash_\rho \mathit{Do}(a)$;
  two-phase action selection (provable action $\to$ consistent action $\to$ nothing).
  Elegant but impractical; the transduction problem and calculative rationality.
\item \textbf{BDI}: two phases --- deliberation ($\mathit{brf}$, $\mathit{options}$,
  $\mathit{filter}$) and means-ends reasoning (planning). Beliefs are subjective and fallible,
  desires may be inconsistent, intentions are commitments that persist and constrain further
  deliberation.
\item \textbf{STRIPS}: action descriptor $\langle P_a,D_a,A_a\rangle$, planning problem
  $\langle B_0,O,G\rangle$, $B_i = (B_{i-1}\setminus D_{a_i})\cup A_{a_i}$;
  a plan is \emph{admissible} ($B_{i-1}\models P_{a_i}$) vs.\ \emph{correct} (moreover
  $B_n \models G$). Be able to work through the blocks world.
\item \textbf{The BDI loop} and the role of $\mathit{reconsider}$ and $\mathit{sound}$.
\item \textbf{Commitment}: blind / single-minded / open-minded; bold vs.\ cautious agent;
  effectiveness $=$ achieved/all intentions; slow environment $\to$ bold, fast $\to$ cautious.
\item \textbf{PRS}: the first BDI implementation; a plan = Goal + Context + Body; a plan library
  instead of planning; the intention stack; deliberation by meta-plans or by utility; descendants
  AgentSpeak/Jason, JAM, JACK, Jadex.
\end{itemize}

\section{Reactive and hybrid architectures}
\label{sec:reaktivni}

\subsection{The reactive approach and its premises}

During the 1980s and 1990s it became clear that minor adjustments to the symbolic approach were not
enough, and alternative AI paradigms arose. Their common denominator is the \textbf{rejection of
symbolic representation} and of reasoning based on syntactic manipulation. Three key theses:

\begin{defbox}[Three principles of the reactive approach]
\textbf{Situatedness} --- intelligent behaviour depends on the environment in which the agent is
\emph{situated}; it cannot be studied in isolation.

\textbf{Embodiment} --- intelligent behaviour is not just logic, it is the product of an agent that
\emph{has a body} and interacts physically with the world.

\textbf{Emergence} --- intelligent behaviour \emph{emerges} from the interaction of many simple
behaviours; it is nowhere explicitly programmed.
\end{defbox}

\textbf{Characteristics of reactive agents:} they are \emph{behavioural} (the emphasis is on
developing and combining individual behaviours), \emph{situated}, and \emph{reactive} (the agent
essentially only reacts to the environment, it does not deliberate). The representation is
\textbf{subsymbolic}: connectionist networks, finite automata, simple IF--THEN rules.

\subsection{Symbolic reactive planning: the subsumption architecture}

\begin{defbox}[Subsumption architecture]
Rodney Brooks, 1986--1991; the most successful reactive approach. Brooks's three theses:
\begin{enumerate}
\item Intelligent behaviour can be generated \emph{without symbolic representation}.
\item Intelligent behaviour can be generated \emph{without explicit abstract reasoning}.
\item Intelligence is an \emph{emergent property} of certain complex systems.
\end{enumerate}

The intelligence of an agent is realised by a set of simple goal-oriented \textbf{behaviours}, for
which the following holds:
\begin{itemize}
\item Each behaviour is itself an action selection mechanism (a \emph{situation $\to$ action} pair).
\item Each behaviour receives percepts and transforms them into actions, is responsible for one
  partial goal, and has the simple structure of a rule.
\item Behaviours run \emph{in parallel} and \emph{compete} with one another for control of the agent.
\item Behaviours are arranged into a \textbf{subsumption hierarchy}, which defines their priorities
  (a higher layer \uv{subsumes} = suppresses a lower one).
\end{itemize}
\end{defbox}

\noindent\textbf{The action selection mechanism in the subsumption architecture:}
\begin{algorithmic}[1]
\Function{action}{percept $p$}
  \State $\mathit{fired} \gets \{\, b \in \mathit{Behaviours} \mid p \text{ satisfies the condition of } b \,\}$
    \Comment{behaviours that \uv{fire}}
  \ForAll{$b \in \mathit{fired}$}
    \If{there is no $b' \in \mathit{fired}$ such that $b' \prec b$}
      \Comment{$b$ is inhibited by nobody}
      \State \Return the action belonging to behaviour $b$
    \EndIf
  \EndFor
  \State \Return \textbf{null} \Comment{nothing fired, the agent does nothing}
\EndFunction
\end{algorithmic}

Formally: a behaviour is a pair $\langle c, a\rangle$, where $c \subseteq \mathit{Per}$ is a
condition and $a \in A$ an action; on the set of behaviours $R$ an inhibition relation
$\prec\ \subseteq R \times R$ (a strict total order) is defined. Mind the direction ---
in Wooldridge and in the lecture course, $b_1 \prec b_2$ means that \textbf{$b_1$ inhibits $b_2$},
that is, $b_1$ is \emph{lower} in the subsumption hierarchy and takes \emph{precedence}.
The notation $b_1 \prec b_2 \prec \dots \prec b_n$ thus lists behaviours from the highest priority
to the lowest.

\textbf{Properties:} the mechanism is \emph{simple} (although with dozens of behaviours the
priorities tend to be tricky) and \emph{very fast} --- it can be implemented in hardware and has
constant complexity.

\medskip
\noindent\textbf{Example: Steels's Mars explorer.}
The goal is to collect rare rock samples on a distant planet using a swarm of robotic explorers.
The means: the base emits a navigation signal (the robots need to do nothing more than detect the
\emph{gradient} of the signal); each robot carries radioactive crumbs for \emph{indirect
communication} (stigmergy) with the others. Direct communication is not needed.

\emph{First iteration} --- a random walk and the return of a single robot (priority decreases from
top to bottom):
\begin{center}\small
\begin{tabular}{ll}
\toprule
R1 & \textbf{if} I see an obstacle \textbf{then} change direction \\
R2 & \textbf{if} I carry a sample \textbf{and} I am at the base \textbf{then} drop the sample \\
R3 & \textbf{if} I carry a sample \textbf{and} I am not at the base \textbf{then} follow the gradient uphill \\
R4 & \textbf{if} I see a sample \textbf{then} pick it up \\
R5 & \textbf{if} true \textbf{then} move at random \\
\bottomrule
\end{tabular}
\end{center}

\emph{Second iteration} --- better exploration by means of stigmergy. Samples often occur in
clusters; a robot that has found one \uv{tells} the others by dropping crumbs on its way home:
\begin{center}\small
\begin{tabular}{ll}
\toprule
R6 & \textbf{if} I carry a sample \textbf{and} I am at the base \textbf{then} drop the sample \\
R7 & \textbf{if} I carry a sample \textbf{and} I am not at the base \textbf{then} drop 2 crumbs
     and follow the gradient uphill \\
R8 & \textbf{if} I sense a crumb \textbf{then} pick up 1~crumb and follow the gradient downhill \\
\bottomrule
\end{tabular}
\end{center}
Priorities (in the notation \uv{the left one inhibits the right one}, i.e.\ from the highest
priority): $\mathrm{R1} \prec \mathrm{R6} \prec \mathrm{R7} \prec \mathrm{R4} \prec
\mathrm{R8} \prec \mathrm{R5}$. A robot following a trail picks up one crumb at a time, so the trail
gradually \emph{evaporates} once the cluster has been exhausted. The swarm exhibits cooperative
behaviour even though no robot has a model of the world or sends any messages.

\subsection{Connectionist reactive planning: the agent network architecture}

The subsumption architecture resolves conflicts by \emph{fixed} priority, which is rigid. The
connectionist approach replaces fixed priorities with \textbf{continuous activation spread through a
network} --- action selection is the result of the network's dynamics, not of a fixed ordering. This
is \uv{reactive planning}: the agent behaves as if it were planning, but it never constructs an
explicit plan.

\begin{defbox}[Agent network architecture / \emph{spreading activation networks}]
Pattie Maes, 1989--1991. The agent is a set of \textbf{competence modules}, similar to Brooks's
behaviours. Each module $i$ has:
\begin{itemize}
\item \textbf{preconditions} $c_i$ --- what must hold for the module to be \emph{executable};
\item \textbf{effects}: an \emph{add list} $a_i$ and a \emph{delete list} $d_i$ (as in STRIPS);
\item an \textbf{activation level} $\alpha_i$ and a global \textbf{activation threshold} $\theta$.
\end{itemize}
The modules are connected in a directed graph on the basis of their conditions --- a match of pre-
and postconditions forms edges, to which further edges representing temporal succession and
conflicts are added. At each step activation is propagated; \textbf{the executable module with the
highest activation is selected}, provided it exceeds the threshold $\theta$.
\end{defbox}

\textbf{Three sources and three paths of activation spreading:}
\begin{itemize}
\item \textbf{From the state of the world}: modules whose preconditions are currently satisfied
  receive activation --- \uv{this can be done right now}.
\item \textbf{From goals}: modules whose add list contains a goal receive activation ---
  \uv{this leads to the goal}.
\item \textbf{From protected goals}: modules whose delete list contains an already satisfied goal are
  \emph{inhibited} --- \uv{this would destroy what I already have}.
\end{itemize}
Activation further spreads between modules along three types of link:
\begin{itemize}
\item \textbf{Successor links}: module $x$ sends activation to module $y$ if $y$ adds something that
  $x$ needs as a precondition --- \emph{forward} chaining (\uv{once I am done, things can go on}).
\item \textbf{Predecessor links}: a module $x$ that is not executable sends activation to modules
  that would satisfy its missing precondition --- \emph{backward} chaining (\uv{help me so that I can
  run}).
\item \textbf{Conflictor links}: module $x$ \emph{inhibits} module $y$, which destroys its
  preconditions.
\end{itemize}
The behaviour of the network is governed by five global parameters: $\theta$ (the activation
threshold, the only threshold in the system), $\phi$ (the activation injected into the network by the
state of the world), $\gamma$ (the activation injected by goals), $\delta$ (the activation removed by
protected goals), and $\pi$ (the \emph{average} activation level to which the network is normalised
after each step, so that the total energy does not change).
The ratio $\gamma/\phi$ expresses precisely the \textbf{proactiveness vs.\ reactivity} trade-off;
$\theta$ determines \textbf{deliberateness vs.\ speed} (a high threshold $\to$ the agent hesitates
longer but decides better). If in a given step no executable module exceeds the threshold, $\theta$
is lowered by 10~\% and the attempt is repeated --- the network therefore cannot \uv{freeze}.

\textbf{How the activation in one step is computed.} The update proceeds in four phases:
(1)~\emph{input from the environment} --- each true fact distributes a quantum $\phi$ evenly among
the modules that have it among their preconditions; (2)~\emph{input from goals} --- each goal
distributes $\gamma$ among the modules that add it, and each protected goal \emph{removes} $\delta$
from the modules that would destroy it; (3)~\emph{spreading between modules} --- an executable module
sends part of its activation forward to its successors (again divided evenly among the recipients), a
non-executable module sends activation backward to modules that can satisfy its missing precondition,
and conflictor links subtract the corresponding share of activation; (4)~\emph{normalisation} --- the
activations of all modules are rescaled so that their average is $\pi$ (the total energy of the
network remains constant). Then the executable module with activation above the threshold $\theta$ is
selected; the module that has been executed resets its activation to zero.

\emph{A note on the lecture course:} the slides present the activation threshold as a property of an
individual module (and speak of it as a priority); in Maes's original model the threshold $\theta$ is
\emph{global} and the role of the \uv{priority} is played by the current activation level
$\alpha_i$, which changes dynamically.

\medskip\noindent
\textbf{Comparison with subsumption:} the ordering of behaviours is not fixed but arises
\emph{dynamically} from the current state of the world and the goals. The network is thus able to
generate even multi-step \uv{plan-like} sequences (activation chains backwards from the goal along
predecessor links) without any explicit plan data structure ever being created --- hence the name
\emph{reactive planning}. The price is sensitivity to the parameter settings and the impossibility of
guaranteeing optimality or completeness.

\subsection{Limitations of reactive architectures}

\begin{itemize}
\item Reactive agents \textbf{have no model of the world}; everything must be derived from the
  current percepts. Sometimes they therefore have to \emph{change the environment} in order to
  \uv{remember} a piece of information (the radioactive crumbs) --- externalised memory, stigmergy.
\item They have only a \textbf{short-term view}: they act according to the current state and local
  information, and it is difficult to take global conditions and long-term goals into account.
\item Emergence \textbf{is not a good engineering practice}: behaviours have to be tuned
  experimentally, and both a methodology and any possibility of verification are missing.
\item The design is manageable for a few layers; \textbf{with dozens of behaviours it becomes
  intractable} (the number of interactions grows quadratically).
\end{itemize}

\subsection{Hybrid architectures}

Neither a purely reactive nor a purely deliberative architecture is ideal. Hybrid architectures
combine both:
\begin{itemize}
\item A \textbf{deliberative / planning component} works at the symbolic level, creating
  representations and plans.
\item A \textbf{reactive component} provides immediate responses without complex computation.
\end{itemize}
The components are usually arranged into \textbf{layers} (layered architectures), with the reactive
layers as a rule taking \emph{precedence} over the deliberative ones (safety before optimality).

\begin{defbox}[Horizontal vs.\ vertical layering]
\textbf{Horizontal layering:} all layers are connected \emph{in parallel} to both the sensors and
the effectors. Each layer is like a separate agent and proposes an action.
\begin{itemize}
\item[$+$] Conceptually simple: for $n$ kinds of behaviour, $n$ layers suffice.
\item[$-$] The layers influence one another and their proposals conflict; a \textbf{mediator
  function} is needed to resolve the conflict. It is the bottleneck: it has to consider $m^n$
  combinations ($n$ layers, $m$ possible proposals) and is a critical point of failure.
\end{itemize}

\textbf{Vertical layering:} the layers are connected \emph{in series}.
\begin{itemize}
\item \emph{One-pass}: a percept enters the lowest layer, passes upwards, and the action emerges from
  the topmost layer. A natural ordering and hierarchy of behaviours.
\item \emph{Two-pass}: percepts travel \emph{bottom-up}, control commands (actions)
  \emph{top-down}. The flow resembles management in a real company.
\item[$+$] The need for mediation disappears, and the complexity of the interactions is linear
  ($n-1$ interfaces).
\item[$-$] It is not robust: the failure of any one layer brings down the whole agent.
\end{itemize}
\end{defbox}

\begin{defbox}[TouringMachines (Ferguson, 1992)]
A \textbf{horizontal} three-layer architecture for a simulated agent in a traffic environment.
\begin{itemize}
\item The \textbf{modelling layer} --- symbolic models of the agent, the environment, and the other
  agents; it predicts and resolves conflicts, sets goals, and passes them to the planning layer.
\item The \textbf{planning layer} --- proactive behaviour; it selects from pre-prepared plans
  (similarly to PRS) and assembles a route plan from them.
\item The \textbf{reactive layer} --- classical reactive \emph{situation $\to$ action} rules, a fast
  immediate response (obstacle avoidance).
\end{itemize}
The layers work in parallel and their outputs are unified by a set of \textbf{control rules} --- a
concrete form of the mediator function; these can suppress the inputs and outputs of the layers
(\emph{censor} and \emph{suppressor} rules).
\end{defbox}

\begin{defbox}[InteRRaP (Müller, 1994)]
A \textbf{vertical two-pass} three-layer architecture.
\begin{itemize}
\item The \textbf{cooperative planning layer} --- social interaction, joint plans with other agents.
\item The \textbf{local planning layer} --- ordinary individual planning.
\item The \textbf{behaviour-based layer} --- reactive behaviour.
\end{itemize}
Each layer has \emph{its own knowledge base} at the corresponding level of abstraction (knowledge
about the world / plans / joint plans and models of the other agents). The layers communicate by two
operations: \emph{bottom-up activation} (a lower layer cannot handle the situation and passes it
upwards) and \emph{top-down execution} (a higher layer uses a lower one to carry out its decisions).
\end{defbox}

\medskip
\noindent\textbf{A real-world example: Stanley} (a Volkswagen Touareg R5, Stanford). The autonomous
vehicle that won the DARPA Grand Challenge 2005 (212~km across the Mojave Desert); a direct
predecessor of today's autonomous vehicles. A layered architecture:
\begin{itemize}
\item the sensor layer $\to$ the abstract perception layer $\to$ the planning and control layer
  (the route plan and vehicle control) $\to$ the vehicle interface layer;
\item plus a user interface layer (the panel, starting) and a global services layer (file system,
  communication, clock).
\end{itemize}

\subsection{Complex architectures: IDA and the global workspace}
\label{ssec:ida}

The last family of architectures goes in the opposite direction to reactive minimalism: it attempts
to integrate \emph{everything} into a single agent --- action selection, memory, deliberation,
emotions.

\begin{defbox}[IDA --- \emph{Intelligent Distribution Agent} (S.~Franklin)]
One of the most complex agent architectures ever built. It arose for a real task of the US Navy ---
assigning personnel to new billets (hence \uv{distribution}); the architecture itself is distributed
as well.

\textbf{The starting point: global workspace theory} (Baars, 1988--97) --- one of the theories of
consciousness:
\begin{itemize}
\item The mind \emph{is} a multiagent system: it consists of many simple, mostly unconscious
  processes performing specialised operations.
\item The processes communicate \emph{rarely} and only through a shared \emph{blackboard}-type
  memory, organised as an associative array (cf.\ chapter~\ref{sec:cdps}).
\item The processes dynamically form higher-order \textbf{coalitions}.
\item \textbf{The action selection mechanism:} the coalition that best matches the current inputs
  reaches \uv{consciousness} (the global workspace) and is executed.
\item There is a \emph{hierarchy of contexts} representing the world at various levels --- the
  context of goals, of percepts, of the senses, of concepts, and the cultural context.
\end{itemize}
\textbf{Assessment:} it combines many approaches from MAS and from the rest of AI (action selection,
memory, deliberation, emotions), and that is precisely its weakness --- many of the decisions are
\emph{ad hoc}, and tuning such a heterogeneous collection of models into an optimally cooperating
whole is very difficult.
\end{defbox}

For comparing action selection mechanisms, IDA is useful as a \emph{fourth} type of arbitration:
subsumption decides by \textbf{fixed priority}, the Maes network by \textbf{activation level},
layered architectures by a \textbf{mediator function}, and IDA by the \textbf{competition of
dynamically arising coalitions} for access to the global workspace.

\subsection{Comparison of architectures}

\begin{center}
\small
\begin{tabular}{@{}p{0.15\textwidth}
  >{\raggedright\arraybackslash}p{0.19\textwidth}
  >{\raggedright\arraybackslash}p{0.19\textwidth}
  >{\raggedright\arraybackslash}p{0.17\textwidth}
  >{\raggedright\arraybackslash}p{0.17\textwidth}@{}}
\toprule
& \textbf{Deductive} & \textbf{BDI / PRS} & \textbf{Reactive} & \textbf{Hybrid} \\
\midrule
Representation & symbolic (FOL) & symbolic (beliefs, plans) & subsymbolic / rules
  & mixed, per layer \\
Action selection & proving $\mathit{Do}(a)$ & deliberation + plan from a library
  & priority / activation & layers + arbitration \\
World model & complete & partial, revisable & none & per layer \\
Speed & very low & medium & high (const.) & high in reflex \\
Reactivity & poor & medium (tunable) & excellent & good \\
Proactiveness & good & excellent & none & good \\
Design & difficult but formal & systematic & experimental & complicated \\
\bottomrule
\end{tabular}
\end{center}

\subsection{Exam summary}

\begin{itemize}
\item The reactive approach rests on the triple \emph{situatedness --- embodiment --- emergence};
  it rejects symbolic representation.
\item \textbf{Subsumption} (Brooks): a set of simple \emph{situation~$\to$~action} behaviours running
  in parallel, arranged in a fixed priority order. Be able to present Steels's Mars explorer and the
  role of stigmergy (the crumbs).
\item \textbf{Agent network architecture} (Maes) = connectionist reactive planning:
  modules with precondition, add, and delete lists; activation spreads along successor,
  predecessor, and conflictor links, and the sources of activation are the state of the world,
  goals, and inhibition from protected goals; the executable module with the highest activation
  above the threshold is selected. The parameters $\gamma/\phi$ tune proactiveness vs.\ reactivity.
\item Limitations of reactive agents: no world model, short-term view, non-engineering design, an
  unscalable number of layers.
\item \textbf{Hybrid}: horizontal (parallel layers + a mediator, $m^n$ combinations, a bottleneck)
  vs.\ vertical (serial, one-pass/two-pass, $n-1$ interfaces, fragile).
\item \textbf{TouringMachines} = horizontal, 3~layers (modelling, planning, reactive) + control
  rules. \textbf{InteRRaP} = vertical two-pass, 3~layers (cooperative planning, local planning,
  behaviour-based) + a separate KB per layer. \textbf{Stanley} as a real-world deployment.
\item \textbf{IDA} (Franklin) + \emph{global workspace theory} (Baars): the mind as a MAS of simple
  unconscious processes communicating through a blackboard; coalitions of processes compete for entry
  into \uv{consciousness} and the winner is executed; a hierarchy of contexts. The most complex
  architecture, with many ad hoc decisions.
\end{itemize}

\section{The pathfinding problem}
\label{sec:cesta}

\subsection{Formulation and context}

Pathfinding is the most common concrete subtask a situated agent has to solve --- whether as part of
the planning layer of a hybrid architecture, as the action \texttt{move\_towards(P)} in a BDI plan, or
as a service provided to other agents.
In a multiagent context two further specifics appear that the single-agent version does not know:
the environment is \textbf{dynamic} (the map changes, obstacles appear) and the other agents are
\textbf{moving obstacles} with which it is necessary to coordinate.

\begin{defbox}[The pathfinding problem]
We are given a weighted graph $G = (V,E)$ with non-negative edge weights $w : E \to \mathbb{R}_0^+$,
an initial vertex $s \in V$ and a goal vertex $g \in V$ (or a set of goals $G_{\mathrm{goal}}$).
We seek a path $s = v_0, v_1, \dots, v_k = g$ minimising $\sum_{i} w(v_{i-1},v_i)$.

The graph is usually given \emph{implicitly} by a successor function (a 4-/8-connected grid, a
navigation mesh, a visibility graph, a grid of the robot's configuration space).
\end{defbox}

\subsection{A* and related algorithms}

\emph{Best-first} search maintains for every vertex $n$ the values
\begin{align*}
g(n) &= \text{the cost of the best path } s \to n \text{ found so far}, \\
h(n) &= \text{a heuristic estimate of the cost } n \to g_{\mathrm{goal}}, \\
f(n) &= g(n) + h(n) \quad \text{--- an estimate of the cost of the best path } s \to g_{\mathrm{goal}}
  \text{ through } n .
\end{align*}

\noindent\textbf{The A* algorithm} (Hart, Nilsson, Raphael, 1968):
\begin{algorithmic}[1]
\State $\mathit{OPEN} \gets \{s\}$ (a priority queue ordered by $f$), $\mathit{CLOSED} \gets \emptyset$
\State $g(s) \gets 0$, $f(s) \gets h(s)$, all other $g \gets \infty$
\While{$\mathit{OPEN} \ne \emptyset$}
  \State $n \gets \arg\min_{x \in \mathit{OPEN}} f(x)$; remove $n$ from $\mathit{OPEN}$
  \If{$n$ is a goal} \Return the reconstructed path \EndIf
  \State add $n$ to $\mathit{CLOSED}$
  \ForAll{$m \in \mathit{succ}(n)$}
    \State $g_{\mathrm{new}} \gets g(n) + w(n,m)$
    \If{$g_{\mathrm{new}} < g(m)$}
      \State $g(m) \gets g_{\mathrm{new}}$; $\mathit{parent}(m) \gets n$;
             $f(m) \gets g(m) + h(m)$
      \State insert/update $m$ in $\mathit{OPEN}$ (and remove it from $\mathit{CLOSED}$ if it was there)
    \EndIf
  \EndFor
\EndWhile
\State \Return \uv{no path exists}
\end{algorithmic}

\begin{defbox}[Admissibility and consistency of a heuristic]
A heuristic $h$ is \textbf{admissible} if it never overestimates:
$h(n) \le h^*(n)$ for all $n$, where $h^*(n)$ is the true cost of the cheapest path from $n$ to the
goal.

A heuristic is \textbf{consistent} (monotone) if for every edge $(n,m)$ the triangle inequality
holds:
\[
h(n) \le w(n,m) + h(m), \qquad h(g_{\mathrm{goal}}) = 0 .
\]
\textbf{Properties:} consistency $\Rightarrow$ admissibility. With an admissible heuristic, A* is
\emph{optimal} (it finds the cheapest path). With a consistent heuristic, $f$ is moreover
non-decreasing along every path and every vertex is expanded \emph{at most once} (there is no need to
reopen vertices from $\mathit{CLOSED}$). A* is also \emph{optimally efficient}: no other optimal
algorithm using the same heuristic expands fewer vertices. This claim has to be taken with a grain of
salt, however: it holds for vertices with $f(n) < C^*$ (where $C^*$ is the cost of the optimum) ---
these must be expanded by every optimal algorithm --- whereas for vertices with $f(n) = C^*$ it
depends on tie-breaking, and a luckier strategy may expand fewer of them.
\end{defbox}

\textbf{Special cases and variants:}
\begin{itemize}
\item $h \equiv 0$ $\Rightarrow$ A* degenerates into \textbf{Dijkstra's algorithm}
  (that is, uniform-cost search).
\item Typical heuristics on a grid: Manhattan $|\Delta x| + |\Delta y|$ (movement in 4~directions),
  octile $\max + (\sqrt{2}-1)\min$ (8~directions), Euclidean (an arbitrary angle).
\item \textbf{Weighted A*}: $f = g + \varepsilon h$, $\varepsilon > 1$ --- faster, but
  \emph{$\varepsilon$-suboptimal} (the path found is at most $\varepsilon$ times more expensive).
\item \textbf{IDA*} --- iterative deepening on $f$; memory $O(d)$ instead of exponential.
\item \textbf{JPS} (\emph{Jump Point Search}) --- on uniform grids it skips symmetric paths, giving an
  order-of-magnitude speed-up without loss of optimality.
\item \textbf{Theta*} / \emph{any-angle} --- allows edges off the grid (line-of-sight checks), giving
  more realistic trajectories.
\item \textbf{Hierarchical} planning (HPA*) --- the map is divided into clusters and the path is first
  sought in an abstract graph.
\end{itemize}

\medskip
\noindent\textbf{A small example.} A $3\times3$ grid, movement in 4~directions, edge cost $1$,
start $(0,0)$, goal $(2,2)$, an obstacle at $(1,1)$; the Manhattan heuristic
$h(x,y) = |2-x| + |2-y|$ is both admissible and consistent. The start has $f=0+4=4$; its successors
$(1,0)$ and $(0,1)$ have $g=1$, $h=3$, $f=4$; then $(2,0)$ and $(0,2)$ with $g=2,h=2,f=4$, then
$(2,1)$ and $(1,2)$ with $g=3,h=1,f=4$, and finally the goal $(2,2)$ with $g=4,h=0,f=4$. All vertices
on the optimal path have $f=4$ --- a typical manifestation of a consistent heuristic that estimates
the remainder of the path exactly.

\subsection{Pathfinding in a dynamic environment}

An agent typically moves and discovers along the way that the map is different from what it thought
(a new obstacle, a changed cost). Replanning with A* from scratch at every step is expensive. The
solution is \textbf{incremental} search algorithms, which upon a local change recompute only the
affected part.

\begin{defbox}[LPA*, D* and D*~Lite]
\textbf{LPA*} (\emph{Lifelong Planning A*}, Koenig, Likhachev, 2001) --- an incremental version of A*
for a \emph{fixed} start and goal. In addition to $g(n)$ it maintains for each vertex the value
\emph{rhs}$(n)$ (\emph{right-hand side}), a \emph{lookahead} value derived from the $g$ values of its
predecessors:
\[
\mathit{rhs}(n) = \begin{cases}
0 & n = s,\\
\min_{p \in \mathit{pred}(n)} \big(g(p) + w(p,n)\big) & \text{otherwise.}
\end{cases}
\]
A vertex is \textbf{locally consistent} if $g(n) = \mathit{rhs}(n)$. If the cost of an edge changes,
some vertices become inconsistent, are placed into a priority queue, and the algorithm propagates the
change only where it is needed.

\textbf{D*} (\emph{Dynamic A*}, Stentz, 1994) --- the original algorithm for a robot that moves and
discovers obstacles. It searches \emph{from the goal towards the start} (so that the values remain
usable when the robot moves).

\textbf{D*~Lite} (Koenig, Likhachev, 2002) --- LPA* searched from the goal, with a correction of the
priority queue as the robot moves. Functionally equivalent to D*, but substantially simpler; the
\emph{de facto} standard in mobile robotics. Orders of magnitude faster than repeated calls to A* ---
typically only a small neighbourhood of the change is recomputed.

\textbf{Other variants:} \emph{Anytime Repairing A*} (ARA*) --- quickly returns a suboptimal solution
and gradually improves it; \emph{AD*} --- a combination of anytime and dynamic replanning;
\emph{Field D*} --- interpolation for continuous motion.
\end{defbox}

\medskip
\noindent\textbf{And what about continuous space?} The preceding algorithms need a graph. In the
robot's \emph{continuous configuration space} (position $\times$ orientation $\times$ joints) the
graph is either produced by discretisation or by \emph{sampling}:
\begin{itemize}
\item \textbf{PRM} (\emph{probabilistic roadmap}) --- randomly sample free configurations, connect
  nearby pairs with a local planner, and search the resulting \uv{road network} in the classical way
  (A*). Suitable for repeated queries in a static map.
\item \textbf{RRT} (\emph{rapidly-exploring random tree}) and \textbf{RRT*} --- a tree grown
  incrementally towards random samples; single-shot queries, and it can also handle kinodynamic
  constraints. RRT is \emph{probabilistically complete} (it finds a solution with probability tending
  to 1), and RRT* is moreover \emph{asymptotically optimal}.
\item \textbf{Potential fields} --- the goal attracts, obstacles repel; very fast and reactive (close
  to subsumption), but they get stuck in \emph{local minima}.
\end{itemize}
What these planners have in common as far as MAS is concerned is that they tend to be the
\emph{local / low-level} layer of a hybrid architecture, while the coordination between agents (MAPF)
is solved over a discrete graph.

\textbf{An alternative: planning with free space.} In an unknown environment the
\emph{freespace assumption} is often used: unknown cells are considered traversable, the agent follows
the planned path, and as soon as it hits something it replans. This is exactly the scenario for which
D* was designed.

\subsection{Multi-agent path finding (MAPF)}

\begin{defbox}[Multi-Agent Path Finding, MAPF]
Input: a graph $G=(V,E)$ and $k$ agents, where agent $i$ has a start $s_i$ and a goal $g_i$ (all
distinct). Time is discrete; at each step every agent either traverses an edge or waits in place.

\textbf{Conflicts} --- a solution has to be \emph{conflict-free}:
\begin{itemize}
\item a \emph{vertex conflict} --- two agents are in the same vertex at the same time;
\item an \emph{edge / \uv{swap} conflict} --- two agents exchange their positions in a single step
  (they pass through each other);
\item depending on the variant also a \emph{following} conflict (an agent enters a vertex that its
  neighbour is just leaving) and \emph{cyclic} rotations.
\end{itemize}
\textbf{The objective function:} \emph{sum-of-costs} $\sum_i \mathit{cost}_i$, or \emph{makespan}
$\max_i \mathit{cost}_i$.

\textbf{Complexity:} finding an \emph{optimal} solution is NP-hard (for both objective functions);
deciding whether \emph{some} solution exists is polynomial on general graphs.
\end{defbox}

\textbf{Approaches to MAPF} split into coupled and decoupled ones:

\begin{center}
\small
\begin{tabular}{p{0.22\textwidth} p{0.36\textwidth} p{0.34\textwidth}}
\toprule
\textbf{Type} & \textbf{Principle} & \textbf{Properties} \\
\midrule
\textbf{coupled}
  & A* in the \emph{joint} state space $V^k$ (a joint state = the positions of all agents)
  & optimal and complete, but the state space grows exponentially with $k$; in practice up to
    a few agents \\
\textbf{decoupled}
  & each agent plans on its own, conflicts are resolved afterwards (priorities, reservations)
  & fast, scales well, but \emph{neither complete nor optimal} \\
\textbf{hybrid / search over conflicts}
  & searching the space of \emph{constraints} rather than of positions (CBS, ICTS, M*)
  & optimal and yet practically usable for dozens of agents \\
\textbf{reductions}
  & translation into SAT, ILP, ASP, or CSP and use of an off-the-shelf solver
  & optimal, very strong on dense instances \\
\bottomrule
\end{tabular}
\end{center}

\begin{defbox}[Prioritized planning]
The simplest decoupled approach (Erdmann, Lozano-Pérez, 1987):
\begin{enumerate}
\item Order the agents by some priority.
\item For $i = 1,\dots,k$: plan the path of agent $i$ with A* in \emph{space--time}
  (a state = a pair $\langle$vertex, time$\rangle$) so that it \emph{avoids} the already planned
  trajectories of agents $1,\dots,i-1$; these are stored in a \textbf{reservation table}.
\end{enumerate}
This variant is called \textbf{Cooperative A*} (Silver, 2005); the extension
\emph{Hierarchical Cooperative A*} uses an abstract heuristic ignoring the other agents, and
\emph{Windowed HCA*} plans only over a limited time window (necessary for real-time games).

\textbf{The weakness:} the algorithm is \emph{incomplete} --- there exist instances that are solvable
but not solvable under any ordering of priorities (classically, two agents in a narrow corridor where
one has to step aside into an alcove).
\end{defbox}

\begin{defbox}[Conflict-Based Search (CBS)]
Sharon et al., 2012--2015. A two-level optimal algorithm.

\textbf{The low level:} for each agent an optimal path in space--time is sought using A* while
respecting a set of \emph{constraints} of the form
$\langle a_i, v, t\rangle$ (\uv{agent $i$ must not be in vertex $v$ at time $t$}), or possibly edge
constraints $\langle a_i, (u,v), t\rangle$.

\textbf{The high level:} it searches a \emph{constraint tree} in best-first order by cost:
\begin{enumerate}
\item The root: an empty set of constraints, each agent has its individually optimal path.
\item Take the node with the lowest solution cost; check whether the paths are conflict-free.
  If they are, \textbf{return the solution} --- it is optimal.
\item Otherwise select the first conflict $\langle a_i, a_j, v, t\rangle$ and \emph{branch} into two
  children: add the constraint $\langle a_i, v, t\rangle$ to the first one and
  $\langle a_j, v, t\rangle$ to the second. In each child replan \emph{only} the agent concerned.
\end{enumerate}
\textbf{The key idea:} the joint state space $V^k$ is never searched; branching occurs only on
conflicts that have actually arisen, and these tend to be few. CBS is \emph{complete} and
\emph{optimal} (sum-of-costs).

\textbf{Improvements:} ICBS/CBSH (stronger heuristics for the high level, classification of conflicts
into cardinal/semi-cardinal/non-cardinal), \emph{meta-agent} CBS (agents with many mutual conflicts
are merged into a single meta-agent and solved in a coupled way), ECBS (\emph{bounded suboptimal},
orders of magnitude faster at the price of a factor of $1+\varepsilon$).
\end{defbox}

\medskip
\noindent\textbf{An example run of CBS.} A corridor $1 - 2 - 3$ with an alcove $4$ attached to $2$;
agent $a_1$ goes $1 \to 3$, agent $a_2$ goes $3 \to 1$.
\begin{itemize}
\item \textbf{The root} (an empty set of constraints): the individually optimal paths are
  $a_1: (1,2,3)$ and $a_2: (3,2,1)$, cost $2+2=4$. At time $t=1$ both are in vertex $2$ ---
  a vertex conflict $\langle a_1,a_2,2,1\rangle$.
\item \textbf{The first branching:} the left child adds the constraint $\langle a_1,2,1\rangle$, the
  right one $\langle a_2,2,1\rangle$; in each of them \emph{only} the agent concerned is replanned.
  On the left, therefore, $a_1 = (1,1,2,3)$ (it waits), and the cost of the node is $3+2 = 5$. But
  this creates a \emph{swap} conflict: between $t=1$ and $t=2$ agent $a_1$ traverses the edge
  $1\to2$ and $a_2$ the same edge $2\to1$. Symmetrically on the right.
\item \textbf{Further branching} on the edge conflict forces one of the agents to pull into the
  alcove. The optimal solution has cost $\mathbf{7}$: $a_1 = (1,1,2,3)$ for $3$ and
  $a_2 = (3,2,4,2,1)$ for $4$. (Neither cost $5$ nor $6$ exists: an agent going straight through
  vertex $2$ always causes either a vertex or a swap conflict --- try out all four combinations of
  paths of length~3.)
\end{itemize}
The example illustrates two typical properties of CBS: the cost of the nodes along a branch is
\emph{non-decreasing} ($4 \to 5 \to 7$) and \emph{both} agents have to give ground, not just the
\uv{yielding} one.
Without the alcove the instance is outright \textbf{unsolvable} and basic CBS would branch forever ---
the completeness of CBS is guaranteed only for solvable instances, and unsolvability has to be tested
separately.

\subsection{Exam summary}

\begin{itemize}
\item A*: $f = g+h$; \emph{admissibility} $h \le h^*$ $\Rightarrow$ optimality;
  \emph{consistency} $h(n) \le w(n,m)+h(m)$ $\Rightarrow$ every vertex is expanded once.
  $h\equiv 0$ gives Dijkstra; weighted A* is $\varepsilon$-suboptimal.
\item Dynamic environment: LPA* (the values $g$ and \emph{rhs}, local consistency
  $g=\mathit{rhs}$), D* and D*~Lite (search from the goal, incremental repair),
  the freespace assumption.
\item MAPF: the definition, vertex and edge (swap) conflicts, sum-of-costs vs.\ makespan;
  optimal MAPF is NP-hard.
\item Decoupled solution: prioritized planning / Cooperative A* with a reservation table,
  planning in space--time --- fast, but incomplete.
\item \textbf{CBS}: the low level = A* with constraints for a single agent, the high level = binary
  branching of the constraint tree on the conflict found; optimal and complete, and it avoids the
  exponential joint space.
\end{itemize}

\section{Knowledge and ontologies}
\label{sec:ontologie}

\subsection{From an agent to a multiagent system}

We now know how to design a single agent. For agents to form a system they have to communicate with
one another, and for that they need to resolve:
\begin{itemize}
\item \textbf{knowledge representation} --- in what terms knowledge is expressed;
\item \textbf{a shared vocabulary} --- so that the same word means the same thing (this is the task of
  \emph{ontologies});
\item \textbf{standard messages} --- the syntax and semantics of communication (\emph{ACL});
\item \textbf{predictable behaviour during communication} --- \emph{interaction protocols};
\item \textbf{technical problems} --- distribution, agent mobility, asynchrony, unreliable
  communication channels.
\end{itemize}

\subsection{Reasoning about knowledge: epistemic logic and common knowledge}

Before turning to the \emph{representation} of knowledge (ontologies), we should say how agents'
knowledge is formally \emph{reasoned about}.

\begin{defbox}[Epistemic logic and common knowledge]
An agent's knowledge is modelled by a modal operator $K_i\varphi$ (\uv{agent $i$ \emph{knows} that
$\varphi$}) with \textbf{Kripke possible-worlds semantics}: $K_i\varphi$ holds in world $w$ if and
only if $\varphi$ holds in \emph{all} worlds that agent $i$ cannot distinguish from $w$ (the
indistinguishability relation --- an equivalence).
\textbf{Knowledge} is axiomatised by the system \textbf{S5} (KT45); the key axiom is
\textbf{T}: $K_i\varphi \Rightarrow \varphi$ (what I know is true).
\textbf{Belief} (the operator $B_i$, cf.\ BDI and the semantics of FIPA SL) is axiomatised by the
system \textbf{KD45}, where axiom~T is replaced by the weaker \textbf{D}:
$\neg B_i\bot$ --- beliefs must be consistent, but they \emph{may be false}.
The difference between knowledge and belief is therefore formally exactly axiom~T.

For a group of agents $G$ one defines: $E_G\varphi$ --- \emph{everyone} in $G$ knows $\varphi$;
\textbf{common knowledge}
$C_G\varphi \equiv E_G\varphi \wedge E_G E_G\varphi \wedge E_G E_G E_G\varphi \wedge \dots$
(everyone knows that everyone knows that\dots); \emph{distributed knowledge} $D_G\varphi$
--- $\varphi$ follows from the union of the knowledge of all members of the group.

\textbf{The coordinated attack} (the two generals problem):
two generals negotiate a simultaneous attack via an unreliable messenger. No finite number of
delivered acknowledgements creates common knowledge of the agreement --- after $k$
messages at most $E_G^k\varphi$ holds, never $C_G\varphi$. The lesson for MAS:
\emph{unreliable communication never creates common knowledge}; perfect coordination therefore
requires either reliable channels or weaker goals (time-bounded or probabilistic variants of common
knowledge).

\emph{The problem of logical omniscience}: in Kripke semantics an agent automatically knows all the
logical consequences of its knowledge as well as all tautologies --- a real, computationally bounded
agent does not. This is the main objection against the straightforward use of epistemic logic as an
implementation tool (cf.\ the difficulties of deductive agents,
chapter~\ref{sec:deliberativni}).
\end{defbox}

\subsection{What an ontology is}

\begin{defbox}[Ontology]
\emph{Philosophically} (from the Greek \emph{to on} = that which is, \emph{logos} = word): the study
of being, according to Aristotle \uv{first philosophy}, a part of metaphysics.

\emph{In computer science}: an \textbf{explicit and formalised description of a part of reality}.
Usually a formal declarative description containing a \emph{glossary} (definitions of concepts) and a
\emph{thesaurus} (definitions of relations between concepts). An ontology is a kind of vocabulary for
storing and exchanging knowledge about a domain in a standard way.

\textbf{Gruber (1995):} \uv{An ontology is a description (like a formal specification of a program) of
the concepts and relationships that can formally exist for an agent or a community of agents.}
The frequently cited shortened version: \emph{an ontology is an explicit specification of a
conceptualisation}.
\end{defbox}

\textbf{A motivating example} (a dialogue between two people --- or two agents). \uv{Have you heard
7777? --- No, what is it? --- It's the new CD by the band SRPR, alt.rock with electronics and great
lyrics. --- I see.} For them to understand each other, they have to share the following:
7777 is an album; it falls into the genres alt.rock and electronica; it has tracks, a performer
(SRPR), and authors of the music and the lyrics; SRPR is a band, it has members, and they play
instruments\dots

\begin{defbox}[The formalism of an ontology]
\begin{itemize}
\item \textbf{Classes} (concepts) --- things that have something in common.
\item \textbf{Individuals} (instances, objects) --- concrete members of classes.
\item \textbf{Properties} --- attributes of classes and individuals.
\item \textbf{Relations between classes}, in particular the \textbf{subclass} (\emph{is-a}) relation
  --- a \emph{transitive} relation.
\item Further relations and properties (basic knowledge about the domain).
\end{itemize}
The \textbf{structural part} (classes, relations, axioms) is usually called the \emph{ontology}
proper; together with \textbf{facts about concrete individuals} it forms a \textbf{knowledge base}.
\end{defbox}

\textbf{An ontology of ontologies by strength of formalisation} (from weak to strong):
a vocabulary (a controlled vocabulary of selected terms) $\to$ a glossary (definitions of meanings in
natural language) $\to$ a thesaurus (synonyms) $\to$ an informal hierarchy (Amazon, wikis) $\to$
a formal \emph{is-a} hierarchy (class subsumption) $\to$ classes with properties $\to$
value restrictions (\uv{every person has exactly one mother}) $\to$ arbitrary logical constraints
(which leads to complex inference algorithms).

\textbf{An ontology of ontologies by scope:}
\begin{itemize}
\item \textbf{Application} ontologies --- the most common in practice, hard to reuse.
\item \textbf{Domain} ontologies --- a popular research subject in the semantic web.
\item \textbf{Upper ontologies} --- these are supposed to describe \uv{everything}
  (Thing, Living thing, \dots) and serve as a foundation for domain ontologies:
  BFO, GFO, UFO; WordNet and IDEAS are not suitable for formal inference.
  There are methodological objections against the very effort to build a universal ontology
  (Wittgenstein, \emph{Tractatus logico-philosophicus}).
\end{itemize}

\subsection{Ontology languages}

Ontology languages are formal declarative languages for knowledge representation; they contain facts
and inference rules and most often rest on first-order logic or on \emph{description logic}.

\textbf{Frames} --- a historical predecessor (Minsky), popular in classical AI and expert systems.
They correspond to object-oriented programming:

\begin{center}\small
\begin{tabular}{ll}
\toprule
\textbf{Frame terminology} & \textbf{Object-oriented terminology} \\
\midrule
Frame & class of objects \\
Slot & property / attribute of an object \\
Trigger, accessor, mutator & method \\
\bottomrule
\end{tabular}
\end{center}

\textbf{XML} was not created for ontologies, but it is used for them: its main advantage is the
possibility of defining one's own tags, which naturally represent a vocabulary:
\begin{quote}\small\ttfamily
<product type=\textquotedbl CD\textquotedbl>\\
\hspace*{1em}<title>7777</title>\\
\hspace*{1em}<artist>SRPR</artist>\\
</product>
\end{quote}
What it lacks, however, is a defined semantics --- it describes only the structure of a document, not
its meaning.

\begin{defbox}[RDF --- \emph{Resource Description Framework}]
The standard tool for knowledge representation for (but not only for) the web. It is
\textbf{simple}: not very expressive, but with fast inference algorithms.

Everything is expressed by \textbf{triples} \emph{subject -- predicate -- object}, that is,
$\mathit{Predicate}(\mathit{Subject},\mathit{Object})$. For example,
$\mathit{MarriedTo}(\mathit{Karel},\mathit{J\acute{a}ja})$,
$\mathit{FatherOf}(\mathit{Karel},\mathit{P\acute{e}\check{t}a})$.
A set of triples forms a directed labelled graph (an \emph{RDF graph}); resources are identified by
IRIs. Serialisations: RDF/XML, Turtle, N-Triples, JSON-LD. The query language is SPARQL.
The extension \textbf{RDFS} adds classes, \texttt{subClassOf}, \texttt{domain}, \texttt{range}.
\end{defbox}

\begin{defbox}[OWL --- \emph{Web Ontology Language}]
It too originates from the (semantic) web, and exists in versions 1 and~2. It is a \emph{collection of
several formalisms} for describing ontologies; the syntax may be XML, RDF, functional notation,
Manchester syntax, or Turtle. An attempt at a formal and yet practically usable language.

Its basis is \textbf{description logic} --- a \emph{decidable fragment} of first-order logic.

\textbf{The open world assumption} (OWA): what we do not know is \emph{undefined} --- as opposed to
the closed world assumption (databases, Prolog), where everything we do not know is \emph{false}.
\end{defbox}

\begin{defbox}[Description logic]
Three kinds of entity: \textbf{concepts} (classes), \textbf{roles} (properties, predicates), and
\textbf{individuals} (objects). An \textbf{axiom} is a logical expression about roles and concepts
--- unlike frames and object-oriented programming, which describe a class \emph{completely}, axioms
describe a class only \emph{partially} (and they may come from several sources).

It is a whole family of logical systems according to the constructors allowed. The basic
$\mathcal{ALC}$: negation, conjunction and disjunction of concepts, and restricted existential and
universal quantifiers. Extensions: role hierarchies ($\mathcal{H}$), transitivity
($\mathcal{S}$), inverse roles ($\mathcal{I}$), cardinality restrictions ($\mathcal{N}$,
$\mathcal{Q}$), nominals ($\mathcal{O}$).

A knowledge base is divided into the \textbf{T-Box} (the terminology --- axioms about concepts and
roles) and the \textbf{A-Box} (the assertional part --- facts about individuals).
\end{defbox}

\textbf{OWL profiles.} \emph{OWL~1 Lite} ($\mathcal{SHIF}$) --- the simplest, close to RDF, with many
restrictions for the sake of fast machine processing; \emph{OWL~1 DL} ($\mathcal{SHOIN}$) ---
corresponds to description logic (it makes it possible to express, e.g., the disjointness of classes),
is complete and decidable; \emph{OWL~1 Full} --- the greatest expressive power, but many problems are
undecidable. \emph{OWL~2} ($\mathcal{SROIQ}$) adds profiles: \emph{EL} (inference in polynomial time,
large medical ontologies such as SNOMED), \emph{QL} (queries over a knowledge base, connection to
relational databases), and \emph{RL} (axioms in the form of rules).

\begin{defbox}[KIF --- \emph{Knowledge Interchange Format}]
Knowledge representation in first-order logic, with a LISP-like prefix syntax. It was designed within
the DARPA \emph{Knowledge Sharing Effort} project as the \textbf{inner language} of messages (the
content), whereas KQML was the \emph{outer} language (the envelope). Examples:
\begin{quote}\small\ttfamily
(salary 015-46-3946 widgets 72000)\\
(> (* (width chip1) (length chip1)) (* (width chip2) (length chip2)))\\
(forall (?F ?T) (=> (and (instance ?F Farmer) (instance ?T Tractor)) (likes ?F ?T)))\\
(exists (?F ?T) (and (instance ?F Farmer) (instance ?T Tractor) (likes ?F ?T)))
\end{quote}
\end{defbox}

\textbf{DAML+OIL} --- DARPA Agent Markup Language + Ontology Interchange Language, the direct
predecessor of OWL, abandoned in 2006.

\subsection{Exam summary}

\begin{itemize}
\item Epistemic logic: $K_i\varphi$ + possible worlds; knowledge = S5 (axiom~T: what I know is
  true), belief = KD45 (consistency only). Common knowledge $C_G\varphi$
  (\uv{everyone knows that everyone knows\dots}) does not arise from unreliable communication
  (the coordinated attack); logical omniscience as the main weakness.
\item An ontology = an explicit formalised description of a part of reality (glossary + thesaurus);
  Gruber's definition. The formalism: classes, individuals, properties, relations (in particular the
  transitive \emph{is-a}), facts. Ontology + facts = a knowledge base.
\item The scale of formalisation from a vocabulary up to arbitrary logical constraints; the division
  into application, domain, and upper ontologies.
\item RDF = \emph{subject--predicate--object} triples, simple and fast; RDFS adds classes.
\item OWL = description logic (a decidable fragment of FOL), the \emph{open world assumption},
  T-Box vs.\ A-Box; the Lite/DL/Full profiles ($\mathcal{SHIF}$, $\mathcal{SHOIN}$),
  OWL~2 = $\mathcal{SROIQ}$ + the EL/QL/RL profiles.
\item KIF = FOL in LISP-like syntax, the content of a message; KQML = the envelope.
\end{itemize}

\section{Agent communication: speech acts, ACL, protocols}
\label{sec:komunikace}

\subsection{Why agent communication is different}

In object-oriented programming, \uv{communication} means calling a method of an object with
parameters --- the callee \emph{must} comply. Agents cannot force another agent to do something, nor
can they directly change its internal variables. They have to \textbf{perform a communicative act}
with the aim of
\begin{itemize}
\item exchanging information, or
\item \emph{influencing} other agents to do something.
\end{itemize}
The other agents have their own goals and it is entirely up to them what they do with the information,
request, or query they receive. This is the fundamental difference: \textbf{communication in MAS is an
action, not a call}.

\subsection{Speech act theory}

\begin{defbox}[Speech acts --- Austin, 1962]
Some parts of the use of language have the character of \textbf{actions}, because they change the
state of the world in the same way as physical actions do --- these are speech acts:
\begin{itemize}
\item \uv{I now pronounce you man and wife.}
\item \uv{I declare war.}
\end{itemize}
This is a \emph{pragmatic} theory of language: how speech is used to achieve goals.
Typical performative verbs: \emph{request}, \emph{inform}, \emph{promise}.
\end{defbox}

\begin{defbox}[The three layers of a speech act]
\begin{itemize}
\item The \textbf{locutionary act} --- \emph{what was said}; the utterance itself.
  \uv{Make me some tea.}
\item The \textbf{illocutionary act} --- \emph{what was meant};
  the locution + the performative meaning (a query, a request, \dots). \uv{She asked me for tea.}
\item The \textbf{perlocutionary act} --- \emph{what actually happened};
  the effect of the speech act. \uv{She got me to make her some tea.}
\end{itemize}
For MAS the \textbf{illocutionary} part is the crucial one --- that is the content of a message (the
performative).
\end{defbox}

\textbf{Searle} elaborated the conditions under which a speech act is \uv{felicitous}.
For the act \emph{SPEAKER requests HEARER to perform an action}:
\begin{itemize}
\item \textbf{Normal input/output conditions}: the hearer can hear, it is not a scene in a film\dots
\item \textbf{Preparatory conditions} (what must hold for the speaker to be able to choose this act):
  the hearer is able to perform the action; the speaker believes the hearer is able to perform it; it
  is not obvious that the hearer would perform it even without being asked.
\item \textbf{Sincerity}: the speaker genuinely wants the action to be performed.
\end{itemize}

\begin{defbox}[Searle's categories of speech acts]
The older division: \emph{request}, \emph{advice}, \emph{statement}, \emph{promise}, \dots

The newer (and standard) division into five classes:
\begin{itemize}
\item \textbf{Assertives} (representatives) --- they inform the hearer (\emph{inform}).
\item \textbf{Directives} --- they request the hearer to perform an action (\emph{request}).
\item \textbf{Commissives} --- the speaker commits to something (\emph{promise}).
\item \textbf{Expressives} --- the speaker expresses a mental state, an emotion
  (\uv{thank you}).
\item \textbf{Declarations} --- they change the state of affairs (war, marriage).
\end{itemize}
A speech act always has two components: a \textbf{performative verb} (request, query, inform)
and a \textbf{propositional content} (\uv{the window is closed}).
\end{defbox}

\begin{defbox}[The planning theory of speech acts (Cohen, Perrault, 1979)]
\uv{\dots modelling [speech acts] in a planning system as operators defined\dots{}
in terms of the beliefs and goals of speakers and hearers. Speech acts are thus treated in exactly the
same way as physical actions.}

The semantics of speech acts is therefore given in the spirit of \textbf{STRIPS} (preconditions --
delete -- add), except that modal operators for \emph{beliefs}, \emph{abilities}, and \emph{wants} are
used.
\end{defbox}

\noindent\textbf{Example: Request and Inform.}
\begin{center}\small
\begin{tabular}{p{0.16\textwidth} p{0.38\textwidth} p{0.38\textwidth}}
\toprule
& \textbf{Request(S,H,A)} & \textbf{Inform(S,H,$\varphi$)} \\
\midrule
\emph{cando} precondition
  & $(S\ \mathrm{believe}\ (H\ \mathrm{cando}\ A))\ \wedge$ \newline
    $(S\ \mathrm{believe}\ (H\ \mathrm{believe}\ (H\ \mathrm{cando}\ A)))$
  & $(S\ \mathrm{believe}\ \varphi)$ \\
\emph{want} precondition
  & $(S\ \mathrm{believe}\ (S\ \mathrm{want}\ \mathrm{requestInstance}))$
  & $(S\ \mathrm{believe}\ (S\ \mathrm{want}\ \mathrm{informInstance}))$ \\
Effect
  & $(H\ \mathrm{believe}\ (S\ \mathrm{believe}\ (S\ \mathrm{want}\ A)))$
  & $(H\ \mathrm{believe}\ (S\ \mathrm{believe}\ \varphi))$ \\
\bottomrule
\end{tabular}
\end{center}
Note that the effect of \emph{Inform} is \emph{not} \uv{the hearer comes to believe $\varphi$}, but
only \uv{the hearer comes to believe that the speaker believes $\varphi$}. An autonomous agent is not
obliged to believe what anyone tells it --- and that is precisely the difference from a method call.

\subsection{KQML and FIPA-ACL}

\begin{defbox}[KQML --- \emph{Knowledge Query and Manipulation Language}]
It arose in the 1990s within the DARPA \emph{Knowledge Sharing Effort} (KSE) project together with KIF.
\begin{itemize}
\item \textbf{KQML} = the \emph{outer} communication language (\uv{the envelope of the letter}); it
  carries the \emph{illocutionary} part of the message.
\item \textbf{KIF} = the \emph{inner} language; the propositional content of the message, knowledge
  representation.
\end{itemize}
A message is a list consisting of a \emph{performative} and named parameters:
\begin{quote}\small\ttfamily
(ask-one\\
\hspace*{1em}:content (PRICE IBM ?price)\\
\hspace*{1em}:receiver stock-server\\
\hspace*{1em}:language LPROLOG\\
\hspace*{1em}:ontology NYSE-TICKS)
\end{quote}
\textbf{Parameters:} content, force, reply-with, in-reply-to, sender, receiver, language, ontology.

\textbf{Performatives:} achieve, advertise, ask-about, ask-one, ask-all, ask-if,
break, sorry, error, broadcast, forward, recruit-one, recruit-all, reply, subscribe, tell, \dots
\end{defbox}

\textbf{Criticism of KQML:} a \emph{commit} performative was missing, the semantics was not precisely
defined, the set of performatives was unsystematic and needlessly large, and the transport mechanism
was not standardised.

\begin{defbox}[FIPA-ACL --- \emph{Agent Communication Language}]
FIPA (the \emph{Foundation for Intelligent Physical Agents}, today an IEEE standards committee)
designed ACL as a \textbf{simplification and refinement of KQML}: a smaller and more systematic set of
performatives and a \emph{formally defined semantics}. A practical implementation: the JADE platform.

\textbf{The structure of a message} (13 fields --- the performative and 12 parameters; only
\texttt{performative} is mandatory):
\begin{center}\small
\begin{tabular}{@{}p{0.38\textwidth} p{0.54\textwidth}@{}}
\toprule
\textbf{Field} & \textbf{Meaning} \\
\midrule
\texttt{performative} & the type of the communicative act (illocutionary force) --- \emph{mandatory} \\
\texttt{sender}, \texttt{receiver}, \texttt{reply-to} & agent identifiers (AID) \\
\texttt{content} & the content of the message itself \\
\texttt{language} & the language of the content (SL, KIF, \dots) \\
\texttt{encoding}, \texttt{ontology} & the encoding and ontology for interpreting the content \\
\texttt{protocol} & the interaction protocol the message is part of \\
\texttt{conversation-id}, \texttt{reply-with}, \texttt{in-reply-to} & matching messages into a conversation \\
\texttt{reply-by} & the deadline by which a reply is expected \\
\bottomrule
\end{tabular}
\end{center}
\begin{quote}\small\ttfamily
(inform\\
\hspace*{1em}:sender agent1\\
\hspace*{1em}:receiver agent2\\
\hspace*{1em}:content (price good2 150)\\
\hspace*{1em}:language sl\\
\hspace*{1em}:ontology hpl-auction)
\end{quote}
\end{defbox}

\textbf{FIPA-ACL performatives} (22 in total, the most important ones):
\begin{center}\small
\begin{tabular}{p{0.24\textwidth} p{0.70\textwidth}}
\toprule
\textbf{Group} & \textbf{Performatives} \\
\midrule
passing information & \texttt{inform}, \texttt{inform-if}, \texttt{inform-ref},
  \texttt{confirm}, \texttt{disconfirm} \\
requesting information & \texttt{query-if}, \texttt{query-ref}, \texttt{subscribe} \\
requesting action & \texttt{request}, \texttt{request-when}, \texttt{request-whenever},
  \texttt{cancel} \\
negotiation & \texttt{cfp} (\emph{call for proposals}), \texttt{propose},
  \texttt{accept-proposal}, \texttt{reject-proposal} \\
performing actions & \texttt{agree}, \texttt{refuse}, \texttt{failure} \\
errors, routing & \texttt{not-understood}, \texttt{proxy}, \texttt{propagate} \\
\bottomrule
\end{tabular}
\end{center}

\begin{defbox}[The semantics of FIPA-ACL: FP and RE]
The semantics is defined in the language SL (\emph{Semantic Language}) --- a modal logic
with the operators $B_i\varphi$ (agent $i$ believes $\varphi$), $U_i\varphi$ (it is uncertain),
$I_i\varphi$ (it has the intention), $\mathrm{Feasible}$, $\mathrm{Done}$. For every performative the
following are given:
\begin{itemize}
\item \textbf{FP} (the \emph{feasibility precondition}) --- what must hold for the sender to be able
  to perform the act;
\item \textbf{RE} (the \emph{rational effect}) --- the effect for the sake of which it performs the
  act. The recipient is \emph{not obliged} to bring about the RE (autonomy!) --- the RE is only the
  sender's reason.
\end{itemize}
For example, for $\langle i, \mathrm{inform}(j,\varphi)\rangle$:
\[
\text{FP: } B_i\varphi \wedge \neg B_i\big(\mathit{Bif}_j\varphi \vee \mathit{Uif}_j\varphi\big),
\qquad
\text{RE: } B_j\varphi .
\]
That is: \uv{I believe $\varphi$ and I do not believe that $j$ already knows about $\varphi$ (that it
would know whether it holds or not --- $\mathit{Bif}$ --- nor that it would hold an uncertain opinion
about it --- $\mathit{Uif}$); I want $j$ to come to believe $\varphi$.}
For $\langle i, \mathrm{request}(j,\alpha)\rangle$:
$\text{FP: } \mathit{FP}(\alpha)[i \backslash j] \wedge B_i \mathrm{Agent}(j,\alpha)
\wedge \neg B_i I_j \mathrm{Done}(\alpha)$, $\text{RE: } \mathrm{Done}(\alpha)$.

\textbf{A practical problem:} this semantics is unverifiable --- it is impossible to check from the
outside that an agent really believes what it claims (the \emph{sincerity assumption}). This is a
well-known weakness of the mentalistic semantics of ACL; the alternative is a \emph{social} semantics
based on public commitments.
\end{defbox}

\subsection{Interaction protocols}

A single message is not enough; agents hold \textbf{conversations} with a prescribed structure.
An interaction protocol is a pattern of admissible sequences of messages;
FIPA standardises about a dozen of them (the FIPA-IP specification, AUML notation).

\begin{center}\small
\begin{tabular}{p{0.26\textwidth} p{0.68\textwidth}}
\toprule
\textbf{Protocol} & \textbf{Course} \\
\midrule
\textbf{FIPA-Request}
  & $I \to P$: \texttt{request}. $P \to I$: \texttt{refuse} (end), or \texttt{agree};
    after execution \texttt{inform-done} / \texttt{inform-result}, or possibly \texttt{failure}.
    The reply may also be \texttt{not-understood}. \\
\textbf{FIPA-Query}
  & $I \to P$: \texttt{query-if} / \texttt{query-ref}. $P$: \texttt{refuse} or \texttt{agree},
    then \texttt{inform} with the result or \texttt{failure}. \\
\textbf{FIPA-Contract-Net}
  & see below: \texttt{cfp} $\to$ $n\times$(\texttt{propose} / \texttt{refuse}) $\to$
    \texttt{accept-proposal} to the winner + \texttt{reject-proposal} to the others $\to$
    \texttt{inform-done} / \texttt{failure}. \\
\textbf{FIPA-Iterated-Contract-Net}
  & like CNET, but after receiving the bids the initiator may announce \emph{another round} of
    \texttt{cfp} (negotiating better terms). \\
\textbf{FIPA-Subscribe}
  & $I$: \texttt{subscribe}; $P$ sends \texttt{inform} every time the monitored value changes,
    until a \texttt{cancel} arrives. \\
\textbf{FIPA-Brokering / Recruiting}
  & $I$ asks a \emph{broker}, which finds suitable providers, delegates the request to them, and
    returns the results (or possibly lets them answer directly). \\
\textbf{Auction protocols}
  & FIPA-English-Auction, FIPA-Dutch-Auction (see chapter~\ref{sec:aukce}). \\
\bottomrule
\end{tabular}
\end{center}

\textbf{A note on design.} A protocol determines which messages are legal in a given state of the
conversation --- thereby it introduces \emph{predictability} into otherwise autonomous behaviour and
makes it possible to implement an agent as a finite automaton over the states of the conversation
(which is exactly what JADE does with classes such as
\texttt{AchieveREInitiator}, \texttt{ContractNetResponder}, and the like).

\subsection{Exam summary}

\begin{itemize}
\item Agent communication is an \emph{action}, not a method call --- the recipient is not obliged to
  comply.
\item Austin: speech acts; locution / illocution / perlocution; Searle: felicity conditions
  and the five categories (assertives, directives, commissives, expressives, declarations); a
  performative verb + a propositional content.
\item Cohen \& Perrault: STRIPS-style semantics of speech acts (pre / effect) with modal
  operators; be able to write down Request(S,H,A) and Inform(S,H,$\varphi$) --- the effect of Inform is
  \emph{$H$ believes that $S$ believes $\varphi$}.
\item KQML (the envelope) + KIF (the content); parameters and performatives; the criticism (no commit,
  unclear semantics).
\item FIPA-ACL: the structure of a message (performative, sender/receiver, content, language, ontology,
  protocol, conversation-id, reply-by), 22 performatives, semantics via FP and RE in the language SL;
  the problem of unverifiability.
\item Protocols: FIPA-Request, FIPA-Query, Contract Net (+ the iterated version), Subscribe,
  Brokering, auction protocols.
\end{itemize}

\section{Distributed problem solving and cooperation}
\label{sec:cdps}

\subsection{Coherence and coordination}

Agents have different goals, they are autonomous, they operate over time, and they decide at run
time; they therefore have to be capable of \emph{dynamic coordination}. In doing so they share two
things: \textbf{tasks} (task sharing) and \textbf{information} (result sharing).

\begin{defbox}[Coherence and coordination]
\textbf{Coherence} --- how well the system functions \emph{as a whole}
(the quality of the solution, the efficiency of resource use, degradation under failures).

\textbf{Coordination} --- how well the agents manage to \emph{minimise the overhead}
associated with synchronisation: avoiding needlessly duplicated work, conflicts over resources, and
unnecessary communication.

A system may be coherent even with poor coordination (the agents do the work three times over, but the
result is correct) --- at the price of waste.
\end{defbox}

\begin{defbox}[CDPS --- \emph{Cooperative Distributed Problem Solving}]
Lesser et al., 1980s. The cooperation of individual agents in solving a problem that exceeds their
individual capabilities (information, resources, computational capacity).

The key assumption: the agents \emph{implicitly share a common goal} and have no conflicts ---
\textbf{benevolence}. The measure of success is the \emph{performance of the system as a whole}; an
agent helps the whole even when this is disadvantageous for itself. Benevolence \emph{enormously
simplifies} the design of the system.
\end{defbox}

\begin{center}\small
\begin{tabular}{p{0.14\textwidth} p{0.38\textwidth} p{0.40\textwidth}}
\toprule
& \textbf{Characteristics} & \textbf{Difference} \\
\midrule
\textbf{PPS} \newline (\emph{parallel problem solving})
  & Bond, Gasser, 1980s; the emphasis is on parallel solving, with homogeneous and simple processors
  & the agents are not autonomous, this is essentially a parallel algorithm \\
\textbf{CDPS}
  & heterogeneous agents with sophisticated capabilities, a shared goal, benevolence
  & conflicts are not resolved, because they do not arise \\
\textbf{MAS}
  & a society of agents with \emph{their own} goals, which they \emph{do not share}
  & it is necessary to address: why and how to cooperate, how to identify and resolve conflicts, how
    to negotiate and bargain \\
\bottomrule
\end{tabular}
\end{center}

\subsection{Task sharing and result sharing}

\textbf{The general CDPS procedure} has three phases:
\begin{enumerate}
\item \textbf{Problem decomposition} --- hierarchical, recursive. The questions are \emph{how} to
  decompose and \emph{who} performs the decomposition. An extreme variant is the ACTORS model: a new
  agent is created for every subproblem, down to the level of individual instructions.
\item \textbf{Solving the subproblems} --- during this the agents typically share information and may
  need to synchronise their actions.
\item \textbf{Solution synthesis} --- again hierarchical, the composition of partial results.
\end{enumerate}

\textbf{Task sharing} requires an \emph{agreement} among the agents (who does what) --- the typical
mechanism is the Contract Net. \textbf{Result sharing} may be \emph{proactive} (an agent sends a
result it believes will be useful) or \emph{reactive} (it sends it upon request).

\begin{defbox}[The benefits of result sharing]
\begin{itemize}
\item \textbf{Confidence} --- independent solutions of the same problem can be compared;
  agreement increases confidence in the result.
\item \textbf{Completeness} --- the agents share their \emph{local} views and thereby a more global
  picture of the problem arises.
\item \textbf{Precision} --- sharing may make the overall solution more precise.
\item \textbf{Timeliness} --- the obvious advantage of distribution: a reduction in time.
\end{itemize}
\end{defbox}

\subsection{The Contract Net}

\begin{defbox}[Contract Net Protocol (CNET)]
Smith and Davis, 1977--1980. The metaphor of \textbf{task sharing by means of contracts} following the
model of a market: the agent in need is the \emph{manager}, the others are \emph{contractors}. Four
phases:
\begin{enumerate}
\item \textbf{Recognition} --- an agent finds that it has a problem it cannot handle on its own (it
  lacks the capability, the resources, or the information, or the solution would be worse).
\item \textbf{Announcement} --- it sends out (typically by broadcast) an announcement of the task:
  a \emph{task description} (which may even be executable), \emph{constraints} (deadline, price,
  quality), and \emph{meta-information} (who may apply).
\item \textbf{Bidding} --- the recipients decide whether they have the interest and the capabilities,
  and send a bid (a \emph{tender}) with a price / time / quality.
\item \textbf{Awarding and expediting} --- the manager selects the winner and concludes a contract with
  it; the contractor performs the task and returns the result.
\end{enumerate}
\textbf{Properties:} simple, decentralised, robust against failures; a contractor may itself become
the manager of a subtask, which leads to \emph{hierarchical cascades of subcontracts}.
It distributes the load dynamically. It is the most frequently implemented coordination mechanism of
all (in FIPA it is standardised as FIPA-Contract-Net; in JADE it is available as a ready-made
behaviour).

\textbf{Limits:} broadcasting is expensive with many agents; it assumes benevolence (a contractor does
not lie about its capabilities) --- otherwise it is necessary to use auction mechanisms that are robust
against manipulation (chapter~\ref{sec:aukce}).
\end{defbox}

\subsection{Blackboard systems}

\begin{defbox}[Blackboard system (BBS)]
The very first scheme for cooperative problem solving (HEARSAY-II, speech recognition, 1970s).
Results are shared through a common data structure --- the \textbf{blackboard}.

The metaphor: several experts sit around a blackboard and can both write on it and read from it. Tasks
appear on the blackboard dynamically; as soon as one of the experts sees a task it knows how to solve,
it writes a partial solution on the blackboard. This continues until the solution of the whole task
appears on the blackboard.

\textbf{The roles:}
\begin{itemize}
\item The \textbf{blackboard} --- shared memory; typically structured into several \emph{levels of
  abstraction} (a hierarchy of goals and actions). The formalism for representing information is
  essential.
\item The \textbf{experts} (\emph{knowledge sources}) --- agents solving a particular type of task;
  they react to the goals on the blackboard, take their goal from the blackboard, and once selected
  they perform an action. They have a list of capabilities and an efficiency.
\item The \textbf{arbiter} (the \emph{control component}) --- it selects which expert may approach the
  blackboard. It may be purely reactive, or it may reason about plans and maximise expected utility.
  It is responsible for solving the problem at a higher level.
\end{itemize}
\textbf{Mutual exclusion over the blackboard is necessary} --- the bottleneck of the system.
\end{defbox}

\textbf{Pros and cons of BBS:}
\begin{itemize}
\item[$+$] A simple mechanism for coordination, cooperation, and the sharing of both tasks and
  results.
\item[$+$] The experts \emph{need not know about each other} and yet they cooperate (\emph{loose
  coupling}); adding an expert requires no change to the others.
\item[$+$] Messages on the blackboard can be rewritten --- delegating tasks, creating subtasks,
  changing experts.
\item[$-$] All agents must have access to the blackboard and share the same architecture; it tends to
  be \uv{crowded} around the blackboard (distributed hash tables partly address this).
\end{itemize}

\medskip
\noindent\textbf{Example: BBWar.} The blackboard is a hash table mapping required capabilities to
tasks; \uv{open missions} are published on the blackboard. The experts form a hierarchy: a
\emph{commander} agent looks for tasks of the type \uv{conquer a city} (\texttt{ATTACK-CITY}) and
converts them into several tasks \uv{attack a position} (\texttt{ATTACK-LOCATION}); soldiers of
various types then pick from the blackboard the missions they have the capabilities for.

\medskip
\noindent\textbf{Example: FELINE} (Wooldridge, Jennings, 1990) --- a distributed expert system from
the time before KQML and ACL. Knowledge sharing and the distribution of subtasks; each agent is a
rule-based system that maintains:
\emph{Skills} --- \uv{I can prove/refute the following\dots} and \emph{Interests} ---
\uv{I am interested in whether the following holds\dots}. Communication is a triple
(sender, recipient, content), where the content is a hypothesis + a speech act (request, response,
inform).

\subsection{Distributed constraints (DisCSP and DCOP)}

Both CNET and the blackboard are \emph{protocols}. Alongside them stands the formal branch of
distributed problem solving, in which the task is given as a system of constraints spread among the
agents.

\begin{defbox}[DisCSP and DCOP]
A \textbf{distributed CSP} (distributed constraint satisfaction problem): the variables
$x_1,\dots,x_n$ with domains $D_i$ are divided among the agents (typically one variable per agent)
and the constraints link the variables of \emph{different} agents. An assignment satisfying all the
constraints is sought; no agent sees the whole problem, and each communicates only with its neighbours
in the constraint graph.

\textbf{DCOP} (distributed constraint optimization) is the optimisation variant: the constraints are
\emph{soft}, i.e.\ given by cost functions $f_{ij}(x_i,x_j)$, and $\sum_{ij} f_{ij}$ is maximised ---
that is, social welfare (chapter~\ref{sec:hry}).

\textbf{Algorithms:} \emph{ABT} (\emph{asynchronous backtracking}, Yokoo) --- the agents are ordered
and send each other \texttt{ok?} messages (their assignment downwards) and \texttt{nogood} messages
(back upwards upon a conflict), where the learned \emph{nogoods} prevent cycling;
\emph{ADOPT} (asynchronous, with bounds, optimal), \emph{DPOP} (dynamic programming over a
pseudo-tree --- linearly many messages, but exponentially large ones), and \emph{Max-Sum}
(message passing over a factor graph, approximate but very scalable).

\textbf{Typical applications:} meeting scheduling, assigning tracking tasks to sensors,
frequency allocation, coordination in MAPF. Unlike in CDPS, benevolence is \emph{not} required here:
DCOP can be combined with the mechanisms of chapter~\ref{sec:aukce}.
\end{defbox}

\subsection{Joint intentions and social laws}

What exactly does it mean for a team of agents to do something \emph{jointly}? The answer is the
multiagent generalisation of commitment from chapter~\ref{ssec:commitment}:

\begin{defbox}[The theory of joint intentions (Cohen--Levesque)]
A team $G$ has a \textbf{joint persistent goal} $\varphi$ if:
(1)~all members believe that $\varphi$ does not hold yet and want it to hold;
(2)~this is mutually known among them;
(3)~each member holds the goal until it learns that $\varphi$ is \emph{achieved},
\emph{unachievable}, or \emph{irrelevant} --- and whoever finds this out takes on the
\textbf{commitment to let the others know} (the state of the team becomes that agent's individual goal
to announce). It is precisely the communicative obligation in point (3) that distinguishes team
activity from a chance coincidence of individual intentions (Searle's example: runners who all dash
for the same bandstand when it starts raining are not a team). Practical models of team cooperation
build on this theory: Jennings's model of \emph{joint commitments}
with \emph{conventions} --- explicit rules stating when an agent may abandon a commitment and whom it
must inform --- and the \textbf{STEAM} system (Tambe); an alternative formalisation is
\emph{SharedPlans} (Grosz--Kraus).
\end{defbox}

\textbf{Social laws and conventions} (Shoham--Tennenholtz): coordination need not take place at run
time through communication --- it can be ensured \emph{offline} by rules that constrain the actions of
all agents in advance (the model example: \uv{drive on the right}). A social law is
\emph{useful} if it leaves every agent the possibility of achieving its goals;
the synthesis of a useful social law is NP-hard in general. Conventions may also
\emph{emerge} from below through repeated interaction (cf.\ Axelrod's tournaments,
chapter~\ref{sec:hry}). Wooldridge summarises four basic coordination techniques:
\textbf{organisational structuring} (roles and hierarchies, cf.\ Gaia,
chapter~\ref{sec:metodologie}), \textbf{contracting} (CNET), \textbf{multiagent
planning}, and \textbf{negotiation} (chapter~\ref{sec:aukce}).

\subsection{Agent interaction in the environment}

Agents can help and hinder one another even \emph{without communication}, simply by affecting the
environment:
\begin{itemize}
\item Robots can carry a brick only if they push from two sides simultaneously (a positive
  interaction, requiring synchronisation).
\item Robots crowd at a door and cannot open it (a negative interaction, requiring coordination).
\end{itemize}
Agents can form communities, hierarchies, and hostilities. \textbf{It is necessary to know the types
of interaction} in a specific MAS --- otherwise effective control mechanisms cannot be designed.
The simplest model case, which the following chapter deals with, is the interaction of two
rational (selfish) agents in an environment resembling a game.

\subsection{Exam summary}

\begin{itemize}
\item Coherence (the performance of the whole) vs.\ coordination (minimising synchronisation
  overhead).
\item CDPS: benevolence + a shared goal; distinguish it from PPS (homogeneous processors, a parallel
  algorithm) and from MAS (own goals, conflicts, negotiation).
\item The three phases of CDPS: decomposition --- solving the subproblems --- synthesis. Task sharing
  (agreement) vs.\ result sharing (proactive/reactive) and its four benefits: confidence,
  completeness, precision, timeliness.
\item CNET: recognition --- announcement --- bidding --- awarding/expediting; the manager
  and the contractor, hierarchical subcontracts; standardised as FIPA-Contract-Net.
\item BBS: the blackboard (a hierarchy of abstraction, mutex), the experts (skills), the arbiter
  (selecting an expert); loose coupling between the experts; access to the blackboard as a bottleneck.
\item DisCSP / DCOP: variables and constraints divided among the agents, nobody sees the whole
  problem; ABT (\texttt{ok?} / \texttt{nogood}), ADOPT, DPOP, Max-Sum; DCOP maximises social
  welfare.
\item A joint persistent goal: the team holds the goal until it is achieved /
  unachievable / irrelevant, and whoever finds this out must announce it to the others; social
  laws = offline coordination by constraining actions (\uv{drive on the right}).
\end{itemize}

\section{Agent interaction and game theory}
\label{sec:hry}

\subsection{The selfish agent}

Why should an agent be honest about its capabilities? Why should it complete an assigned task?
If the system is homogeneous (designed entirely by us), benevolence is a good strategy. Usually,
however, the system contains agents with different interests --- a conflict between the common goal
and the goals of individuals (e.g.\ air traffic control: everybody wants safety, but every airline
wants to land first). Sometimes it is not even clear what the common interest is. In that case it is
better to consider \textbf{self-interested} agents.

\begin{defbox}[Self-interested agent]
The set of possible outcomes $O = \{o_1, o_2, \dots\}$ is \emph{common} to all agents,
but the preferences are not: each agent $i$ has its own \textbf{utility function}
$u_i : O \to \mathbb{R}$, which orders the outcomes.

\textbf{Strategic thinking:} maximise your utility under the assumption that everybody else acts
rationally as well. This does \emph{not} maximise the common utility, but it is a \emph{robust}
strategy.

\emph{A note:} money is not a good utility for a human being --- the utility function of money is
non-linear and differs for the rich and the poor.
\end{defbox}

\subsection{Normal-form games}

\begin{defbox}[Two-player game in normal form]
Agents $i$ and $j$ with action sets $A_i$, $A_j$. The outcome is determined by both:
$e : A_i \times A_j \to O$. An agent chooses a \textbf{strategy} $s_i$:
\begin{itemize}
\item \textbf{pure} --- a deterministic choice of a single action;
\item \textbf{mixed} --- a probability distribution over pure strategies (a random choice).
\end{itemize}
A game is written down as a \textbf{payoff matrix}: rows = the strategies of player $i$,
columns = the strategies of player $j$, and a cell contains the pair $u_i / u_j$.

In general: a game in normal form is a triple $\langle N, (A_i)_{i \in N}, (u_i)_{i \in N}\rangle$,
where $N$ is the set of players.
\end{defbox}

\begin{defbox}[Dominance]
A strategy $s_i$ \textbf{strictly dominates} a strategy $s_i'$ if it gives agent $i$ a
\emph{strictly} better outcome against \emph{all} strategies of the opponent;
it \textbf{weakly dominates} it if it gives an outcome at least as good against all of them and a
strictly better one against at least one (this is also how the lecture course defines dominance). A
strategy is \textbf{dominant} if it dominates all the other strategies of agent $i$.

The difference is not academic: truthfulness in the Vickrey auction is only \emph{weakly} dominant
(chapter~\ref{sec:aukce}).

A rational agent plays a dominant strategy if one exists. \emph{Iterated elimination of strictly
dominated strategies} (IESDS) is the basic solution technique: nobody will ever play a dominated
strategy, so it can be struck out of the matrix, which may render further strategies dominated.
When eliminating \emph{strictly} dominated strategies the order does not matter and the result is
unique; with \emph{weakly} dominated strategies the order \textbf{does} matter and some equilibria may
be lost.
\end{defbox}

\begin{defbox}[Nash equilibrium]
A pair of strategies $(s_1^*, s_2^*)$ is a \textbf{Nash equilibrium} if:
\begin{itemize}
\item when agent $i$ plays strategy $s_1^*$, it is best for agent $j$ to play $s_2^*$, and at the same
  time
\item when agent $j$ plays strategy $s_2^*$, it is best for agent $i$ to play $s_1^*$.
\end{itemize}
That is, $s_1^*$ and $s_2^*$ are \emph{mutual best responses}.
In general, for $n$ players: a profile $s^* = (s_1^*,\dots,s_n^*)$ is a NE if for every $i$ and every
strategy $s_i$ of that player
\[
u_i(s_i^*, s_{-i}^*) \ge u_i(s_i, s_{-i}^*),
\]
where $s_{-i}$ denotes the strategies of all the others. \textbf{No player has a unilateral reason to
deviate.}

A naive search for equilibria for $n$ agents and $m$ strategies requires going through $m^n$ profiles.
This definition concerns \emph{pure} strategies --- not every game has a NE in pure strategies
(rock--paper--scissors) and some games have several.
\end{defbox}

\begin{theorem}[Nash's theorem, 1950]
Every game with a finite number of players and a finite number of strategies has at least one Nash
equilibrium \emph{in mixed strategies}.
\end{theorem}

\textbf{Computational complexity.} Finding a NE is a \emph{total} search problem (a solution always
exists, the only question is how to find it). Since 2006 we have known that it is
\textbf{PPAD-complete} --- very roughly speaking \uv{somewhat less hopeless than NP-completeness}
(NP-completeness does not even make sense for a problem that always has a solution). For three and more
players this was proved by Daskalakis, Goldberg, and Papadimitriou; shortly afterwards Chen and Deng
supplied the remaining and hardest case of \emph{two} players.

\begin{defbox}[Pareto efficiency and social welfare]
A strategy profile is \textbf{Pareto optimal} (Pareto efficient) if there is no other profile that
would improve the outcome of one agent without worsening the outcome of some other one.
A non-Pareto-optimal solution can therefore be improved \uv{for free}.

The \textbf{social welfare} of a profile is the sum of the utilities of all agents:
$\mathit{sw}(o) = \sum_{i \in N} u_i(o)$. Maximising social welfare makes sense
in cooperative scenarios --- when the agents are on the same team, have a single owner, or solve a
single task; the more homogeneous the system, the better. The maximum of social welfare is always
Pareto optimal, but not conversely.
\end{defbox}

\subsection{Classical games}

\begin{defbox}[The Prisoner's dilemma]
\begin{center}
\begin{tabular}{c|cc}
$i \backslash j$ & $C$ (cooperate) & $D$ (defect) \\
\hline
$C$ & $3/3$ & $0/4$ \\
$D$ & $4/0$ & $1/1$ \\
\end{tabular}
\end{center}
$C$ = \emph{cooperate}, cooperate with one's accomplice; $D$ = \emph{defect}, inform on him and get a
lower sentence. The canonical notation for the payoffs: $R$ (\emph{reward} for mutual cooperation),
$P$ (\emph{punishment} for mutual defection), $T$ (\emph{temptation} --- I defect against a
cooperator), $S$ (\emph{sucker's payoff} --- I am betrayed). The conditions of the prisoner's dilemma:
\[
T > R > P > S \qquad\text{and often in addition}\qquad 2R > T + S .
\]
\begin{itemize}
\item $R > P$: mutual cooperation is better than mutual defection.
\item $T > R$ and $P > S$: \emph{defection is a dominant strategy for both agents}.
\item $(D,D)$ is the \textbf{only Nash equilibrium}.
\item $(D,D)$ is at the same time the \textbf{only non-Pareto-optimal} outcome of the game.
\item $(C,C)$ maximises social welfare ($3+3=6$).
\end{itemize}
\textbf{The rationality of each individual therefore leads to the collectively worst outcome.}
\end{defbox}

\textbf{Consequences of the prisoner's dilemma:}
\begin{itemize}
\item \textbf{The tragedy of the commons}: a shared resource
  (pasture, fisheries, the atmosphere) is depleted because every individual maximises their own
  benefit.
\item The question of \emph{what being rational actually means} --- and whether people are rational
  (experimentally, people cooperate more often than the theory predicts).
\item \textbf{The shadow of the future}: in the \emph{iterated} prisoner's dilemma the situation
  changes.
\end{itemize}

\begin{defbox}[The iterated prisoner's dilemma and Axelrod's tournaments]
If the PD is played \emph{a finite number of times and both players know it}, backward induction yields
$D$ in every round (in the last round there is no reason to cooperate, hence none in the
next-to-last either, \dots). If it is played \emph{infinitely} (or with an unknown end, i.e.\ with a
probability of continuing), cooperation becomes rational --- the \emph{folk theorem}: every payoff
profile better than the minimax can be supported as an equilibrium of the repeated game.

In 1980 Robert Axelrod organised tournaments of strategies for the iterated PD. The winner (twice) was
the simplest strategy submitted, \textbf{Tit For Tat} (TFT):
\emph{cooperate in the first round, then always repeat the opponent's last move}. From the results
Axelrod derived four properties of successful strategies:
\begin{itemize}
\item being \textbf{nice} --- do not defect first;
\item being \textbf{retaliating} --- respond to defection immediately;
\item being \textbf{forgiving} --- forgive an opponent who returns to cooperation;
\item being \textbf{non-envious} --- do not strive to beat a \emph{particular} opponent; TFT never
  scored more points than its counterpart, and yet it won the tournament.
\end{itemize}
As a fifth property Axelrod lists \textbf{clarity}: a strategy must be legible to the opponent,
otherwise no cooperation can be established with it.
The weakness of TFT: under noise (a move executed incorrectly) two TFT players get stuck in mutual
retaliation --- hence the variants \emph{Generous TFT} (which occasionally forgives) and
\emph{Tit for Two Tats}.
\end{defbox}

\begin{defbox}[Coordination game]
\begin{center}
\begin{tabular}{c|cc}
$i \backslash j$ & Ballet & Wrestling \\
\hline
Ballet & $1/2$ & $0/0$ \\
Wrestling & $0/0$ & $2/1$ \\
\end{tabular}
\end{center}
Also known as the \emph{Battle of the Sexes} (BoS) or \emph{Bach or Stravinsky}. The payoffs favour
\emph{agreement} of strategies, but the players differ in which agreement they care about more.
\begin{itemize}
\item Nash equilibria in pure strategies: (Ballet, Ballet) and (Wrestling, Wrestling).
\item Social welfare and Pareto optimality: again both of these pairs.
\item \textbf{The mixed NE}: each player plays their \emph{more preferred} option
  with probability $2/3$ and the less preferred one with $1/3$. The expected utility in the mixed
  equilibrium is only $2/3$ --- \emph{less} than in either pure equilibrium.
\end{itemize}

\emph{Computing the mixed NE (the general recipe for $2\times2$ games):} player $j$ chooses the
probability $q$ of playing Ballet so that player $i$ is \emph{indifferent} between their actions:
\[
u_i(\text{Ballet}) = 1\cdot q + 0\cdot(1-q), \qquad u_i(\text{Wrestling}) = 0\cdot q + 2\cdot(1-q).
\]
The equality $q = 2 - 2q$ gives $q = 2/3$, i.e.\ player $j$ plays Ballet (their preferred option)
with probability $2/3$. Symmetrically, player $i$ plays Wrestling with probability $2/3$.
The expected utility of player $i$ is then $1 \cdot \frac{2}{3} = \frac{2}{3}$.
\textbf{The key rule: in a mixed NE each player chooses their probabilities so as to make the
\emph{opponent} indifferent.}
\end{defbox}

\begin{defbox}[Anti-coordination game]
\begin{center}
\begin{tabular}{c|cc}
$i \backslash j$ & Chicken (swerve) & Dare (hold course) \\
\hline
Chicken & $5/5$ & $1/6$ \\
Dare & $6/1$ & $0/0$ \\
\end{tabular}
\end{center}
Also known as \emph{Chicken}, \emph{Hawk--Dove}, or \emph{Snowdrift}. The payoffs favour
the choice of \emph{different} strategies. In coordination games sharing a resource benefits everyone;
in anti-coordination games sharing is costly (a collision of cars).
\begin{itemize}
\item Nash equilibria in pure strategies: (Chicken, Dare) and (Dare, Chicken).
\item The maximum of social welfare: (Chicken, Chicken) --- $5+5=10$, which is \emph{not} a NE
  (both players benefit from deviating to Dare: $6 > 5$).
\item \textbf{The symmetric mixed NE}: if the opponent plays Chicken with probability $p$, then
  $u_i(\text{Chicken}) = 5p + 1(1-p) = 1+4p$ and $u_i(\text{Dare}) = 6p$; from the indifference
  $1+4p = 6p$ it follows that $p = 1/2$. Both therefore play both actions with probability $1/2$
  and the expected utility is $3$ --- less than the $5$ under (Chicken, Chicken), but nobody can
  improve unilaterally.
\end{itemize}
\end{defbox}

\medskip
\noindent\textbf{An addendum: a game with no pure NE.} \emph{Matching pennies} (and likewise
rock--paper--scissors) is a zero-sum game:
\begin{center}
\begin{tabular}{c|cc}
$i \backslash j$ & Heads & Tails \\
\hline
Heads & $1/-1$ & $-1/1$ \\
Tails & $-1/1$ & $1/-1$ \\
\end{tabular}
\end{center}
No pure NE exists. The mixed NE: both play $(1/2, 1/2)$, with expected utility $0$.
The motif is analogous to the Condorcet paradox from the following chapter: individually rational
choices have no \uv{stable} deterministic outcome.

\subsection{Cooperative (coalitional) games}

Everything so far has concerned \textbf{non-cooperative} games: the basic unit is the \emph{individual}
and their strategies, and agreements are not enforceable. \textbf{Cooperative} game theory changes the
unit to the \emph{group} and asks: who joins forces with whom, and how do they divide the joint gain?
For MAS this is a formal model of team formation and reward division.

\begin{defbox}[Coalitional game with transferable utility]
A pair $\langle N, v\rangle$, where $N$ is a set of agents and the \textbf{characteristic function}
$v : 2^N \to \mathbb{R}$ assigns to every \emph{coalition} $C \subseteq N$ the value its members can
secure by their own efforts ($v(\emptyset)=0$). The utility is \emph{transferable}: the value $v(C)$
can be divided arbitrarily among the members (\uv{in money}).

A game is \textbf{superadditive} if $v(C \cup D) \ge v(C) + v(D)$ for disjoint $C,D$ ---
in that case it always pays off to form the \emph{grand coalition} $N$, and only the question of the
division remains.

A \textbf{solution} is a vector $(x_1,\dots,x_n)$ with $\sum_i x_i = v(N)$ (\emph{feasibility}).
\end{defbox}

\begin{defbox}[The core and the Shapley value]
\textbf{The core} --- the set of divisions that \emph{no} coalition wants to leave:
\[
\text{core} = \Big\{ x : \sum_{i \in N} x_i = v(N) \ \wedge\
\forall C \subseteq N:\ \sum_{i \in C} x_i \ge v(C) \Big\}.
\]
It is the analogue of the Nash equilibrium for groups (\emph{stability}). The core may, however, be
\textbf{empty} --- in which case no stable division exists --- and it may also be very large,
so it does not tell us which division is \emph{fair}.

\textbf{The Shapley value} (1953) --- the answer to the question of fairness. An agent receives its
\emph{average marginal contribution} over all the orders in which the coalition can be assembled:
\[
\varphi_i(v) = \sum_{C \subseteq N \setminus \{i\}}
\frac{|C|!\,(|N|-|C|-1)!}{|N|!}\,\big(v(C \cup \{i\}) - v(C)\big).
\]
It is the \emph{only} division satisfying simultaneously: efficiency ($\sum_i \varphi_i = v(N)$),
symmetry (interchangeable agents receive the same), the \emph{dummy} property (an agent with a zero
marginal contribution receives nothing), and additivity.

\emph{Example (gloves):} agents 1 and 2 each have a left glove, agent 3 a right one; a pair has value
1, a single glove 0. Thus $v(\{1,3\}) = v(\{2,3\}) = v(N) = 1$, and 0 otherwise. The core contains a
\emph{single} point $(0,0,1)$ --- agent 3 takes everything, because agents 1 and 2 are dependent on it.
The Shapley value is $(1/6,\,1/6,\,2/3)$ --- fairer, but unstable.

\textbf{The computational side} (why this interests MAS): the characteristic function has $2^n$ values,
so even its \emph{representation} is a problem (compact languages and games over graphs are used);
deciding the non-emptiness of the core and computing the Shapley value are both hard in general,
and the \emph{coalition structure generation} problem (finding the optimal partition into coalitions
when the game is not superadditive) is NP-hard.
\end{defbox}

\subsection{Exam summary}

\begin{itemize}
\item A selfish agent maximises \emph{its own} expected utility under the assumption that the others do
  the same; this does not maximise the common utility.
\item Normal-form games, pure vs.\ mixed strategies, dominance and dominant strategies,
  iterated elimination of dominated strategies.
\item \textbf{Nash equilibrium} = mutual best responses, nobody has a reason to deviate unilaterally;
  a naive search takes $m^n$; \textbf{Nash's theorem} (existence in mixed strategies);
  finding one is PPAD-complete.
\item \textbf{Pareto optimality} (one agent cannot be made better off without making another worse off)
  vs.\ \textbf{social welfare} (the sum of utilities); the maximum of sw is Pareto optimal, but not
  conversely.
\item \textbf{The prisoner's dilemma}: $T>R>P>S$; $D$ dominant; $(D,D)$ the only NE and the only
  non-Pareto-optimal outcome; $(C,C)$ maximises sw. The tragedy of the commons.
  The iterated PD, the shadow of the future, Axelrod and TFT (nice, retaliating, forgiving, clear).
\item \textbf{Coordination game} (BoS): two pure NE + a mixed one, $2/3$--$1/3$, with worse utility;
  \textbf{anti-coordination} (Chicken): two asymmetric NE, and sw is maximised by a non-equilibrium
  profile.
\item Be able to compute the mixed NE for $2\times2$: choose the probabilities so as to make the
  \emph{opponent} indifferent.
\item \textbf{Cooperative games}: the characteristic function $v : 2^N \to \mathbb{R}$,
  superadditivity and the grand coalition; \textbf{the core} (no coalition wants to leave, may be
  empty) vs.\ \textbf{the Shapley value} (the average marginal contribution; efficiency, symmetry,
  dummy, additivity). The gloves example: core $(0,0,1)$, Shapley $(1/6,1/6,2/3)$.
\end{itemize}

\section{Social choice and voting}
\label{sec:volby}

\subsection{Preference aggregation}

The previous chapter dealt with how agents behave when they have conflicting interests. Now we ask the
converse question: how to \emph{derive} a collective decision from individual preferences. This is the
subject of \textbf{social choice theory}.

\begin{defbox}[Social welfare vs.\ social choice]
A \textbf{social welfare function} aggregates individual preferences into a \emph{complete ordering}
of the candidates $>_{sw}$.

A \textbf{social choice function} selects only \emph{one} winner (or a set of winners).

\emph{Example:} voters $A: o_2 > o_1 > o_3$, $B: o_3 > o_2 > o_1$, $C: o_2 > o_3 > o_1$.
Social welfare: $o_2 > o_3 > o_1$. Social choice: $o_2$.
\end{defbox}

\begin{defbox}[Ballot and preference schedule]
A \textbf{ballot} --- how a voter expresses their preferences. Traditionally:
\uv{pick your favourite from the list}. A \textbf{preference ballot}
--- the voter ranks \emph{all} the options by preference.

The individual ballots are combined into a \textbf{preference schedule};
from it a \textbf{voting scheme} determines the winner.
\end{defbox}

\subsection{Voting systems}

\begin{itemize}
\item \textbf{Majority} --- the winner gets more than 50~\% of the votes; trivial for two
  candidates.
\item \textbf{Plurality} --- a candidate gets a point for every first place, and the one with the most
  points wins; in practice the ballot contains only the first choice, and in the event of a tie a
  \emph{run-off} between the two best follows.
\item \textbf{Plurality with elimination} (\emph{instant runoff}, IRV) --- repeatedly eliminate the
  candidate with the fewest first places and \emph{redistribute} their votes to the next choice of the
  voters concerned, until the winner has a majority. It makes use of the already recorded preferences
  without a second round and suppresses insincere voting (the Australian House of Representatives, the
  IOC).
\item \textbf{Sequential majority} --- the candidates play tournament rounds in pairs and the winner
  advances (a knockout bracket). \emph{The order of the matches affects the result}
  (agenda setting).
\item \textbf{The Condorcet method} --- simulate \emph{all} pairwise contests; expensive, and it need
  not produce a winner.
\item \textbf{The Borda count} --- every voter submits complete preferences and a
  candidate gets $N-1$ points for a first place, $N-2$ for a second, \dots, $0$ for the last, and the
  points are summed. This is \emph{consensual} rather than majoritarian voting (Slovenia, Nauru,
  Iceland, academic institutions).
\item \textbf{Slater's system} --- from the pairwise majority contests build a tournament graph and
  look for an \emph{acyclic} ordering minimising the number of reversed edges. \textbf{NP-hard}.
\item \textbf{Combinations with Borda} (all of which elect the Condorcet winner if one exists):
  \emph{Black} (the Condorcet method, otherwise the Borda winner); \emph{Baldwin} (repeatedly eliminate
  the candidate with the lowest Borda score and recompute); \emph{Nanson} (repeatedly eliminate all
  candidates below the average).
\end{itemize}

\medskip
\noindent\textbf{The plurality winner need not be a majority winner.} Example:
$o_1$ --- 40~\%, $o_2$ --- 30~\%, $o_3$ --- 30~\%. Candidate $o_1$ wins, but 60~\% of the voters did not
want them. This leads to \textbf{tactical (insincere) voting}:
if $o_2$ and $o_3$ are similar (from the same party), the voters of $o_2$ will vote for $o_3$
in order to defeat $o_1$. In practice plurality is used with the requirement of an absolute majority
and a possible second round between the two best (the Czech presidential and Senate elections).

\begin{defbox}[The Condorcet paradox]
There are situations in which \emph{whichever outcome we choose, a majority of voters will be
dissatisfied with it}.

\emph{Example:} $V_1: A > B > C$, \quad $V_2: B > C > A$, \quad $V_3: C > A > B$.
Pairwise contests: $A$ beats $B$ (2:1), $B$ beats $C$ (2:1), $C$ beats $A$ (2:1).
\textbf{The collective preference is cyclic}, even though all the individual preferences are
linear. There is no \emph{Condorcet winner}.

The situation is analogous to the non-existence of a Nash equilibrium in pure strategies in
rock--paper--scissors: individually rational (linear) preferences compose into a cycle
that has no stable deterministic top. This is not a formal reduction, but it is the same motif.

The Marquis de Condorcet (1743--94), a French philosopher and mathematician, active in the French
Revolution, died in prison. He proposed several \emph{criteria of fairness} and thereby founded the
mathematical theory of voting.
\end{defbox}

\subsection{A worked example: one preference schedule, four winners}
\label{ssec:volby-priklad}

A typical exam problem: given a preference schedule, compute the winner under the individual systems.
The following schedule is chosen so that \emph{four} systems give \emph{four different}
winners. Candidates $A, B, C, D$, and 17 voters in total:

\begin{center}\small
\begin{tabular}{lcccc}
\toprule
\textbf{Number of voters} & 5 & 6 & 4 & 2 \\
\midrule
1st place & $A$ & $C$ & $D$ & $B$ \\
2nd place & $D$ & $B$ & $B$ & $A$ \\
3rd place & $B$ & $D$ & $A$ & $C$ \\
4th place & $C$ & $A$ & $C$ & $D$ \\
\bottomrule
\end{tabular}
\end{center}

\noindent\textbf{1. Plurality.} First places: $C\,{=}\,6$, $A\,{=}\,5$,
$D\,{=}\,4$, $B\,{=}\,2$. \textbf{Winner $C$} (nobody has an absolute majority, i.e.\ 9 votes).

\medskip
\noindent\textbf{2. Plurality with elimination (IRV).}
\begin{itemize}
\item Round 1: nobody has 9 votes; we eliminate $B$ (2 first places). Its voters
  ($B\!>\!A\!>\!C\!>\!D$) have $A$ as their next choice: $A = 5+2 = 7$, $C = 6$, $D = 4$.
\item Round 2: still nobody has 9; we eliminate $D$ (4). Its voters ($D\!>\!B\!>\!A\!>\!C$)
  have, after the elimination of $B$, the next choice $A$: $A = 7+4 = 11$, $C = 6$.
\item $A$ has 11 out of 17, i.e.\ a majority --- \textbf{winner $A$}.
\end{itemize}

\medskip
\noindent\textbf{3. The Borda count} (4 candidates, hence $3/2/1/0$ points):
\begin{align*}
A &= 5\cdot3 + 6\cdot0 + 4\cdot1 + 2\cdot2 = 23, &
B &= 5\cdot1 + 6\cdot2 + 4\cdot2 + 2\cdot3 = 31, \\
C &= 5\cdot0 + 6\cdot3 + 4\cdot0 + 2\cdot1 = 20, &
D &= 5\cdot2 + 6\cdot1 + 4\cdot3 + 2\cdot0 = 28 .
\end{align*}
\textbf{Winner $B$} (a check on the sum: $23{+}31{+}20{+}28 = 102 = 17\cdot(3{+}2{+}1{+}0)$).

\medskip
\noindent\textbf{4. Pairwise contests (Condorcet).} It suffices to go through the contests of
candidate $D$:
\[
D : A = 10 : 7, \qquad D : B = 9 : 8, \qquad D : C = 9 : 8 .
\]
(E.g.\ $D$ vs.\ $B$: groups \uv{4} and \uv{5} vote for $D$, giving $9$, and groups
\uv{6} and \uv{2} vote for $B$, giving $8$.) Candidate $D$ beats everyone --- it is the
\textbf{Condorcet winner}, and it would therefore also win \emph{sequential majority} voting under any
bracket. The remaining contests give $B\!:\!A = 12\!:\!5$, $B\!:\!C = 11\!:\!6$, $A\!:\!C = 11\!:\!6$,
so the tournament graph is acyclic and Slater's ordering is $D > B > A > C$ (with no reversed
edges).

\medskip
\noindent\textbf{What to take away from the example.}
\begin{itemize}
\item The same preference schedule gives \textbf{four different winners}: plurality $C$, IRV $A$,
  Borda $B$, Condorcet/sequential $D$. \uv{Who won the election} is therefore not a property of the
  voters, but of the choice of \emph{rules}.
\item The plurality winner $C$ is here even the \textbf{Condorcet loser} --- it loses
  \emph{every} pairwise contest ($11\!:\!6$, $11\!:\!6$, $9\!:\!8$). This is the strongest possible
  illustration of why plurality does not satisfy the Condorcet winner condition.
\item Neither Borda nor IRV elected the Condorcet winner $D$ --- neither of them satisfies the
  Condorcet winner condition. This is precisely why the Black / Baldwin / Nanson combinations were
  devised.
\item Borda's ordering ($B > D > A > C$) differs from the ordering by pairwise majority
  ($D > B > A > C$) --- the difference between \emph{consensual} and \emph{majoritarian} voting.
\end{itemize}
The differing results for the same schedule are a concrete form of what Arrow's theorem says in
general.

\subsection{Properties of voting systems}

\begin{defbox}[Criteria]
\textbf{The Pareto property} (unanimity): if \emph{every} voter
prefers $X$ to $Y$, then $X >_{sw} Y$ should hold.
\emph{Satisfied by:} the majority system and the Borda count. \emph{Not satisfied by:} sequential
majority voting.

\textbf{The Condorcet winner condition}: a candidate who beats all the others in pairwise
comparison (the \emph{Condorcet winner}) should be the overall winner.
\emph{Satisfied by:} sequential majority voting and the Borda combinations (Black, Baldwin, Nanson).
\emph{Not satisfied by:} plurality, or the Borda count on its own.

\textbf{Independence of irrelevant alternatives} (IIA): whether $X >_{sw} Y$ should depend
\emph{only} on the relative ordering of $X$ and $Y$ in the voters' preferences. If the ordering of
other candidates $Z$, $W$ changes, the relation between $X$ and $Y$ must not be reversed.
\emph{Not satisfied by:} plurality, sequential majority voting, or the Borda count ---
practically no common system satisfies it.

\textbf{Unrestricted domain (universality)}: the system must
admit \emph{arbitrary} preferences of \emph{all} voters (and no other inputs), and for them
deterministically return a \emph{complete} ordering of the social choices.

\textbf{Dictatorship}: the result is determined by the preferences of a single voter.
\emph{Non-dictatorship:} there is no voter $i$ such that whenever $i$ strictly prefers $x$
to $y$, $x$ is socially chosen over $y$.
\end{defbox}

\begin{theorem}[Arrow's impossibility theorem, 1951]
For three or more candidates there exists no preferential voting system that would convert
individual ordered rankings into a complete and transitive social ordering and at the same time
satisfy all of the following criteria:
\begin{itemize}
\item unrestricted domain,
\item non-dictatorship,
\item Pareto efficiency (unanimity),
\item independence of irrelevant alternatives.
\end{itemize}
\end{theorem}

\textbf{A simpler (equivalent) formulation:} for an election with more than two candidates, the only
voting procedure satisfying the Pareto property and IIA is a \emph{dictatorship}, in which the social
outcome is simply given by one of the voters.

\textbf{Interpretation.} This is a \emph{negative} result showing the fundamental limits of collective
decision making. It certainly does \emph{not} follow from it that the only workable system is a
dictatorship --- what follows is that every real voting system has to sacrifice some of the criteria
(most often IIA) and that every system can be manipulated in some situation.

\medskip
\noindent\textbf{The Gibbard--Satterthwaite theorem} (an addendum): every social choice function
for three or more candidates that is \emph{surjective} and \emph{non-dictatorial} is
\textbf{manipulable} --- there is a situation in which it pays a voter to vote insincerely.
This is fundamental for MAS: an agent is a computer, and tactical voting poses no moral problem for it.
This is why MAS looks for mechanisms that are \emph{strategy-proof},
which is the topic of the following chapter on auctions.

\subsection{Exam summary}

\begin{itemize}
\item Social welfare (a complete ordering) vs.\ social choice (a single winner);
  the ballot and the preference schedule.
\item Systems: majority, plurality, plurality with elimination (IRV), sequential majority
  (the order decides), the Condorcet method, Borda ($N-1,\dots,0$), Slater (NP-hard),
  the Black / Baldwin / Nanson combinations.
\item \textbf{The Condorcet paradox}: $A>B>C$, $B>C>A$, $C>A>B$ gives a cyclic collective
  preference; a Condorcet winner need not exist.
\item Be able to work through a preference schedule (section~\ref{ssec:volby-priklad}): the same
  schedule gives plurality $C$, IRV $A$, Borda $B$, and Condorcet $D$ --- and the plurality winner may
  be the Condorcet \emph{loser}.
\item Criteria: Pareto (satisfied by majority and Borda, not by sequential), the Condorcet winner
  (satisfied by sequential and the Borda combinations), IIA (satisfied by almost nothing),
  universality, non-dictatorship.
\item \textbf{Arrow's theorem}: for $\ge 3$ candidates one cannot have universality,
  non-dictatorship, Pareto, and IIA at the same time. The simpler version: Pareto + IIA
  $\Rightarrow$ dictatorship.
\item Gibbard--Satterthwaite: every reasonable voting system can be manipulated by insincere
  voting.
\end{itemize}

\section{Resource allocation, auctions, and negotiation}
\label{sec:aukce}

\subsection{Auctions as a mechanism of resource allocation}

\begin{defbox}[Auction]
An auction is a \textbf{market institution} in which the messages from traders contain price
information --- an offer to buy at a given price (a \emph{bid}), or an offer to sell at a given price
(an \emph{ask}) --- and which gives priority to higher offers to buy and lower offers to sell.

In MAS an auction is above all a \textbf{mechanism for dividing scarce resources among agents}.
\emph{In the real world:} eBay, Sotheby's, mining licences, frequency bands for mobile operators.
\emph{In computer science:} allocating processor time, bandwidth, computational tasks,
advertising slots.

\textbf{The efficiency of an auction:} it allocates the resource to the agent that wants it most (that
has the highest valuation). The seller wants to \emph{maximise} the price, the buyer to
\emph{minimise} it; each buyer has its own utility function.
\end{defbox}

\begin{defbox}[Classification of auctions]
\textbf{By the valuation of the item:}
\begin{itemize}
\item \emph{private value} --- the value depends only on the buyer
  (a concert ticket);
\item \emph{common value} --- the value is the same for everybody, they just do not know it
  (mining rights) --- hence the \emph{winner's curse}:
  the winner is the one who erred the most in the upward direction;
\item \emph{correlated value} --- a combination.
\end{itemize}
\textbf{By protocol:} the winner pays the first / the second price; \emph{open cry}
vs.\ \emph{sealed bid}; a single round vs.\ several rounds.
\end{defbox}

\subsection{Mechanism design}
\label{ssec:mechanism}

An auction is not merely \uv{how bidding is done}; it is a special case of a more general problem. Game
theory asks \emph{how agents will behave under given rules}; \textbf{mechanism design} is
\uv{inverse game theory}: \emph{which rules to choose so that rational selfish agents will do of their
own accord what we want} --- typically tell the truth.

\begin{defbox}[A mechanism and its desirable properties]
A \textbf{mechanism} is a pair: an \emph{outcome selection rule}, which from the agents' declared
preferences $\hat v = (\hat v_1,\dots,\hat v_n)$ selects an outcome $x(\hat v)$,
and a \emph{payment rule} $p_i(\hat v)$. The utility of an agent is
$u_i = v_i(x(\hat v)) - p_i(\hat v)$ (\emph{quasi-linear} utility).
Desirable properties:
\begin{itemize}
\item \textbf{Incentive compatibility} (IC) --- it pays the agent to declare the truth,
  $\hat v_i = v_i$. The strongest variant: truthfulness is a \emph{dominant} strategy
  (\emph{strategy-proof}, DSIC) --- this is how Vickrey and VCG work.
\item \textbf{Individual rationality} --- nobody loses by participating
  ($u_i \ge 0$), otherwise the agents would not take part.
\item \textbf{Efficiency} --- the chosen outcome maximises social welfare $\sum_i v_i(x)$.
\item \textbf{Budget balance} --- the payments do not \uv{evaporate}
  ($\sum_i p_i = 0$, or a non-negative revenue for the organiser).
\item \textbf{Computational tractability} --- both the outcome and the payments must be computable
  (typically a MAS concern, which economics is not interested in).
\end{itemize}
It is not always possible to have all these properties at once: the Myerson--Satterthwaite theorem says
that for bilateral trade one cannot have IC, individual rationality, efficiency, and budget balance
simultaneously. VCG sacrifices the budget, other mechanisms sacrifice efficiency.
\end{defbox}

\begin{defbox}[The revelation principle]
For every mechanism in which the agents play some (possibly very complicated) equilibrium
strategies, there exists a \textbf{direct} mechanism in which the agents simply \emph{declare their
preferences truthfully} and the outcome is the same --- the \uv{lying} is built into the mechanism
itself.

\emph{The consequence for us:} in looking for a good mechanism it suffices to restrict attention to
truthful direct mechanisms. This is why the theory revolves around Vickrey and VCG.
\emph{The limits:} the Gibbard--Satterthwaite theorem (chapter~\ref{sec:volby}) shows that without
money (with ordinal preferences only) a truthful non-dictatorial mechanism is impossible ---
it is only \emph{payments} in a quasi-linear setting that allow Vickrey and VCG to circumvent this
barrier.
\end{defbox}

\subsection{The four basic types of auction}

\begin{center}\small
\begin{tabular}{@{}p{0.15\textwidth}
  >{\raggedright\arraybackslash}p{0.23\textwidth}
  >{\raggedright\arraybackslash}p{0.24\textwidth}
  >{\raggedright\arraybackslash}p{0.27\textwidth}@{}}
\toprule
\textbf{Auction} & \textbf{Protocol} & \textbf{Rational strategy} & \textbf{Notes} \\
\midrule
\textbf{English}
  & open cry, \emph{ascending} price, the first (highest) price is paid
  & \emph{dominant}: keep bidding up by $\varepsilon$ until the price reaches my valuation,
    then withdraw
  & The most common one (antiques, art, the internet --- about 95~\% of auctions). The typical case in
    which the buyer overpays. \\
\textbf{Dutch}
  & open cry, \emph{descending} price, the first to accept pays the current price
  & \emph{no} dominant strategy exists; wait until the price falls just below my valuation ---
    but this is a gamble
  & Fast; perishable goods (flowers, fish). Popular with sellers.
    It makes it impossible to estimate either the market price or the preferences of the others. \\
\textbf{FPSB} \newline (\emph{first-price sealed bid})
  & a single round, sealed envelopes, the highest bid wins and pays its own price
  & \emph{no} dominant strategy exists; bid \emph{less} than one's valuation
    (\emph{bid shading}) --- the optimum depends on an estimate of the others
  & Sales of real estate and government bonds. It does not force buyers to use their utility function
    and does not reveal the market price. \emph{Strategically equivalent to the Dutch auction.} \\
\textbf{Vickrey} \newline (\emph{second-price sealed bid})
  & a single round, sealed envelopes, the highest bid wins but pays the \emph{second} highest
  & \emph{dominant}: bid one's valuation \textbf{truthfully}
  & Theoretically the most popular one; Google AdWords, stamp auctions.
    \emph{Equivalent to the English auction} (with private values). \\
\bottomrule
\end{tabular}
\end{center}

\begin{defbox}[Why truthfulness is a dominant strategy in the Vickrey auction]
Let my valuation be $v$, let me bid $b$, and let the highest bid of the others be $p$ (which I do not
know). I win exactly when $b > p$, and then I pay $p$; my profit is $v - p$.
\begin{itemize}
\item \textbf{I bid $b > v$ (overstate):} compared with $b=v$ something changes only when
  $v < p < b$. Then I win, but I pay $p > v$ --- a profit of $v - p < 0$. \emph{This harms me.}
\item \textbf{I bid $b < v$ (understate):} compared with $b=v$ something changes only when
  $b < p < v$. Then I lose, although I would have won with a profit of $v - p > 0$. \emph{I forgo a
  profit.}
\item In no case \emph{can} I gain more by bidding $b \ne v$ than by bidding $b = v$.
\end{itemize}
Hence $b = v$ is a (weakly) dominant strategy --- regardless of what the others do.
The auction is \textbf{strategy-proof} (incentive compatible)
and at the same time \emph{efficient} (the bidder with the highest valuation wins).
\end{defbox}

\medskip
\noindent\textbf{Example --- a comparison of revenues.} Three buyers with valuations
$A = 80$~CZK, $B = 60$~CZK, $C = 30$~CZK. In an ideal scenario (with no additional information and no
cheating) the auctions turn out as follows:
\begin{center}\small
\begin{tabular}{lll}
\toprule
\textbf{Auction} & \textbf{Winner} & \textbf{Price} \\
\midrule
English & $A$ & $\approx 61$ (bidding continues until $B$ withdraws at 60) \\
Dutch & $A$ & $80$ or slightly less (79) --- $A$ must not risk accepting too late \\
FPSB & $A$ & $80$ (or $60+\varepsilon$ if $A$ knows the valuations of the others) \\
Vickrey & $A$ & $60$ (the second highest bid) \\
\bottomrule
\end{tabular}
\end{center}
In all cases the item is obtained by the agent with the highest valuation (the auctions are
\emph{efficient}), but they differ in the price and in the information the auction reveals.

\emph{Why does FPSB come out at 80 when one is supposed to bid below one's valuation?} Because the
scenario assumes \emph{zero} information about the others: agent $A$ does not know how high $B$ will
go, and any retreat from 80 risks losing. The degree of shading is a function of one's beliefs about
the others --- which is why FPSB (and the Dutch auction) \textbf{has no dominant strategy}. If $A$
knows the valuations of the others, it bids $60+\varepsilon$ and the seller's revenue drops to the
level of the Vickrey auction (cf.\ the revenue equivalence theorem below). In the English and Vickrey
auctions, by contrast, a dominant strategy does exist and does not depend on what the others do.

\begin{defbox}[The revenue equivalence theorem]
Vickrey (1961), Myerson, Riley--Samuelson. Under the assumptions: risk-neutral buyers,
independent private values from the same continuous distribution, the item is obtained by the buyer
with the highest valuation, and the buyer with the lowest possible valuation has zero expected
profit --- \textbf{all four basic auctions give the seller the same expected revenue.}

The differences in practice therefore stem from violations of the assumptions: risk aversion (in which
case FPSB/Dutch are better), common values and the winner's curse (in which case the English auction is
better because it reveals information), the possibility of collusion among buyers (the English auction
is more vulnerable), and the cost of communication (sealed bid is cheaper).
\end{defbox}

\subsection{Combinatorial auctions and VCG}

\begin{defbox}[Combinatorial auction]
\emph{Several items are auctioned simultaneously} and buyers submit bids on \emph{subsets}
(\emph{bundles}). The motivation: the values of items are not additive ---
\emph{complementarity} (a left and a right shoe, two adjacent frequency blocks:
$v(\{x,y\}) > v(x)+v(y)$) and \emph{substitution} (two tickets for the same flight:
$v(\{x,y\}) < v(x)+v(y)$).

\textbf{The winner determination problem} (WDP): select a set of mutually disjoint bids that maximises
the sum of the prices. This is a generalisation of set packing --- an \textbf{NP-hard}
problem (and moreover hard to approximate). It is solved with ILP solvers and branch and bound
(CABOB, CPLEX).

The second problem is \emph{representational}: there are $2^m - 1$ possible bids on subsets;
compact languages are used (OR, XOR, OR-of-XOR).
\end{defbox}

\begin{defbox}[The VCG mechanism (Vickrey--Clarke--Groves)]
A generalisation of the Vickrey auction to an arbitrary allocation. The principle:
\begin{enumerate}
\item Every agent $i$ declares its valuation $\hat v_i(\cdot)$.
\item An allocation $x^*$ maximising \emph{social welfare}
  $\sum_i \hat v_i(x)$ is chosen.
\item Agent $i$ pays the \textbf{externality it imposes on the others}:
\[
p_i \;=\; \underbrace{\max_{x} \sum_{j \ne i} \hat v_j(x)}_{\text{the welfare of the others if } i
\text{ did not exist}} \;-\; \underbrace{\sum_{j \ne i} \hat v_j(x^*)}_{\text{the welfare of the others
under the chosen allocation}} .
\]
\end{enumerate}
\textbf{Properties:} truthfulness is a dominant strategy (\emph{incentive compatible}),
the allocation maximises social welfare (\emph{efficient}), and the agents have a non-negative profit
(\emph{individually rational}). For a single item VCG reduces exactly to the Vickrey auction.

\textbf{Disadvantages:} computational cost (the WDP has to be solved $n+1$ times), a possible zero or
negative revenue for the seller, vulnerability to collusion and false identities (\emph{shill
bidding}), and the need to reveal one's true valuations (a privacy problem).
\end{defbox}

\medskip
\noindent\textbf{A VCG example.} The items $\{x, y\}$ are auctioned. The bids:
agent 1: $\{x\} \to 5$; agent 2: $\{y\} \to 4$; agent 3: $\{x,y\} \to 8$.
\begin{itemize}
\item The optimal allocation: $x \to$ agent 1, $y \to$ agent 2, total welfare $5+4=9 > 8$.
\item The payment of agent 1: without agent 1 the best allocation would be \uv{$\{x,y\} \to$ agent 3}
  with value $8$; under the allocation $x^*$ the others have value $4$ (agent 2 alone).
  Hence $p_1 = 8 - 4 = 4$. The profit of agent 1 is $5 - 4 = 1$.
\item Symmetrically $p_2 = 8 - 5 = 3$, and the profit of agent 2 is $4-3=1$.
\item The seller's total revenue is $4+3 = 7$, which is \emph{less} than the $8$ that agent 3 offered ---
  a typical illustration that VCG does not maximise the seller's revenue.
\end{itemize}

\subsection{Further types of auction and resource allocation}

\begin{itemize}
\item An \emph{all-pay auction} --- everybody pays and the highest bid wins; a model of lobbying,
  bribes, sporting contests, and patent races.
\item A \emph{silent} auction --- a written variant of the English one. The \emph{Amsterdam} auction ---
  it starts as an English one, and once two bidders remain it switches to a Dutch one with a doubled
  price. \emph{Tsukiji} (the Tokyo fish market) --- everybody bids at once, and ties are resolved by
  rock--paper--scissors.
\item A \emph{double auction} --- bids are submitted by buyers and sellers alike and are matched; a
  model of a stock exchange.
\end{itemize}

\begin{defbox}[Resource allocation in general (MARA)]
\emph{Multiagent resource allocation}: dividing a set of resources among agents according to their
preferences. An auction is only one of the mechanisms (the \emph{market} one); others are centralised
optimisation, negotiation, or fair division. Criteria for the quality of an allocation:
\begin{itemize}
\item \textbf{utilitarian} --- maximises $\sum_i u_i$ (social welfare);
\item \textbf{egalitarian} --- maximises $\min_i u_i$ (the worst-off agent);
\item \textbf{Nash welfare} --- maximises $\prod_i u_i$ (a compromise);
\item \textbf{proportionality} --- each of the $n$ agents receives at least $1/n$ of its total value;
\item \textbf{envy-freeness} --- no agent prefers another agent's share to its own.
\end{itemize}
For divisible resources (\uv{cake cutting}) there are classical protocols
(\emph{cut and choose} for two agents, \emph{Dubins--Spanier} for $n$); for indivisible
items both proportionality and envy-freeness are in general unattainable.
\end{defbox}

\subsection{Negotiation}

An auction requires a central auctioneer. The alternative is direct \textbf{negotiation}
between agents. The classical model:

\begin{defbox}[The monotonic concession protocol and the Zeuthen strategy]
The \textbf{negotiation set} --- the individually rational and Pareto-optimal deals.

The \textbf{monotonic concession protocol}: in each round both agents simultaneously propose a deal. It
ends if one agent's proposal is at least as good for the other as that agent's own proposal; otherwise
at least one of them must \emph{concede} in the next round (propose a deal that is better for the
counterpart), or else the negotiation breaks down.

The \textbf{Zeuthen strategy} answers the question of \emph{who should concede}: the one who has the
smaller \uv{willingness to risk} a conflict,
\[
\mathrm{Risk}_i = \frac{u_i(\text{own proposal}) - u_i(\text{counterpart's proposal})}
{u_i(\text{own proposal})},
\]
that is, the ratio of what it loses by accepting the other's proposal to what it would lose by a
conflict. The agent with the \emph{lower} $\mathrm{Risk}$ concedes. It should concede just enough for
the ratios to be reversed. A pair of Zeuthen strategies is a Nash equilibrium, and the resulting deal
maximises the \emph{Nash bargaining solution} (the product of the utilities).
\end{defbox}

Rosenschein and Zlotkin introduced for this type of negotiation the formalism of
\textbf{task-oriented domains} (TOD): a pair of agents is assigned sets of tasks, a \emph{deal}
$\delta$ is a redistribution of the tasks between them, and the utility $u_i(\delta)$ is the cost the
agent saves by the deal compared with performing its own tasks alone; the \emph{conflict deal} $\Theta$
means that they fail to agree and each carries its own tasks itself. In TOD it may pay to
\textbf{deceive}: an agent may declare \emph{phantom} and \emph{decoy} tasks that it does not have (in
order to look more heavily loaded and concede less), or conversely \emph{conceal} some of its tasks
(\emph{hidden tasks}); whether lying pays off depends on the structural properties of the domain.
Further models of negotiation are \textbf{Rubinstein's alternating offers protocol}, in which the
patience of the agents is modelled by a discount factor and in equilibrium the deal is struck right in
the first round, and \emph{argumentation-based negotiation}, in which the agents exchange not only
proposals but also reasons (threats, promises, appeals).

\subsection{Exam summary}

\begin{itemize}
\item An auction = a mechanism for allocating scarce resources; an efficient auction gives the resource
  to whoever values it most. The dimensions: private/common value, open cry/sealed bid, first/second
  price, the number of rounds.
\item \textbf{Mechanism design} = inverse game theory (I design the rules, not the strategy).
  The properties: incentive compatibility, individual rationality, efficiency, budget
  balance, computability --- one cannot have everything at once. \textbf{The revelation principle}: it
  suffices to consider direct truthful mechanisms. Without payments it runs into
  Gibbard--Satterthwaite.
\item The four basic types and their strategies: the English auction (bid up by $\varepsilon$
  up to your valuation) and Vickrey (bid truthfully) have a \emph{dominant} strategy;
  the Dutch auction (wait for your valuation $-\varepsilon$) and FPSB (bid below your valuation)
  \textbf{do not}. Be able to give the \textbf{proof of truthfulness} for Vickrey.
\item Equivalences: English $\leftrightarrow$ Vickrey, Dutch $\leftrightarrow$ FPSB.
  \textbf{The revenue equivalence theorem} and when it fails (risk aversion, common values,
  the winner's curse, collusion).
\item Be able to do the example with valuations 80/60/30 and compute the price in all four auctions.
\item Combinatorial auctions: bids on subsets (complementarity/substitution),
  the WDP is NP-hard. \textbf{VCG}: the allocation maximising social welfare + a payment equal to
  the externality the agent imposes on the others; truthful and efficient, but expensive and with a low
  revenue.
\item Resource allocation: the utilitarian / egalitarian / Nash criterion, proportionality,
  envy-freeness.
\item Negotiation: the monotonic concession protocol + the Zeuthen strategy (the one with the lower
  risk concedes); task-oriented domains and the three types of deception (phantom / decoy / hidden
  tasks); Rubinstein's alternating offers with discounting.
\end{itemize}

\section{Agent learning and reinforcement learning}
\label{sec:uceni}

\subsection{Why and how agents learn}

Intelligent agents in MAS should be \emph{adaptive}. Learning is a change in an agent's behaviour on
the basis of experience; there is a whole range of methods:

\begin{defbox}[Methods for agent learning]
\begin{itemize}
\item \textbf{By the learning signal} (the three classical paradigms of machine learning):
  \emph{supervised learning} (the agent receives the correct answers --- e.g.\ learning a model of the
  opponent from observed moves, classifying percepts), \emph{unsupervised learning}
  (finding structure in observations --- clustering situations, discovering roles),
  and \textbf{reinforcement learning} (only a scalar reward from the environment is available;
  the most natural for autonomous agents, see the rest of this chapter).
\item \textbf{Imitation learning} (\emph{learning from demonstration}): the agent learns a
  policy from trajectories demonstrated by an expert --- either by direct copying
  (\emph{behavioural cloning}), or by \emph{inverse RL}, in which the reward function is
  reconstructed from the expert's behaviour.
\item \textbf{Evolutionary methods}: a population of controllers (parameters of reactive behaviours,
  neural network weights, rules) is bred by a genetic algorithm according to the utility
  achieved --- the typical way in which the reactive architectures of
  chapter~\ref{sec:reaktivni} are \uv{tuned experimentally}; this also includes \emph{learning
  classifier systems} and neuroevolution.
\item \textbf{By what the agent learns}: a model of the environment (the transition function),
  a model of the opponent (\emph{opponent modelling}), a value function, the policy directly,
  or, in PRS, the ratings of plans.
\item \textbf{Individual vs.\ social learning}: the agent learns only from its own
  experience, or also by imitating and sharing experience with the others (learning
  in a swarm, shared experience in MARL).
\end{itemize}
\end{defbox}

The most common and theoretically most developed approach to learning in MAS
is \textbf{reinforcement learning} (RL) --- a change of behaviour by trial and error on the basis of
\emph{rewards} from the environment. RL fits naturally into the formalism of chapter
\ref{sec:architektura}: the policy $\pi$ is an action selection mechanism and the reward
is an operationalisation of utility.

We distinguish:
\begin{itemize}
\item \textbf{Single-agent learning} --- easier, and classical RL algorithms can be used (Q-learning).
\item \textbf{Multiagent RL} (\emph{MARL}) --- the agents learn \emph{simultaneously}; this is
  formalised by Markov games (Nash and Pareto optimality), and the state of the art is deep MARL
  (MADRL).
\end{itemize}

\subsection{The Markov decision process}

\begin{defbox}[MDP --- \emph{Markov Decision Process}]
An MDP is a quintuple $\langle S, A, T, R, \gamma\rangle$:
\begin{itemize}
\item $S$ --- the state space; $A$ --- the action space;
\item $T : S \times A \times S \to [0,1]$ --- the transition function,
  $T(s,a,s') = P(s_{t+1}=s' \mid s_t=s, a_t=a)$;
\item $R : S \times A \times S \to \mathbb{R}$ --- the reward function (the immediate reward for a
  transition);
\item $0 \le \gamma < 1$ --- the discount factor, a trade-off between immediate and future reward.
\end{itemize}
The \textbf{Markov property} holds: the next state and reward depend only on the current state
and action, not on the whole history. (Compare this with the state transformer function
$\tau : R^A \to 2^E$ from the abstract architecture, which \emph{does} have a history --- an MDP is
its Markovian simplification enriched with probabilities and rewards.)

A \textbf{policy} $\pi : S \to A$ (or stochastically $\pi(a|s)$).
The goal is to find an \emph{optimal policy} $\pi^*$ maximising the expected discounted sum of
future rewards:
\[
\mathbb{E}\Big[ \sum_{t=0}^{\infty} \gamma^t R(s_t, a_t, s_{t+1}) \Big].
\]
The \textbf{value function} $V^\pi : S \to \mathbb{R}$ assigns to a state the expected utility
if the agent follows the policy $\pi$; the \textbf{action-value function}
$Q^\pi(s,a)$ does the same for a state--action pair.
\end{defbox}

\begin{defbox}[The Bellman equations]
For a given policy (the \emph{Bellman equation for evaluation}):
\[
V^\pi(s) = \sum_{s'} T(s,\pi(s),s') \big[ R(s,\pi(s),s') + \gamma V^\pi(s') \big].
\]
For the optimal policy (the \emph{Bellman optimality equation}):
\begin{align*}
V^*(s) &= \max_{a \in A} \sum_{s'} T(s,a,s') \big[ R(s,a,s') + \gamma V^*(s')\big], \\
Q^*(s,a) &= \sum_{s'} T(s,a,s')\big[R(s,a,s') + \gamma \max_{a'} Q^*(s',a')\big].
\end{align*}
The optimal policy is obtained from $Q^*$ as $\pi^*(s) = \arg\max_a Q^*(s,a)$.
\end{defbox}

\noindent\textbf{Value iteration} --- if we know the complete description of the MDP:
\begin{algorithmic}[1]
\State initialise $V_0(s) \gets 0$ for all $s$
\Repeat
  \ForAll{$s \in S$}
    \State $V_{k+1}(s) \gets \max_{a} \sum_{s'} T(s,a,s')\big[R(s,a,s') + \gamma V_k(s')\big]$
  \EndFor
\Until{$\max_s |V_{k+1}(s) - V_k(s)| < \varepsilon$}
\State \Return $\pi(s) = \arg\max_a \sum_{s'} T(s,a,s')[R(s,a,s') + \gamma V(s')]$
\end{algorithmic}
The Bellman operator is a contraction with coefficient $\gamma$, so the iteration converges to $V^*$
geometrically fast. An alternative is \textbf{policy iteration}:
alternately \emph{evaluate} the current policy (solve a system of linear equations) and \emph{improve}
it (greedily according to $V^\pi$); it converges in finitely many steps, because there are finitely
many policies.

\subsection{Q-learning and SARSA}

In reality we usually \textbf{do not know} either $T$ or $R$. The agent therefore learns from
experience by interacting with the environment in discrete steps --- this is \emph{model-free} RL.

\begin{defbox}[Q-learning (Watkins, 1989)]
The agent maintains a \emph{table} $Q(s,a)$ --- an estimate of the discounted sum of future rewards if
action $a$ is performed in state $s$. After every transition $s \xrightarrow{a} s'$ with reward $r$:
\[
Q(s,a) \;\leftarrow\; Q(s,a) + \alpha\Big[\, \underbrace{r + \gamma \max_{a'} Q(s',a')}
_{\text{the target (TD target)}} - Q(s,a) \Big],
\qquad 0 \le \alpha \le 1 \text{ is the learning rate,}
\]
equivalently $Q(s,a) \leftarrow (1-\alpha)Q(s,a) + \alpha\big(r + \gamma\max_{a'}Q(s',a')\big)$.

\textbf{Off-policy:} the target uses $\max_{a'}$, so it learns the value of the \emph{optimal} policy
regardless of the policy the agent actually acts by. Provided that every pair
$(s,a)$ is visited infinitely often and $\alpha$ decreases suitably, $Q \to Q^*$.
\end{defbox}

\begin{defbox}[The exploration--exploitation dilemma; $\varepsilon$-greedy]
The agent has to \emph{exploit} what it already knows, but at the same time \emph{explore}
new actions. The standard solution:
\[
\pi(s) = \begin{cases}
\text{a random action from } A & \text{with probability } \varepsilon,\\
\arg\max_a Q(s,a) & \text{with probability } 1-\varepsilon .
\end{cases}
\]
$\varepsilon$ is typically decreased over time (\emph{annealing}). Alternatives:
softmax/Boltzmann $\pi(a|s) \propto e^{Q(s,a)/\tau}$, optimistic initialisation of $Q$,
UCB (upper confidence bounds).
\end{defbox}

\medskip
\noindent\textbf{A concrete Q-learning step.} Let $\alpha = 0.5$, $\gamma = 0.9$.
The agent is in state $s_1$, performs action $a_1$, receives reward $r = 10$, and moves to $s_2$.
The current values: $Q(s_1,a_1) = 4$, $Q(s_2,a_1) = 6$, $Q(s_2,a_2) = 8$.
\begin{align*}
\max_{a'} Q(s_2,a') &= \max(6,\,8) = 8, \\
\text{TD target} &= r + \gamma \max_{a'} Q(s_2,a') = 10 + 0.9 \cdot 8 = 17.2, \\
\text{TD error} &= 17.2 - Q(s_1,a_1) = 17.2 - 4 = 13.2, \\
Q(s_1,a_1) &\leftarrow 4 + 0.5 \cdot 13.2 = \mathbf{10.6}.
\end{align*}

\begin{defbox}[SARSA]
\emph{State--Action--Reward--State--Action}. The update uses the \emph{actually chosen}
next action $a'$ (according to the same policy the agent acts by):
\[
Q(s,a) \leftarrow Q(s,a) + \alpha\big[r + \gamma\, Q(s',a') - Q(s,a)\big].
\]
\textbf{On-policy:} it learns the value of the policy the agent actually executes (including random
exploration). The consequence: SARSA is \emph{more cautious} --- in the \emph{cliff walking} task
Q-learning learns the optimal path right at the edge of the precipice (and with
$\varepsilon$-greedy it occasionally falls in), whereas SARSA learns a safer path further from the
edge, because it factors in the risk of its own exploration.

A generalisation of both: $\mathrm{TD}(\lambda)$ with \emph{eligibility traces}, which spread
credit over several preceding steps.
\end{defbox}

\subsection{Policy-based methods}

\begin{defbox}[Policy gradient, REINFORCE]
Instead of value functions we optimise \emph{the policy directly}. We parameterise the policy by a
vector $z$, i.e.\ $\pi(a \mid s, z)$, and look for the optimal parameters $z^*$
by gradient ascent on the expected return $J(z)$.

\textbf{REINFORCE} (Williams, 1992) --- estimates the gradient from Monte Carlo simulations of whole
episodes:
\[
z_{t+1} = z_t + \alpha\, G_t\, \frac{\nabla_z \pi(A_t \mid S_t, z_t)}{\pi(A_t \mid S_t, z_t)}
= z_t + \alpha\, G_t\, \nabla_z \ln \pi(A_t \mid S_t, z_t),
\]
where $G_t$ is the return from time $t$ onwards. The intuition: increase the probability of actions
that were followed by a high return.

\textbf{Advantages over value-based methods:} it handles \emph{continuous} action spaces, naturally
produces \emph{stochastic} policies (necessary in multiagent games with mixed equilibria),
and has better convergence properties under function approximation.
\textbf{A disadvantage:} the high variance of the gradient estimate (addressed by subtracting a
\emph{baseline}).
\end{defbox}

\begin{defbox}[Actor--critic]
A combination of both approaches:
\begin{itemize}
\item The \textbf{actor} represents the \emph{policy} --- the action selection mechanism; it optimises
  it.
\item The \textbf{critic} learns the \emph{value function} --- it estimates how good the actor's
  performance is and thereby reduces the variance.
\end{itemize}
If the critic learns both $Q(s,a)$ and $V(s)$, the \textbf{advantage function} can be used:
\[
A(s,a) = Q(s,a) - V(s),
\]
i.e.\ the relative advantage of action $a$ over the average --- the policy gradient is then scaled by
$A(s,a)$ instead of $G_t$, which substantially reduces the variance.

\textbf{A3C} (\emph{Asynchronous Advantage Actor-Critic}) --- several worker threads
asynchronously compute gradients over their own copies of the environment and update shared
parameters; the loss function consists of the \emph{policy loss} for the actor and the
\emph{value loss} for the critic (plus an entropy regulariser encouraging exploration). More modern
variants: A2C (synchronous), PPO, SAC, DDPG.
\end{defbox}

\textbf{Deep learning.} \textbf{DQN} (\emph{Deep Q-Network}) replaces the $Q$ table
with a neural network (feed-forward, or recurrent for partially observable environments);
the key tricks are \emph{experience replay} (breaking the correlations between samples) and a
\emph{target network} (stabilising the target). In policy-based methods the policy is represented by a
neural network (the input is the state, the output an action or its distribution, and the weights are
the parameters $z$).

\subsection{Multiagent reinforcement learning}

\begin{defbox}[Markov game (stochastic game)]
A generalisation of the MDP to $N$ agents: $\langle N, S, (A_i)_{i\in N}, T, (R_i)_{i \in N},
\gamma\rangle$.
\begin{itemize}
\item The \textbf{joint action space} is the product of the individual ones:
  $\mathbf{A} = A_1 \times \cdots \times A_N$.
\item $T : S \times \mathbf{A} \times S \to [0,1]$ --- the transition depends on the actions of
  \emph{all} agents.
\item Each agent has \emph{its own} reward function $R_i$ and its own policy $\pi_i$,
  but its value function $V^{\boldsymbol\pi}_i$ depends on the policies of \emph{all} agents.
\end{itemize}
\textbf{The connection with game theory:} an agent's policy is a \emph{best response} to the joint
policies of the others if it maximises its value function. The joint policies may be
in a \textbf{Nash equilibrium} and may be \textbf{Pareto optimal}.
Special cases: $N=1$ gives an MDP; a game with a single state gives a normal-form game;
$N=2$ with zero sum gives a \emph{minimax} game.
\end{defbox}

\begin{defbox}[Minimax-Q and Nash-Q]
\textbf{Minimax-Q} (Littman, 1994) --- for two-agent zero-sum games. Instead of
$\max_{a'} Q(s',a')$, the update uses the value of the \emph{mixed minimax equilibrium}
of the matrix $Q(s',\cdot,\cdot)$:
\[
V(s') = \max_{\sigma \in \Delta(A_1)} \min_{a_2 \in A_2} \sum_{a_1} \sigma(a_1)\, Q(s',a_1,a_2),
\]
which is a linear program. It converges to the minimax (and hence Nash) policy regardless of the
opponent's behaviour --- at the price of \emph{conservatism} (it will not exploit a weaker opponent).

\textbf{Nash-Q} (Hu, Wellman) --- for general-sum games; the agent maintains $Q_i$ for
\emph{all} agents and in the update uses the value of the Nash equilibrium of the stage game.
It converges only under strong (and practically unverifiable) assumptions on the uniqueness of the
equilibria. Related: \emph{Friend-or-Foe Q}, \emph{correlated-Q}.
\end{defbox}

\begin{defbox}[The problems of MARL]
Game theory provides elegant formalisations (for few agents), but practice runs into the following:
\begin{itemize}
\item the environment is only \textbf{partially observable} (formally a Dec-POMDP, whose optimal
  solution is NEXP-complete);
\item the state and action spaces may be \textbf{continuous}; the joint action space grows
  \emph{exponentially} with the number of agents;
\item the agents update their policies \textbf{simultaneously}, so from the point of view of a single
  agent the environment is \textbf{non-stationary} --- the basic convergence assumption of RL is
  violated (\uv{a moving target});
\item the agents may converge to \textbf{suboptimal} solutions or \textbf{oscillate} between several
  Nash equilibria;
\item the \textbf{variance} of the value function estimates may be large;
\item \textbf{credit assignment} --- with a shared reward it is not clear
  which agent contributed.
\end{itemize}
\textbf{The standard countermeasures:} \emph{centralised training with decentralised execution}
(CTDE) --- during training the critic sees the state and the actions of all agents, while at run time
the actor sees only its own local observations (MADDPG, COMA, QMIX); furthermore \emph{parameter
sharing} between homogeneous agents and \emph{opponent modelling}.
\end{defbox}

\textbf{Self-play.} The agent learns by playing against itself (or possibly against earlier versions
of itself --- \emph{fictitious self-play}, leagues of opponents). In zero-sum games with perfect
information, self-play reliably leads to strong strategies: TD-Gammon (Tesauro, 1992),
AlphaGo Zero and AlphaZero (self-play + MCTS search + a deep network), AlphaStar
(a league of opponents for StarCraft~II). The risk: \emph{cyclic} preferences among strategies
(rock--paper--scissors), where naive self-play oscillates --- hence leagues and population-based
methods (\emph{Policy-Space Response Oracles}).

\textbf{LLM agents} (a current trend). Large language models are used as an action selection
mechanism (a policy), as a planner, as a critic in an actor-critic scheme, and for learning with a
human in the loop. Typical applications in agent-based simulations: the spread of information in a
graph of social interactions and synthetic electoral panels (LLM agents represent a population
according to sociological surveys and estimate election results).

\subsection{Exam summary}

\begin{itemize}
\item Methods of agent learning: supervised / unsupervised / reinforcement; imitation learning
  and inverse RL; evolutionary methods (LCS, neuroevolution); what the agent learns (a model of the
  environment, a model of the opponent, a policy); individual vs.\ social learning. RL is the most
  common.
\item \textbf{MDP} $= \langle S,A,T,R,\gamma\rangle$, the Markov property, a policy $\pi$,
  the value functions $V^\pi$, $Q^\pi$; the Bellman optimality equation; value iteration
  (a contraction, converges) and policy iteration.
\item \textbf{Q-learning}: $Q(s,a) \leftarrow Q(s,a) + \alpha[r + \gamma\max_{a'}Q(s',a') - Q(s,a)]$,
  off-policy, tabular; be able to compute one step with numbers.
\item \textbf{SARSA}: $\dots + \alpha[r + \gamma Q(s',a') - Q(s,a)]$, on-policy, more cautious
  (cliff walking). $\varepsilon$-greedy as the solution of the exploration/exploitation dilemma.
\item \textbf{Policy gradient} (REINFORCE, $\nabla \ln \pi$ times the return) and
  \textbf{actor-critic} (the actor = the policy, the critic = the value function, the advantage
  $A = Q - V$); A3C; DQN (replay buffer, target network).
\item \textbf{Markov game} = an MDP for $N$ agents with a joint action space and individual
  rewards; best response, Nash equilibrium, Pareto optimality of the joint policies.
\item \textbf{Minimax-Q} (zero-sum games, an LP in every update), Nash-Q
  (general sum, strong assumptions).
\item The problems of MARL: non-stationarity, the exponential joint action space, partial
  observability, credit assignment; CTDE as the solution. Self-play and AlphaZero.
\end{itemize}

\section{Design methodologies, languages, and MAS platforms}
\label{sec:metodologie}

\subsection{Agent-oriented software engineering}

Object-oriented methodologies (UML) are not sufficient, because agents are autonomous (their methods
cannot be called directly), they have their own goals, and the relations between them are
\emph{organisational} (roles, protocols, commitments) rather than structural. This is why the
methodologies of \textbf{AOSE} (agent-oriented software engineering) arose.

\begin{defbox}[Gaia (Wooldridge, Jennings, Kinny, 2000)]
One of the most cited AOSE methodologies. The key metaphor: a MAS is an \textbf{organisation} and
agents play \textbf{roles}. It proceeds from requirements through analysis to design (leaving the
implementation to the particular platform).

\textbf{The analysis phase} --- two abstract models. The \emph{roles model}: every
role is described by four attributes --- \emph{permissions} (which resources it may read/modify),
\emph{responsibilities} (what it has to ensure; these split into \emph{liveness} properties
\uv{something good happens}, described by a regular expression over activities and protocols, and
\emph{safety} properties \uv{nothing bad happens}, described by invariants), \emph{activities}
(computations performed by the role itself), and \emph{protocols} (interactions with other roles). The
\emph{interactions model}: the definition of the protocols between roles (purpose, initiator,
responder, inputs, outputs).

\textbf{The design phase} --- three concrete models: the \emph{agent model} (the assignment of roles to
agent types and the numbers of instances), the \emph{services model} (the functions provided by a role
--- inputs, outputs, preconditions, effects), and the \emph{acquaintance model} (a directed graph of who
communicates with whom --- a check for bottlenecks and excessive coupling).

\textbf{Limitations:} the original Gaia assumes a closed system with cooperative agents and a static
organisation; the extensions \emph{ROADMAP} and \emph{Gaia v.2} relax these limitations.
\end{defbox}

\begin{defbox}[Prometheus and Tropos]
\textbf{Prometheus} (Padgham, Winikoff, 2002) --- a practical methodology focused on
\textbf{BDI agents}, with the PDT tool and the generation of code skeletons for JACK/Jason. Three
phases: \emph{system specification} (actors and scenarios, goals, \emph{functionalities} as coherent
units of behaviour, the interface with the surroundings --- \emph{percepts} and \emph{actions});
\emph{architectural design} (from the functionalities, \emph{agent types} are derived according to
cohesion and coupling, and protocols and messages are defined); \emph{detailed design} (the inside of
an agent: \emph{capabilities}, plans, events, and data structures --- directly the building blocks of
BDI).

\textbf{Tropos} (Castro, Giorgini, Mylopoulos, 2002) --- built on the \textbf{modelling of actors
and goals} (the $i^*$ framework); it is the only one that also covers the \emph{early requirements}
phase. Five phases: early requirements, late requirements, architectural design, detailed design,
implementation. The concepts: \emph{actor}, \emph{goal} (and \emph{softgoal} = a non-functional
requirement with no sharp satisfaction criterion), \emph{plan}, \emph{resource}, and \emph{dependency}
(an actor depends on another one for a goal, a plan, or a resource). The analysis proceeds by means-ends
decomposition of goals and by \emph{contribution links} ($+$/$-$) to softgoals; the same concepts
pervade the entire development from business analysis down to code.
\end{defbox}

\textbf{Other methodologies} (for orientation only): \textbf{MaSE} (from use cases to automata,
with the agentTool tool), \textbf{INGENIAS} (a metamodel + code generation),
\textbf{ADELFE} (adaptive/emergent systems), \textbf{O-MaSE}, \textbf{PASSI}.

\subsection{The FIPA standard and the architecture of a platform}

\begin{defbox}[The FIPA reference architecture of an agent platform]
\begin{itemize}
\item \textbf{AMS} (\emph{Agent Management System}) --- a mandatory component; the \uv{registry} of the
  platform: it keeps records of agents, assigns them unique \emph{AIDs} (agent identifiers),
  manages the life cycle (create, suspend, resume, terminate), and provides the
  \emph{white pages}.
\item \textbf{DF} (\emph{Directory Facilitator}) --- the \emph{yellow pages}:
  agents register with it the \emph{services} they offer and look up service providers.
\item \textbf{MTS} (\emph{Message Transport Service}) --- the delivery of ACL messages within a
  platform and between platforms (the IIOP and HTTP protocols; bit-efficient, XML, and string
  encodings).
\item \textbf{ACC} (\emph{Agent Communication Channel}) --- the interface to the MTS.
\end{itemize}
The standard further defines the ACL language, content languages (SL), interaction protocols, and
ontology management.
\end{defbox}

\subsection{Languages and platforms}

The historical starting point is \textbf{agent-oriented programming}
(AOP; Shoham 1993) --- a paradigm in which the state of a program consists of
\emph{mental categories} (beliefs, \emph{commitments}, and capabilities) and the program is a set of
\emph{commitment rules} which, on the basis of incoming messages and the current mental
state, create new commitments. The first AOP language was Shoham's \textbf{AGENT-0}:
an agent = capabilities + initial beliefs + initial commitments + commitment rules;
it communicates with three types of message (\texttt{request}, \texttt{unrequest},
\texttt{inform}). A different route is taken by \textbf{Concurrent MetateM} (Fisher): an agent is a
\emph{directly executable specification} in temporal logic --- rules of the form
\uv{past $\Rightarrow$ future} are evaluated at run time and the agent acts so as to
make its specification true. Both languages continue the line of deductive agents
(chapter~\ref{sec:deliberativni}); the most successful representative of agent
languages in practice today is AgentSpeak/Jason (see the table).

\begin{center}\small
\begin{tabular}{p{0.15\textwidth} p{0.20\textwidth} p{0.57\textwidth}}
\toprule
\textbf{System} & \textbf{Type} & \textbf{Characteristics} \\
\midrule
\textbf{JADE}
  & Java, a FIPA platform
  & \emph{Java Agent DEvelopment Framework}, one of the most widespread platforms for MAS, fully
    compliant with FIPA. An agent = a class inheriting from \texttt{Agent} with a \texttt{setup()}
    method; behaviour is composed of \texttt{Behaviour} objects (\texttt{OneShot}, \texttt{Cyclic},
    \texttt{Ticker}, \texttt{FSMBehaviour}, \texttt{SequentialBehaviour}), which are scheduled by a
    cooperative (non-preemptive) scheduler inside the agent's single thread. Messages:
    \texttt{ACLMessage}, \texttt{send()}/\texttt{blockingReceive()}, \texttt{MessageTemplate}.
    Ready-made classes for protocols (\texttt{ContractNetInitiator}, \texttt{AchieveREResponder}),
    debugging tools (RMA, Sniffer, Introspector), containers, and agent mobility. \\
\textbf{Jason}
  & an AgentSpeak(L) interpreter in Java
  & The reference implementation of \textbf{AgentSpeak}, that is, of a \textbf{BDI} language (logic
    programming + BDI), a direct descendant of PRS. An agent program = initial \emph{beliefs},
    \emph{rules}, and a \emph{plan library} of the form \texttt{trigger : context <- body}. The trigger
    is an \emph{event} (\texttt{+b} the addition of a belief, \texttt{+!g} the addition of a goal,
    \texttt{-!g} the failure of a goal). The interpreter cycle: perceive $\to$ update the beliefs $\to$
    select an event $\to$ find the \emph{relevant} plans (by trigger) $\to$ select the
    \emph{applicable} ones (by context) $\to$ create an intention $\to$ execute a step. Together with
    \texttt{CArtAgO} (the environment) and \texttt{Moise} (the organisation) it forms \emph{JaCaMo}. \\
\textbf{NetLogo}
  & an environment for ABM
  & A language derived from Logo for \emph{agent-based modelling}. The entities: \emph{turtles} (mobile
    agents), \emph{patches} (cells of the environment), \emph{links}, and the \emph{observer}. A model
    library (Segregation, Wolf Sheep Predation, Ants, Flocking) and \emph{BehaviorSpace} for parametric
    studies. Intended for emergent behaviour and for teaching, not for FIPA communication. \\
\textbf{SPADE}
  & Python, XMPP
  & \emph{Smart Python Agent Development Environment}. Communication via standard
    \textbf{XMPP} instead of a proprietary MTS; an agent = a class with \emph{behaviours}
    (\texttt{CyclicBehaviour}, \texttt{PeriodicBehaviour}, \texttt{FSMBehaviour}), asyncio, and a
    web interface; popular thanks to its connection to the Python machine learning ecosystem. \\
\bottomrule
\end{tabular}
\end{center}

\textbf{Other environments:} \emph{Jadex}, \emph{JACK}, and \emph{JAM} (further implementations of
BDI/PRS, with Jadex building a BDI layer on top of JADE); \emph{MASON}, \emph{Repast}, \emph{GAMA},
\emph{AnyLogic}, and \emph{Mesa} (simulation platforms for agent-based models in sociology, economics,
and epidemiology); \emph{PettingZoo}, \emph{Gymnasium}, \emph{OpenSpiel}, and \emph{SMAC} (standard
environments for multiagent reinforcement learning); \emph{LangGraph} and \emph{AutoGen} (the
orchestration of LLM agents).

\subsection{Exam summary}

\begin{itemize}
\item Why OO methodologies are not sufficient: autonomy, own goals, organisational (rather than
  structural) relations.
\item \textbf{Gaia}: organisation + roles; analysis = the roles model (permissions, responsibilities
  with liveness and safety, activities, protocols) + the interactions model; design = the agent,
  services, and acquaintance models. The assumption of cooperative agents and a static organisation.
\item \textbf{Prometheus}: system specification (goals, functionalities, percepts/actions) $\to$
  architectural design (agent types, protocols, messages) $\to$ detailed design
  (capabilities, plans, events) --- targeted at BDI.
\item \textbf{Tropos}: actors, goals and softgoals, dependencies; it also covers early requirements
  ($i^*$).
\item \textbf{The FIPA platform}: AMS (white pages, life cycle, AID), DF (yellow pages,
  services), MTS/ACC (the delivery of ACL messages).
\item Agent-oriented programming (Shoham): the state = mental categories, the program =
  commitment rules; AGENT-0; Concurrent MetateM = an executable temporal specification.
\item \textbf{JADE} = Java, FIPA, \texttt{Agent} + \texttt{Behaviour} + \texttt{ACLMessage},
  ready-made protocols and debugging tools. \textbf{Jason/AgentSpeak} = a BDI language,
  a plan of the form \texttt{trigger : context <- body}, a descendant of PRS. \textbf{NetLogo} = ABM
  simulation (turtles/patches/links). \textbf{SPADE} = Python + XMPP.
\end{itemize}

\section*{In closing: how to approach the exam}
\addcontentsline{toc}{section}{In closing: how to approach the exam}

The topic is broad; the examiner usually checks that you understand the \emph{connections}. There are
four axes around which the questions revolve:
\begin{enumerate}
\item \textbf{Reactivity vs.\ proactiveness} --- from the definition of an intelligent agent through
  commitment strategies (bold vs.\ cautious) and the parameters $\gamma/\phi$ in the Maes network all
  the way to the layers of hybrid architectures.
\item \textbf{Individual vs.\ collective rationality} --- a selfish agent leads to $(D,D)$ in the
  prisoner's dilemma, to the tragedy of the commons, and to tactical voting; hence the need for
  mechanism design (Vickrey, VCG), that is, for designing \emph{rules} under which it pays a selfish
  agent to be truthful.
\item \textbf{Symbolic vs.\ subsymbolic representation} --- FOL, STRIPS, ontologies, and ACL versus
  subsumption, activation networks, and neural networks in RL; practice is hybrid.
\item \textbf{Theory vs.\ practice} --- elegant but non-constructive definitions (the optimal agent,
  Nash's theorem, Arrow's theorem) versus usable algorithms (A*, CBS, Q-learning, CNET, JADE).
  Do not forget complexity: the PPAD-completeness of finding a NE, and the NP-hardness of the WDP, of
  optimal MAPF, and of Slater's system.
\end{enumerate}

Small details that are often asked about: the difference between a \emph{desire} and an
\emph{intention}; why the effect of \texttt{inform} is \emph{not} that the recipient comes to believe;
why $(D,D)$ is the only \emph{non}-Pareto-optimal outcome of the prisoner's dilemma; why truthfulness
is dominant in the Vickrey auction; the difference between off-policy Q-learning and on-policy SARSA;
and what CBS does at the high level and what at the low level.



\end{document}
